\kefalaio{Γεωμετρία του επιπέδου}{Από τη γωνία στον κύκλο: όλα τα σχήματα του επιπέδου με ένα κανόνα και έναν διαβήτη}{%
\begin{tabular}{@{}l@{\hspace{1em}}p{0.82\linewidth}@{}}
3.1 & Είδη γωνιών\\[2pt]
3.2 & Μεσοκάθετος ευθύγραμμου τμήματος\\[2pt]
3.3 & Διχοτόμος γωνίας\\[2pt]
3.4 & Είδη τριγώνων και στοιχεία τους\\[2pt]
3.5 & Ιδιότητες ισοσκελούς και ισόπλευρου τριγώνου\\[2pt]
3.6 & Κατακορυφήν, συμπληρωματικές και παραπληρωματικές γωνίες\\[2pt]
3.7 & Σχετικές θέσεις ευθειών στο επίπεδο\\[2pt]
3.8 & Γωνίες που σχηματίζονται από δύο παράλληλες ευθείες και μια τέμνουσα\\[2pt]
3.9 & Άθροισμα γωνιών τριγώνου\\[2pt]
3.10 & Τετράπλευρα, παραλληλόγραμμα, τραπέζια\\[2pt]
3.11 & Είδη παραλληλογράμμων και ιδιότητες: ορθογώνιο, ρόμβος, τετράγωνο\\[2pt]
3.12 & Κύκλος και στοιχεία του\\[2pt]
3.13 & Σχετικές θέσεις ευθείας και κύκλου\\[2pt]
 & Επανάληψη του κεφαλαίου\\
\end{tabular}}
\paragrafos{Στοιχεία και είδη τριγώνων}{Β.3.1, «Στοιχεία τριγώνου – Είδη τριγώνων»}

\begin{anagke}
Στο εργαστήριο τεχνολογίας του σχολείου η Αγγελική και ο Νίκος φτιάχνουν μια μικρή γέφυρα από ξυλάκια. Ο Νίκος πρώτα φτιάχνει ένα τετράγωνο, αλλά αυτό «κουνιέται» και στραβώνει. Η κυρία Ελένη λέει: «Βάλε ένα ξυλάκι διαγώνια και θα γίνουν δύο τρίγωνα. Το τρίγωνο δεν παραμορφώνεται!».

Η Πολυξένη μετράει τα τρίγωνα της γέφυρας και βλέπει ότι δεν είναι όλα ίδια: άλλα έχουν όλες τις πλευρές ίσες, άλλα δύο ίσες, άλλα καμία. Ο Βασίλης προσθέτει ότι άλλα έχουν ορθή γωνία και άλλα όχι. Για να συνεννοούνται, χρειάζονται ονόματα: ποια τρίγωνα λέμε ισοσκελή, ποια ορθογώνια, ποια οξυγώνια.

Σε αυτή την παράγραφο μαθαίνουμε τα στοιχεία του τριγώνου, ταξινομούμε τα τρίγωνα με δύο κριτήρια (πλευρές και γωνίες) και γνωρίζουμε τη διάμεσο, το ύψος και τη διχοτόμο.
\end{anagke}

\begin{lexi}{τρίγωνο, ισοσκελές, σκαληνό, διάμεσος}{ονόματα τριγώνων}
Το \lexiword{τρίγωνο} έχει τρεις γωνίες. Το \lexiword{ισοσκελές} έχει «ίσα σκέλη» (πλευρές), το \lexiword{ισόπλευρο} όλες τις πλευρές ίσες. Το \lexiword{σκαληνό} (από το «σκαληνός» = λοξός, ανώμαλος) δεν έχει δύο ίσες πλευρές. Η \lexiword{διάμεσος} χωρίζει την πλευρά στη «μέση».
\end{lexi}

\begin{metafrasi}
Τρεις κορυφές, τρεις πλευρές, τρεις γωνίες. & $\mathrm{Α},\mathrm{Β},\Gamma$, $\mathrm{ΑΒ},\mathrm{ΒΓ},\Gamma\mathrm{Α}$, $\hat{\mathrm{Α}},\hat{\mathrm{Β}},\hat{\Gamma}$\\
Ένα τρίγωνο με ορθή γωνία. & ορθογώνιο ($\hat{\mathrm{Α}}=90^\circ$)\\
Ένα τρίγωνο με όλες τις πλευρές ίσες. & ισόπλευρο\\
Ένα τμήμα από κορυφή στο μέσο της απέναντι πλευράς. & διάμεσος\\
Ένα τμήμα από κορυφή κάθετο στην απέναντι πλευρά. & ύψος\\
\end{metafrasi}

\tmima{Η θεωρία}

\begin{orismos}{στοιχεία τριγώνου}
Κάθε τρίγωνο $\mathrm{ΑΒΓ}$ έχει τρεις \textbf{κορυφές} $\mathrm{Α},\mathrm{Β},\Gamma$, τρεις \textbf{πλευρές} $\mathrm{ΑΒ},\mathrm{ΒΓ},\Gamma\mathrm{Α}$ και τρεις \textbf{γωνίες} $\hat{\mathrm{Α}},\hat{\mathrm{Β}},\hat{\Gamma}$. Τα $\mathrm{ΑΒ},\mathrm{ΒΓ},\Gamma\mathrm{Α}$ δηλώνουν και τα μήκη των πλευρών. Περίμετρος είναι το άθροισμα $\mathrm{ΑΒ}+\mathrm{ΒΓ}+\Gamma\mathrm{Α}$.
\end{orismos}

\begin{kanonas}{ταξινόμηση τριγώνων}
\textbf{Ως προς τις γωνίες:} \emph{ορθογώνιο} (μία ορθή γωνία), \emph{αμβλυγώνιο} (μία αμβλεία γωνία), \emph{οξυγώνιο} (όλες οι γωνίες οξείες).\\
\textbf{Ως προς τις πλευρές:} \emph{ισόπλευρο} (τρεις ίσες), \emph{ισοσκελές} (δύο ίσες), \emph{σκαληνό} (όλες άνισες).
\end{kanonas}

\begin{center}
\begin{tikzpicture}[scale=0.62]
\draw[thick] (0,0)--(0,2.2)--(2.8,0)--cycle; \node at (1,-0.6) {\small ορθογώνιο};
\begin{scope}[xshift=4.6cm]
\draw[thick] (0,0)--(3,0)--(-0.8,1.6)--cycle; \node at (1.1,-0.6) {\small αμβλυγώνιο};
\end{scope}
\begin{scope}[xshift=9.6cm]
\draw[thick] (0,0)--(3,0)--(1.4,2.4)--cycle; \node at (1.5,-0.6) {\small οξυγώνιο};
\end{scope}
\end{tikzpicture}
\end{center}

\begin{center}
\begin{tikzpicture}[scale=0.62]
\draw[thick] (0,0)--(3,0)--(1.5,2.6)--cycle; \node at (1.5,-0.6) {\small ισόπλευρο};
\begin{scope}[xshift=5cm]
\draw[thick] (0,0)--(3,0)--(1.5,2.6)--cycle; \node at (1.5,-0.6) {\small ισοσκελές};
\end{scope}
\begin{scope}[xshift=10cm]
\draw[thick] (0,0)--(3.3,0)--(0.9,2.2)--cycle; \node at (1.5,-0.6) {\small σκαληνό};
\end{scope}
\end{tikzpicture}
\end{center}

\begin{orismos}{δευτερεύοντα στοιχεία}
\textbf{Διάμεσος} τριγώνου: τμήμα που ενώνει μια κορυφή με το μέσο της απέναντι πλευράς.\\
\textbf{Ύψος}: τμήμα από μια κορυφή, κάθετο στην ευθεία της απέναντι πλευράς.\\
\textbf{Διχοτόμος}: το τμήμα της διχοτόμου μιας γωνίας του τριγώνου από την κορυφή ως την απέναντι πλευρά.
\end{orismos}

Κάθε τρίγωνο έχει τρεις διαμέσους, τρία ύψη και τρεις διχοτόμους. Οι τρεις διάμεσοι τέμνονται σε ένα σημείο (βαρύκεντρο), τα τρία ύψη επίσης (ορθόκεντρο). Στο αμβλυγώνιο τρίγωνο τα ύψη φέρονται εκτός τριγώνου, πάνω στις προεκτάσεις των πλευρών. Στο ορθογώνιο τρίγωνο οι δύο κάθετες πλευρές είναι και ύψη.

\begin{prosoxi}
\begin{itemize}[leftmargin=1.2em,itemsep=0.3ex,topsep=2pt]
\item Το ισόπλευρο είναι και ισοσκελές (έχει δύο ίσες πλευρές), αλλά όχι ανάποδα.
\item Ένα τρίγωνο έχει το πολύ μία ορθή ή μία αμβλεία γωνία.
\item Το ύψος είναι \emph{κάθετο} τμήμα, η διάμεσος πάει στο \emph{μέσο}. Δεν είναι πάντα το ίδιο τμήμα.
\item Το ύψος μπορεί να βρίσκεται έξω από το τρίγωνο.
\end{itemize}
\end{prosoxi}

\tmima{Λυμένα παραδείγματα}

\paradeigma{Η περίμετρος και το είδος}
\tixero Τρίγωνο με πλευρές $5$, $5$ και $8$ cm.\par
\tithelo Το είδος του ως προς τις πλευρές και την περίμετρο.\par
\possindeo Δύο πλευρές είναι ίσες, άρα ισοσκελές. Περίμετρος $5+5+8=18$ cm.\par
\apantisi Ισοσκελές, περίμετρος $18$ cm.\par
\elegxos Το άθροισμα των δύο μικρότερων ($5+5=10$) είναι μεγαλύτερο από την τρίτη ($8$), άρα το τρίγωνο υπάρχει.

\paradeigma{Το ύψος}
\tixero Τρίγωνο $\mathrm{ΑΒΓ}$ με $\mathrm{ΒΓ}=10$ cm και ύψος $\mathrm{ΑΔ}=4$ cm.\par
\tithelo Σχεδίασε τα ύψη όταν το τρίγωνο είναι οξυγώνιο, ορθογώνιο ή αμβλυγώνιο.\par
\possindeo Από κάθε κορυφή φέρνω κάθετο στην απέναντι πλευρά (ή την προέκτασή της, στο αμβλυγώνιο).\par
\apantisi Τα τρία ύψη περνούν από το ίδιο σημείο (ορθόκεντρο).\par
\elegxos Στο ορθογώνιο τρίγωνο το ορθόκεντρο είναι η κορυφή της ορθής γωνίας.

\paradeigma{Η διάμεσος}
\tixero Στο τρίγωνο $\mathrm{ΑΒΓ}$ η πλευρά $\mathrm{ΒΓ}=7{,}4$ cm και η $\mathrm{ΑΜ}$ είναι διάμεσος.\par
\tithelo Το μήκος $\mathrm{ΒΜ}$.\par
\possindeo Το $\mathrm{Μ}$ είναι το μέσο της $\mathrm{ΒΓ}$, άρα $\mathrm{ΒΜ}=7{,}4:2=3{,}7$ cm.\par
\apantisi $3{,}7$ cm.\par
\elegxos $\mathrm{ΒΜ}+\mathrm{ΜΓ}=3{,}7+3{,}7=7{,}4$.

\tmima{Ασκήσεις}

\zoni{Το κατάλαβα}{ονόματα και είδη}

\ex Ονόμασε τα τρίγωνα με πλευρές: α) $4,4,4$ β) $5,5,3$ γ) $3,4,5$ ως προς τις πλευρές.

\ex Βρες την περίμετρο τριγώνου με πλευρές $6{,}5$, $7$ και $8{,}5$ cm.

\ex Σχεδίασε οξυγώνιο, ορθογώνιο και αμβλυγώνιο τρίγωνο και φέρε τα ύψη τους.

\ex Στο τρίγωνο $\mathrm{ΑΒΓ}$ με $\mathrm{ΒΓ}=4{,}4$ cm φέρε τη διάμεσο $\mathrm{ΑΜ}$. Πόσο είναι το $\mathrm{ΒΜ}$;

\ex Σχεδίασε τρίγωνο $\mathrm{ΑΒΓ}$, βρες τα μέσα $\mathrm{Δ}$, $\mathrm{Ε}$, $\mathrm{Ζ}$ των πλευρών και σχεδίασε τις τρεις διαμέσους. Τι παρατηρείς;

\ex Σχεδίασε τρίγωνο και τη διχοτόμο μιας γωνίας του. Μέτρησε τις δύο γωνίες που σχηματίζονται.

\ex Σωστό ή λάθος; α) Το ισόπλευρο είναι ισοσκελές. \quad β) Το αμβλυγώνιο έχει δύο αμβλείες. \quad γ) Το ορθογώνιο τρίγωνο μπορεί να είναι ισοσκελές. \quad δ) Το ισόπλευρο είναι πάντα οξυγώνιο. \quad ε) Ένα σκαληνό δεν μπορεί να είναι ορθογώνιο.

\zoni{Το χρησιμοποιώ}{ξύλινες κατασκευές}

\ex \textbf{Η γέφυρα.} Ένα ξύλινο τρίγωνο της γέφυρας έχει πλευρές $30$, $30$ και $40$ cm. Τι είδους είναι και πόση είναι η περίμετρος;

\ex \textbf{Το μαντήλι.} Ένα τριγωνικό μαντήλι έχει πλευρές $60$, $60$ και $60$ cm. Πόσα εκατοστά κορδόνι χρειάζεται για να το γυρίσει κανείς γύρω γύρω;

\ex \textbf{Το σχοινί.} Ένα σχοινί $24$ μ. σχηματίζει ισόπλευρο τρίγωνο. Πόση είναι η πλευρά του;

\ex \textbf{Το τρίγωνο της κορνίζας.} Ισοσκελές τρίγωνο έχει περίμετρο $50$ cm και βάση $20$ cm. Πόσο μακριές είναι οι δύο ίσες πλευρές;

\ex \textbf{Η κατασκευή.} Μπορείς να φτιάξεις τρίγωνο με ξυλάκια $3$, $4$ και $9$ cm; Και με $3$, $4$ και $5$ cm; Εξήγησε.

\zoni{Το σκέφτομαι}{σύγκριση ειδών}

\ex Σε ποια είδη τριγώνων ως προς τις γωνίες μπορεί να ανήκει ένα ισοσκελές τρίγωνο;

\ex Μπορεί ένα ορθογώνιο τρίγωνο να είναι ισόπλευρο; Εξήγησε.

\ex Πόσα το πολύ ύψη του τριγώνου βρίσκονται εκτός αυτού; (Σκέψου το αμβλυγώνιο.)

\ex Συμπλήρωσε πίνακα: ποιοι συνδυασμοί (ορθογώνιο, αμβλυγώνιο, οξυγώνιο) × (σκαληνό, ισοσκελές, ισόπλευρο) υπάρχουν;

\zoni{Για το σπίτι}{τρίγωνα στην πόλη}

\ex Βρες πέντε τριγωνικά αντικείμενα στο σπίτι (σήμα, φέτα πίτσας, ράφι, στέγη κ.λπ.) και πες το είδος του καθενός.

\ex Κόψε από χαρτί τρίγωνο και δίπλωσε το ώστε να βρεις το μέσο μιας πλευράς και τη διάμεσο.

\paragrafos{Άθροισμα γωνιών τριγώνου}{Β.3.2, «Άθροισμα γωνιών τριγώνου – Ιδιότητες ισοσκελούς τριγώνου»}

\begin{anagke}
Ο Βασίλης κόβει από χαρτί ένα τρίγωνο και σκίζει τις τρεις γωνίες του. Τις βάζει δίπλα δίπλα, με τις κορυφές στο ίδιο σημείο, και οι τρεις μυτούλες σχηματίζουν μια ευθεία γραμμή. «Ευθεία γωνία!», φωνάζει. Η Αγγελική κάνει το ίδιο με ένα άλλο τρίγωνο, πιο λεπτό και μακρύ: πάλι ευθεία. Ο Νίκος το κάνει με αμβλυγώνιο: ξανά ευθεία.

Η Πολυξένη ρωτάει: «Ισχύει πάντα; Και πώς το ξέρουμε χωρίς να κόψουμε όλα τα τρίγωνα του κόσμου;». Η απάντηση είναι ότι το αποδεικνύουμε με όσα ξέρουμε για τις παράλληλες ευθείες. Η Ρήνη προσθέτει: «Και τα ισοσκελή; Αν διπλώσω ένα ισοσκελές στη μέση, οι δύο μισές πέφτουν η μία πάνω στην άλλη!».

Σε αυτή την παράγραφο μαθαίνουμε ότι το άθροισμα των γωνιών κάθε τριγώνου είναι $180^\circ$, τι ισχύει για την εξωτερική γωνία, και τις ιδιότητες του ισοσκελούς και του ισόπλευρου τριγώνου.
\end{anagke}

\begin{lexi}{εξωτερική γωνία, βάση}{ονόματα}
Η \lexiword{εξωτερική γωνία} είναι αυτή που σχηματίζεται από μια πλευρά και την προέκταση της διπλανής. Η \lexiword{βάση} του ισοσκελούς είναι η πλευρά που δεν είναι από τις δύο ίσες, και οι \lexiword{προσκείμενες γωνίες} της είναι οι δύο γωνίες στα άκρα της.
\end{lexi}

\begin{metafrasi}
Οι τρεις γωνίες φτιάχνουν ευθεία γωνία. & $\hat{\mathrm{Α}}+\hat{\mathrm{Β}}+\hat{\Gamma}=180^\circ$\\
Οι δύο οξείες του ορθογωνίου τριγώνου. & συμπληρωματικές: $\hat{\mathrm{Β}}+\hat{\Gamma}=90^\circ$\\
Η εξωτερική γωνία στο $\Gamma$. & $\hat{\varphi}=\hat{\mathrm{Α}}+\hat{\mathrm{Β}}$\\
Οι γωνίες στη βάση του ισοσκελούς. & $\hat{\mathrm{Β}}=\hat{\Gamma}$\\
Όλες οι γωνίες του ισόπλευρου. & $60^\circ$\\
\end{metafrasi}

\tmima{Η θεωρία}

\begin{kanonas}{άθροισμα γωνιών τριγώνου}
Σε κάθε τρίγωνο $\mathrm{ΑΒΓ}$: $\hat{\mathrm{Α}}+\hat{\mathrm{Β}}+\hat{\Gamma}=180^\circ$.
\end{kanonas}

\textbf{Γιατί ισχύει.} Φέρνω από το $\mathrm{Α}$ την ευθεία $x\mathrm{Α}y$ παράλληλη προς τη $\mathrm{ΒΓ}$. Η $x\hat{\mathrm{Α}}\mathrm{Β}$ είναι ίση με την $\hat{\mathrm{Β}}$ (εντός εναλλάξ) και η $y\hat{\mathrm{Α}}\Gamma$ ίση με την $\hat{\Gamma}$ (εντός εναλλάξ). Οι τρεις γωνίες γύρω από το $\mathrm{Α}$ σχηματίζουν ευθεία γωνία, άρα $\hat{\mathrm{Β}}+\hat{\mathrm{Α}}+\hat{\Gamma}=180^\circ$.

\begin{center}
\begin{tikzpicture}[scale=0.9]
\draw[thick] (-0.4,0)--(3.8,0) ;
\node[below left] at (0,0) {$\mathrm{Β}$}; \node[below right] at (3.4,0) {$\Gamma$};
\draw[thick] (0,0)--(1.4,1.8)--(3.4,0)--cycle;
\draw[thick,gblue] (-0.6,1.8)--(3.6,1.8);
\node[above] at (1.4,1.8) {$\mathrm{Α}$};
\node at (0.55,0.2) {\small $\omega$}; \node at (2.9,0.2) {\small $\theta$};
\node at (0.7,1.6) {\small $\omega$}; \node at (2.1,1.6) {\small $\theta$};
\node at (1.4,1.45) {\small $\hat{\mathrm{Α}}$};
\end{tikzpicture}
\end{center}

\begin{kanonas}{συνέπειες}
\begin{itemize}[leftmargin=1.2em,itemsep=0.3ex,topsep=2pt]
\item Στο \textbf{ορθογώνιο} τρίγωνο οι δύο οξείες γωνίες είναι \textbf{συμπληρωματικές}: $\hat{\mathrm{Β}}+\hat{\Gamma}=90^\circ$.
\item Η \textbf{εξωτερική γωνία} τριγώνου ισούται με το άθροισμα των δύο \emph{απέναντι} εσωτερικών: $\hat{\varphi}=\hat{\mathrm{Α}}+\hat{\mathrm{Β}}$ (γιατί $\hat{\varphi}=180^\circ-\hat{\Gamma}$).
\item Στο \textbf{ισόπλευρο} τρίγωνο κάθε γωνία είναι $180^\circ:3=60^\circ$.
\end{itemize}
\end{kanonas}

\begin{kanonas}{ισοσκελές τρίγωνο}
Σε ισοσκελές τρίγωνο ($\mathrm{ΑΒ}=\mathrm{ΑΓ}$, κορυφή $\mathrm{Α}$, βάση $\mathrm{ΒΓ}$):
\begin{itemize}[leftmargin=1.2em,itemsep=0.3ex,topsep=2pt]
\item η ευθεία της διαμέσου προς τη βάση είναι \textbf{άξονας συμμετρίας},
\item η διάμεσος προς τη βάση είναι και \textbf{ύψος} και \textbf{διχοτόμος},
\item οι \textbf{προσκείμενες στη βάση γωνίες είναι ίσες}: $\hat{\mathrm{Β}}=\hat{\Gamma}$.
\end{itemize}
Στο \textbf{ισόπλευρο}: οι ευθείες των διαμέσων είναι άξονες συμμετρίας, κάθε διάμεσος είναι ύψος και διχοτόμος, όλες οι πλευρές και γωνίες είναι ίσες.
\end{kanonas}

\begin{prosoxi}
\begin{itemize}[leftmargin=1.2em,itemsep=0.3ex,topsep=2pt]
\item Στο ισοσκελές, η ιδιότητα «διάμεσος = ύψος = διχοτόμος» ισχύει για το τμήμα προς τη \emph{βάση}, όχι για τα άλλα.
\item Στο ισοσκελές ορθογώνιο τρίγωνο οι δύο οξείες είναι $45^\circ$ η καθεμία.
\item Η εξωτερική γωνία είναι ίση με το άθροισμα των \emph{δύο άλλων} γωνιών (όχι της διπλανής της).
\end{itemize}
\end{prosoxi}

\tmima{Λυμένα παραδείγματα}

\paradeigma{Η τρίτη γωνία}
\tixero Τρίγωνο με γωνίες $\hat{\mathrm{Α}}=47^\circ$ και $\hat{\mathrm{Β}}=68^\circ$.\par
\tithelo Η $\hat{\Gamma}$.\par
\possindeo $\hat{\Gamma}=180^\circ-(47^\circ+68^\circ)=180^\circ-115^\circ=65^\circ$.\par
\apantisi $65^\circ$ (οξυγώνιο τρίγωνο).\par
\elegxos $47+68+65=180$.

\paradeigma{Η εξωτερική γωνία}
\tixero Τρίγωνο με $\hat{\mathrm{Α}}=50^\circ$, $\hat{\mathrm{Β}}=60^\circ$.\par
\tithelo Η εξωτερική γωνία στο $\Gamma$.\par
\possindeo $\hat{\varphi}=\hat{\mathrm{Α}}+\hat{\mathrm{Β}}=110^\circ$. Ή: $\hat{\Gamma}=70^\circ$ και $\hat{\varphi}=180^\circ-70^\circ=110^\circ$.\par
\apantisi $110^\circ$.\par
\elegxos Οι δύο τρόποι δίνουν το ίδιο.

\paradeigma{Ισοσκελές με γωνία κορυφής}
\tixero Ισοσκελές τρίγωνο με γωνία κορυφής $\hat{\mathrm{Α}}=40^\circ$.\par
\tithelo Οι γωνίες της βάσης.\par
\possindeo $\hat{\mathrm{Β}}=\hat{\Gamma}$ και $\hat{\mathrm{Β}}+\hat{\Gamma}=180^\circ-40^\circ=140^\circ$. Άρα $\hat{\mathrm{Β}}=\hat{\Gamma}=70^\circ$.\par
\apantisi $70^\circ$ και $70^\circ$.\par
\elegxos $40+70+70=180$.

\paradeigma{Ορθογώνιο και ισοσκελές}
\tixero Ένα τρίγωνο είναι ορθογώνιο και ισοσκελές.\par
\tithelo Οι γωνίες του.\par
\possindeo Η ορθή είναι $90^\circ$. Οι άλλες δύο είναι ίσες και με άθροισμα $90^\circ$, άρα $45^\circ$ η καθεμία.\par
\apantisi $90^\circ,45^\circ,45^\circ$.\par
\elegxos $90+45+45=180$.

\tmima{Ασκήσεις}

\zoni{Το κατάλαβα}{γωνίες τριγώνων}

\ex Βρες την τρίτη γωνία τριγώνου με γωνίες: α) $35^\circ,\,80^\circ$ β) $90^\circ,\,28^\circ$ γ) $112^\circ,\,41^\circ$.

\ex Ισοσκελές τρίγωνο έχει γωνία κορυφής $50^\circ$. Βρες τις γωνίες της βάσης.

\ex Ισοσκελές τρίγωνο έχει γωνία στη βάση $72^\circ$. Βρες τη γωνία κορυφής.

\ex Πόσο είναι κάθε γωνία ισοπλεύρου τριγώνου; Πόσο η εξωτερική του γωνία;

\ex Οι οξείες ενός ορθογωνίου τριγώνου είναι $x$ και $2x$. Βρες τες.

\ex Σε τρίγωνο η εξωτερική γωνία στο $\Gamma$ είναι $125^\circ$ και $\hat{\mathrm{Α}}=48^\circ$. Βρες τη $\hat{\mathrm{Β}}$ και τη $\hat{\Gamma}$.

\ex Σωστό ή λάθος; α) Τρίγωνο μπορεί να έχει δύο ορθές γωνίες. \quad β) Οι γωνίες βάσης ισοσκελούς είναι ίσες. \quad γ) Σε ισόπλευρο τρίγωνο όλες οι γωνίες είναι $60^\circ$. \quad δ) Οι οξείες γωνίες ορθογωνίου είναι παραπληρωματικές.

\zoni{Το χρησιμοποιώ}{στέγες και σήματα}

\ex \textbf{Η στέγη.} Μια στέγη έχει σχήμα ισοσκελούς τριγώνου με γωνία κορυφής $100^\circ$. Πόση είναι η γωνία που σχηματίζει κάθε πλευρά με το οριζόντιο δοκάρι;

\ex \textbf{Το σήμα.} Ένα τριγωνικό σήμα έχει ισόπλευρο σχήμα. Πόσο είναι το άθροισμα των τριών γωνιών του;

\ex \textbf{Η σκάλα.} Μια σκάλα ακουμπά στον τοίχο και σχηματίζει με αυτόν γωνία $25^\circ$. Ο τοίχος είναι κάθετος στο έδαφος. Πόση γωνία σχηματίζει η σκάλα με το έδαφος;

\ex \textbf{Ο χαρταετός.} Ένα τριγωνικό κομμάτι ενός χαρταετού έχει γωνίες $55^\circ$ και $75^\circ$. Πόση είναι η τρίτη γωνία;

\ex \textbf{Η γωνιά.} Ένας ξύλινος γνώμονας έχει σχήμα ισοσκελούς ορθογωνίου τριγώνου. Πόσες μοίρες είναι οι δύο ίσες γωνίες του;

\zoni{Το σκέφτομαι}{γιατί συμβαίνει}

\ex Εξήγησε με παράλληλη από την κορυφή γιατί το άθροισμα των γωνιών ενός τριγώνου είναι $180^\circ$.

\ex Εξήγησε γιατί τρίγωνο δεν μπορεί να έχει δύο αμβλείες γωνίες.

\ex Αν σε ισοσκελές τρίγωνο μια γωνία είναι $100^\circ$, ποια θέση έχει (κορυφής ή βάσης); Βρες τις άλλες δύο.

\ex Στο ισοσκελές τρίγωνο $\mathrm{ΑΒΓ}$ ($\mathrm{ΑΒ}=\mathrm{ΑΓ}$) η $\mathrm{ΑΔ}$ είναι διάμεσος προς τη $\mathrm{ΒΓ}$ και $\mathrm{ΒΓ}=12$ cm. Πόσο είναι το $\mathrm{ΒΔ}$ και πόση η γωνία $\mathrm{ΑΔΒ}$;

\zoni{Για το σπίτι}{γωνίες με το χέρι}

\ex Κόψε ένα τρίγωνο από χαρτί, σκίσε τις γωνίες του και τοποθέτησέ τες δίπλα δίπλα. Τι σχηματίζουν;

\ex Δίπλωσε ένα ισοσκελές τρίγωνο στη μέση και δες ποιες γωνίες πέφτουν η μία πάνω στην άλλη.

\paragrafos{Διχοτόμος γωνίας}{Β.3.3, «Διχοτόμος γωνίας»}

\begin{anagke}
Στο διάλειμμα ο Νίκος και η Αγγελική μοιράζονται ένα κομμάτι πίτσα σε σχήμα «φέτας» με μυτερή γωνία. «Θέλω ακριβώς τη μισή», λέει η Αγγελική, «και όχι με το μάτι». Ο Νίκος κοιτάζει τη φέτα: «Αν κόψω από την μύτη προς την άκρη, πώς θα ξέρω ότι τα δύο κομμάτια είναι ίσα;». Η Φανή το σημειώνει: ζητούμενο είναι μια ευθεία από την κορυφή που χωρίζει τη γωνία σε δύο ίσες γωνίες.

Η Ρήνη βρίσκει και ένα δεύτερο πρόβλημα στο διπλανό γήπεδο: «Το καφενείο θέλει να βάλει ένα τραπέζι σε μια γωνία, ώστε να απέχει ίσα από τους δύο τοίχους». Ο Ηλίας υποψιάζεται ότι οι δύο προβλήματα είναι ένα: η ευθεία που κόβει τη γωνία στη μέση είναι και η ευθεία των σημείων που ισαπέχουν από τις πλευρές της.

Σε αυτή την παράγραφο μαθαίνουμε τι είναι η διχοτόμος μιας γωνίας, ποια ιδιότητα έχουν τα σημεία της και πώς τη σχεδιάζουμε με μοιρογνωμόνιο, με δίπλωση και με διαβήτη.
\end{anagke}

\begin{lexi}{διχοτόμος, απόσταση}{διχοτόμος}
Η \lexiword{διχοτόμος} είναι αυτή που «κόβει δίχα», δηλαδή στα δύο (δίχα + τέμνω): χωρίζει τη γωνία σε δύο ίσα μέρη. Η \lexiword{απόσταση} σημείου από ευθεία είναι το μήκος του κάθετου τμήματος που φέρνουμε από το σημείο προς την ευθεία.
\end{lexi}

\begin{metafrasi}
Η ημιευθεία που κόβει τη γωνία στη μέση. & η διχοτόμος $\mathrm{Οδ}$ της $x\hat{\mathrm{Ο}}y$\\
Οι δύο μισές γωνίες είναι ίσες. & $x\hat{\mathrm{Ο}}\delta=\delta\hat{\mathrm{Ο}}y$\\
Κάθε μισή είναι το μισό της γωνίας. & $x\hat{\mathrm{Ο}}\delta=\dfrac{x\hat{\mathrm{Ο}}y}{2}$\\
Το σημείο $\mathrm{Α}$ έχει ίσες αποστάσεις από τις δύο πλευρές. & $\mathrm{ΑΒ}=\mathrm{ΑΓ}$\\
\end{metafrasi}

\tmima{Η θεωρία}

\begin{orismos}{διχοτόμος γωνίας}
\textbf{Διχοτόμος} μιας γωνίας $x\hat{\mathrm{Ο}}y$ λέγεται η ημιευθεία $\mathrm{Οδ}$ που έχει αρχή την κορυφή της γωνίας και τη χωρίζει σε δύο ίσες γωνίες: $x\hat{\mathrm{Ο}}\delta=\delta\hat{\mathrm{Ο}}y$.
\end{orismos}

Κάθε γωνία έχει μία και μόνο μία διχοτόμο. Κάθε μία από τις δύο μισές γωνίες έχει μέτρο ίσο με το μισό της αρχικής.

\begin{center}
\begin{tikzpicture}[scale=0.95]
\draw[thick] (5,0)--(0,0)--({5*cos(60)},{5*sin(60)});
\draw[gblue,thick] (0,0)--({5*cos(30)},{5*sin(30)});
\fill (0,0) circle (2pt) node[below left] {$\mathrm{Ο}$};
\node[right] at (5,0) {$x$}; \node[above] at ({5*cos(60)},{5*sin(60)}) {$y$}; \node[right,gblue] at ({5*cos(30)},{5*sin(30)}) {$\delta$};
\draw[gblue] (1.1,0) arc(0:30:1.1); \draw[gblue] (1.4,0) arc(0:30:1.4); \draw[gwarmD] (1.3,0.75) arc(30:60:1.5); \draw[gwarmD] (1.6,0.92) arc(30:60:1.85);
\end{tikzpicture}
\end{center}

Θυμόμαστε ότι \textbf{απόσταση} ενός σημείου $\mathrm{Α}$ από μια ευθεία $\varepsilon$ είναι το μήκος του κάθετου τμήματος $\mathrm{ΑΒ}$ που φέρνουμε από το $\mathrm{Α}$ προς την $\varepsilon$.

\begin{idiotita}{σημεία της διχοτόμου}
Κάθε σημείο της διχοτόμου μιας γωνίας \textbf{ισαπέχει} από τις πλευρές της.
\end{idiotita}

Ισχύει και το αντίστροφο: κάθε σημείο στο εσωτερικό της γωνίας που ισαπέχει από τις πλευρές της ανήκει στη διχοτόμο της. Έτσι η διχοτόμος είναι ακριβώς το σύνολο των σημείων που ισαπέχουν από τις πλευρές της γωνίας, όπως η μεσοκάθετος είναι το σύνολο των σημείων που ισαπέχουν από δύο σημεία.

\begin{center}
\begin{tikzpicture}[scale=0.95]
\draw[thick] (5,0)--(0,0)--({5*cos(60)},{5*sin(60)});
\draw[gblue,thick] (0,0)--({4.5*cos(30)},{4.5*sin(30)});
\fill (0,0) circle (2pt) node[below left] {$\mathrm{Ο}$};
\coordinate (P) at (2.944,1.7); \coordinate (B) at (2.944,0); \coordinate (G) at (1.472,2.55);
\fill (P) circle (2pt) node[above right] {$\mathrm{Α}$};
\draw[gwarmD,thick] (P)--(B) node[midway,right] {$\delta$}; \draw[gwarmD,thick] (P)--(G) node[midway,above right] {$\delta$};
\node[below] at (B) {$\mathrm{Β}$}; \node[left] at (G) {$\Gamma$};
\draw (2.744,0)--(2.744,0.2)--(2.944,0.2);
\draw (1.372,2.377)--(1.545,2.277)--(1.645,2.45);
\end{tikzpicture}
\end{center}

\tmimasmall{Πώς σχεδιάζουμε τη διχοτόμο}

\textbf{Με μοιρογνωμόνιο.} Μετράμε τη γωνία, διαιρούμε το μέτρο με το $2$ και σημειώνουμε αυτή τη γωνία πάνω στη μία πλευρά.

\textbf{Με δίπλωση.} Διπλώνουμε διαφανές χαρτί ώστε οι δύο πλευρές της γωνίας να συμπέσουν. Η πτυχή είναι η διχοτόμος.

\textbf{Με κανόνα και διαβήτη.} Με κέντρο την κορυφή $\mathrm{Ο}$ γράφουμε τόξο που τέμνει τις πλευρές στα $\mathrm{Κ}$ και $\Lambda$. Με κέντρα τα $\mathrm{Κ}$ και $\Lambda$ και την ίδια ακτίνα γράφουμε δύο τόξα που τέμνονται στο $\mathrm{Π}$. Η $\mathrm{ΟΠ}$ είναι η διχοτόμος.

\begin{center}
\begin{tikzpicture}[scale=0.95]
\draw[thick] (5,0)--(0,0)--({5*cos(70)},{5*sin(70)});
\fill (0,0) circle (2pt) node[below left] {$\mathrm{Ο}$};
\draw[gblue] (2,0) arc(0:70:2);
\coordinate (K) at (2,0); \coordinate (L) at ({2*cos(70)},{2*sin(70)});
\fill (K) circle (2pt) node[below] {$\mathrm{Κ}$}; \fill (L) circle (2pt) node[left] {$\Lambda$};
\draw[gwarmD] (K)++(55:1.8) arc(55:95:1.8);
\draw[gwarmD] (L)++(-25:1.8) arc(-25:15:1.8);
\fill (2.478,1.735) circle (2pt) node[above right] {$\mathrm{Π}$};
\draw[gblue,thick] (0,0)--(3.9,2.73);
\end{tikzpicture}
\end{center}

\begin{prosoxi}
\begin{itemize}[leftmargin=1.2em,itemsep=0.3ex,topsep=2pt]
\item Η διχοτόμος είναι ημιευθεία μέσα στη γωνία, με αρχή την κορυφή. Δεν είναι τμήμα που ενώνει δύο σημεία των πλευρών.
\item Η απόσταση σημείου από ευθεία μετριέται πάντα κατά το \emph{κάθετο} τμήμα, όχι κατά ένα λοξό.
\item Σημείο στο εσωτερικό της γωνίας με άνισες αποστάσεις από τις πλευρές δεν ανήκει στη διχοτόμο.
\item Οι δύο μισές γωνίες είναι ίσες σε \emph{μέτρο}, ακόμη κι αν στο σχέδιο οι πλευρές τους έχουν διαφορετικό μήκος.
\end{itemize}
\end{prosoxi}

\tmima{Λυμένα παραδείγματα}

\paradeigma{Η φέτα της πίτσας}
\tixero Η γωνία της κορυφής της φέτας είναι $60^\circ$.\par
\tithelo Με ποια γωνία ξεκινά η κοψιά που τη χωρίζει σε δύο ίσα μέρη;\par
\possindeo Η κοψιά είναι η διχοτόμος της γωνίας, άρα χωρίζει τη γωνία σε δύο ίσες γωνίες.\par
\apantisi $60^\circ:2=30^\circ$ με κάθε πλευρά.\par
\elegxos $30^\circ+30^\circ=60^\circ$.

\paradeigma{Η διχοτόμος και το μέτρο}
\tixero Η $\mathrm{Οδ}$ είναι διχοτόμος της $x\hat{\mathrm{Ο}}y$ και $x\hat{\mathrm{Ο}}\delta=28^\circ$.\par
\tithelo Το μέτρο της $x\hat{\mathrm{Ο}}y$.\par
\possindeo Οι δύο μισές γωνίες είναι ίσες: $\delta\hat{\mathrm{Ο}}y=x\hat{\mathrm{Ο}}\delta=28^\circ$.\par
\apantisi $x\hat{\mathrm{Ο}}y=28^\circ+28^\circ=56^\circ$.\par
\elegxos Η διχοτόμος των $56^\circ$ δίνει δύο γωνίες των $28^\circ$.

\paradeigma{Το ρολόι ως διχοτόμος}
\tixero Στις $4{:}00$ ο λεπτοδείκτης δείχνει το $12$ και ο ωροδείκτης το $4$. Κάθε ώρα είναι $30^\circ$.\par
\tithelo Ποιον αριθμό του ρολογιού δείχνει η διχοτόμος της γωνίας των δεικτών;\par
\possindeo Η γωνία των δεικτών είναι $4\cdot30^\circ=120^\circ$ και η μισή της είναι $60^\circ$, που αντιστοιχεί σε $60^\circ:30^\circ=2$ ώρες από το $12$.\par
\apantisi Η διχοτόμος δείχνει το $2$.\par
\elegxos Από το $12$ ως το $2$ είναι $60^\circ$ και από το $2$ ως το $4$ άλλες $60^\circ$.

\paradeigma{Το σημείο που ισαπέχει}
\tixero Μια γωνία $x\hat{\mathrm{Ο}}y=80^\circ$.\par
\tithelo Να βρω σημείο $\mathrm{Α}$ στο εσωτερικό της που να ισαπέχει από τις πλευρές της και να απέχει $5$ cm από την κορυφή.\par
\possindeo Τα σημεία που ισαπέχουν από τις πλευρές είναι τα σημεία της διχοτόμου. Φέρνω τη διχοτόμο (γωνία $40^\circ$ με κάθε πλευρά) και παίρνω πάνω της σημείο $\mathrm{Α}$ με $\mathrm{ΟΑ}=5$ cm.\par
\apantisi Το $\mathrm{Α}$ είναι πάνω στη διχοτόμο, σε απόσταση $5$ cm από το $\mathrm{Ο}$.\par
\elegxos Μετράω τις κάθετες από το $\mathrm{Α}$ προς τις πλευρές: βγαίνουν ίσες.

\tmima{Ασκήσεις}

\zoni{Το κατάλαβα}{διχοτόμος και ισαπόσταση}

\ex Βρες το μέτρο κάθε μισής γωνίας όταν η γωνία είναι: α) $70^\circ$ \quad β) $124^\circ$ \quad γ) $36^\circ$ \quad δ) $90^\circ$ \quad ε) $180^\circ$.

\ex Σχεδίασε μια γωνία $x\hat{\mathrm{Ο}}y=64^\circ$ και φέρε τη διχοτόμο της: α) με μοιρογνωμόνιο, β) με δίπλωση διαφανούς χαρτιού. Συμπίπτουν;

\ex Η $\mathrm{Οδ}$ είναι διχοτόμος της $x\hat{\mathrm{Ο}}y$ και $x\hat{\mathrm{Ο}}\delta=37^\circ$. Βρες το μέτρο της $x\hat{\mathrm{Ο}}y$.

\ex Σχεδίασε ένα τρίγωνο $\mathrm{ΑΒΓ}$ και φέρε τις διχοτόμους των γωνιών του. Τι παρατηρείς;

\ex Σωστό ή λάθος; α) Η διχοτόμος είναι ημιευθεία. \quad β) Η διχοτόμος χωρίζει τη γωνία σε δύο άνισα μέρη. \quad γ) Κάθε σημείο της διχοτόμου ισαπέχει από τις πλευρές της γωνίας. \quad δ) Η διχοτόμος γωνίας $80^\circ$ σχηματίζει με κάθε πλευρά γωνία $40^\circ$. \quad ε) Η διχοτόμος της ευθείας γωνίας είναι κάθετη στις πλευρές της.

\ex Σχεδίασε μια ευθεία $\varepsilon$ και ένα σημείο $\mathrm{Α}$ εκτός της. Φέρε με τον γνώμονα την κάθετη από το $\mathrm{Α}$ προς την $\varepsilon$ και μέτρησε την απόσταση του $\mathrm{Α}$ από την $\varepsilon$.

\ex Σχεδίασε γωνία $x\hat{\mathrm{Ο}}y=80^\circ$ και πάρε πάνω στη διχοτόμο της ένα σημείο $\mathrm{Α}$ με $\mathrm{ΟΑ}=5$ cm. Μέτρησε τις αποστάσεις του $\mathrm{Α}$ από τις $\mathrm{Ο}x$ και $\mathrm{Ο}y$.

\zoni{Το χρησιμοποιώ}{πίτσες, ρολόγια, δρόμοι}

\ex \textbf{Η πίτσα.} Μια φέτα πίτσα έχει γωνία κορυφής $50^\circ$. Ο Νίκος τη χωρίζει στα δύο με την διχοτόμο. Πόση γωνία έχει καθένα από τα δύο κομμάτια;

\ex \textbf{Το μονοπάτι.} Η γωνία της γωνιάς του φράχτη στον κήπο της θείας Μαρίας είναι $84^\circ$. Ένα μονοπάτι τη χωρίζει σε δύο ίσα παρτέρια. Με ποια γωνία ξεκινά το μονοπάτι από κάθε πλευρά του φράχτη;

\ex \textbf{Το ρολόι.} Στις $8{:}00$, ποιον αριθμό του ρολογιού δείχνει η διχοτόμος της (μικρότερης) γωνίας των δεικτών; Και στις $2{:}00$;

\ex \textbf{Ο ανεμιστήρας.} Ένας ανεμιστήρας της Αγγελικής έχει τρία πτερύγια που σχηματίζουν μεταξύ τους ίσες γωνίες. Πόση γωνία έχει η διχοτόμος δύο διαδοχικών πτερυγίων με καθένα από αυτά;

\ex \textbf{Το περίπτερο.} Δύο ευθείες οδοί τέμνονται σε γωνία $60^\circ$ και ένα περίπτερο πρέπει να ισαπέχει από τις δύο οδούς και να είναι $6$ cm από τη διασταύρωση στο σχέδιο. Σχεδίασε πού μπαίνει.

\zoni{Το σκέφτομαι}{γιατί και πώς}

\ex Γιατί κάθε γωνία έχει μία και μοναδική διχοτόμο; Τι γίνεται με τη διχοτόμο της μηδενικής γωνίας;

\ex Ένα σημείο στο εσωτερικό μιας γωνίας απέχει $3$ cm από τη μία πλευρά και $5$ cm από την άλλη. Ανήκει στη διχοτόμο; Πού να το μετακινήσεις ώστε να ανήκει;

\ex Δύο γωνίες είναι διαδοχικές με άθροισμα $180^\circ$. Η πρώτη είναι $50^\circ$. Σχεδίασε τις διχοτόμους και βρες τη γωνία που σχηματίζουν μεταξύ τους.

\ex Δες τον τρίτο τρόπο, του διαβήτη, για τη διχοτόμο. Σχεδίασε γωνία $70^\circ$ και κατασκεύασε τη διχοτόμο της χωρίς μοιρογνωμόνιο. Έλεγξέ την με μοιρογνωμόνιο.

\zoni{Για το σπίτι}{με χαρτί}

\ex Κόψε από χαρτί μια οξεία και μια αμβλεία γωνία και βρες τη διχοτόμο τους με δίπλωση. Μέτρησε με μοιρογνωμόνιο ότι οι δύο μισές είναι ίσες.

\ex Στο σπίτι σου βρες τρία πράγματα με γωνίες (π.χ. ένα βιβλίο ανοιχτό, ένα ψαλίδι) και πες πού θα περνούσε η διχοτόμος τους.

\paragrafos{Ιδιότητες των τετραπλεύρων}{Β.3.4, «Ιδιότητες παραλληλογράμμου – ορθογωνίου – ρόμβου – τετραγώνου – τραπεζίου – ισοσκελούς τραπεζίου»}

\begin{anagke}
Ο Νίκος έχει ένα μεταλλικό πλαίσιο παραθύρου σε σχήμα ορθογωνίου και θέλει να βεβαιωθεί ότι είναι σωστά στερεωμένο, χωρίς γωνιόμετρο. Η Αγγελική του δείχνει ένα κόλπο: «Μέτρα τις δύο διαγώνιες. Αν είναι ίσες, το πλαίσιο είναι σωστό ορθογώνιο!». Ο Βασίλης δοκιμάζει σε ρόμβο από χαρτόνι: οι διαγώνιες του δεν είναι ίσες, αλλά τέμνονται κάθετα και διχοτομούνται.

Η Μυρτώ διπλώνει ένα χάρτινο παραλληλόγραμμο. Δεν βρίσκει άξονα συμμετρίας, αλλά αν το περιστρέψει μισή στροφή γύρω από το σημείο τομής των διαγωνίων, πέφτει ακριβώς πάνω στον εαυτό του. Η Πολυξένη: «Το κέντρο συμμετρίας του είναι εκεί!».

Σε αυτή την παράγραφο μαθαίνουμε τις ιδιότητες κάθε τετραπλεύρου: πλευρές, γωνίες, διαγώνιες, άξονες και κέντρο συμμετρίας.
\end{anagke}

\begin{lexi}{διαγώνιος, διχοτομούνται}{ονόματα}
Η \lexiword{διαγώνιος} ενώνει δύο μη διαδοχικές κορυφές (διά γωνίας). Δύο τμήματα \lexiword{διχοτομούνται} όταν το καθένα περνά από το μέσο του άλλου.
\end{lexi}

\begin{metafrasi}
Οι διαγώνιες περνούν από το μέσο η μία της άλλης. & διχοτομούνται\\
Οι απέναντι πλευρές είναι ίσες. & $\mathrm{ΑΒ}=\mathrm{ΔΓ}$, $\mathrm{ΑΔ}=\mathrm{ΒΓ}$\\
Οι απέναντι γωνίες είναι ίσες. & $\hat{\mathrm{Α}}=\hat{\Gamma}$, $\hat{\mathrm{Β}}=\hat{\mathrm{Δ}}$\\
Οι διαγώνιες σχηματίζουν ορθή γωνία. & κάθετες: $\mathrm{ΑΓ}\perp\mathrm{ΒΔ}$\\
\end{metafrasi}

\tmima{Η θεωρία}

\begin{kanonas}{παραλληλόγραμμο}
Σε κάθε παραλληλόγραμμο:
\begin{itemize}[leftmargin=1.2em,itemsep=0.3ex,topsep=2pt]
\item το σημείο τομής των διαγωνίων είναι \textbf{κέντρο συμμετρίας},
\item οι διαγώνιες \textbf{διχοτομούνται},
\item οι απέναντι πλευρές είναι \textbf{ίσες},
\item οι απέναντι γωνίες είναι \textbf{ίσες}.
\end{itemize}
\end{kanonas}

\begin{kanonas}{ορθογώνιο}
Εκτός από τις ιδιότητες του παραλληλογράμμου: οι \textbf{διαγώνιες είναι ίσες}, και οι μεσοκάθετοι των πλευρών του είναι \textbf{άξονες συμμετρίας} (δύο άξονες).
\end{kanonas}

\begin{kanonas}{ρόμβος}
Εκτός από τις ιδιότητες του παραλληλογράμμου: οι ευθείες των διαγωνίων είναι \textbf{άξονες συμμετρίας}, οι διαγώνιες είναι \textbf{κάθετες} και \textbf{διχοτόμοι των γωνιών} του.
\end{kanonas}

\begin{kanonas}{τετράγωνο}
Έχει όλες τις ιδιότητες του ορθογωνίου και του ρόμβου: οι διαγώνιες είναι ίσες, κάθετες, διχοτομούνται, διχοτομούν τις γωνίες ($45^\circ$ η καθεμία). Έχει \textbf{$4$ άξονες συμμετρίας} (δύο διαγώνιες και δύο μεσοκάθετοι των πλευρών) και κέντρο συμμετρίας.
\end{kanonas}

\begin{kanonas}{ισοσκελές τραπέζιο}
Η ευθεία που περνά από τα μέσα των βάσεων είναι \textbf{άξονας συμμετρίας} και μεσοκάθετος των βάσεων. Οι \textbf{προσκείμενες στη βάση γωνίες} είναι ίσες.
\end{kanonas}

\begin{center}\small
\begin{tabular}{lcccc}
\toprule
Ιδιότητα & Παραλληλόγραμμο & Ορθογώνιο & Ρόμβος & Τετράγωνο\\
\midrule
Απέναντι πλευρές ίσες & ναι & ναι & ναι & ναι\\
Απέναντι γωνίες ίσες & ναι & ναι & ναι & ναι\\
Διαγώνιες διχοτομούνται & ναι & ναι & ναι & ναι\\
Διαγώνιες ίσες & όχι & ναι & όχι & ναι\\
Διαγώνιες κάθετες & όχι & όχι & ναι & ναι\\
Άξονες συμμετρίας & $0$ & $2$ & $2$ & $4$\\
Κέντρο συμμετρίας & ναι & ναι & ναι & ναι\\
\bottomrule
\end{tabular}
\end{center}

\begin{prosoxi}
\begin{itemize}[leftmargin=1.2em,itemsep=0.3ex,topsep=2pt]
\item Σε ένα (πλάγιο) παραλληλόγραμμο οι διαγώνιες δεν είναι ούτε ίσες ούτε κάθετες, και δεν υπάρχει άξονας συμμετρίας.
\item Ό,τι ισχύει για το παραλληλόγραμμο ισχύει και για το ορθογώνιο, τον ρόμβο και το τετράγωνο, γιατί είναι παραλληλόγραμμα.
\item Οι διαγώνιες του ισοσκελούς τραπεζίου είναι ίσες, αλλά δεν διχοτομούνται.
\end{itemize}
\end{prosoxi}

\tmima{Λυμένα παραδείγματα}

\paradeigma{Το ορθογώνιο παράθυρο}
\tixero Ορθογώνιο παράθυρο με διαστάσεις $90$ cm και $120$ cm.\par
\tithelo Η διαγώνιος και πόσο απέχει το κέντρο συμμετρίας από μια κορυφή.\par
\possindeo Από το πυθαγόρειο: $\mathrm{ΑΓ}=\sqrt{90^2+120^2}=\sqrt{8100+14400}=150$ cm. Οι διαγώνιες διχοτομούνται και είναι ίσες, άρα το κέντρο απέχει $150:2=75$ cm από κάθε κορυφή.\par
\apantisi $150$ cm και $75$ cm.\par
\elegxos Οι τέσσερις αποστάσεις από το κέντρο είναι ίσες.

\paradeigma{Οι διαγώνιες του ρόμβου}
\tixero Ρόμβος με διαγώνιες $10$ και $24$ cm.\par
\tithelo Η πλευρά του.\par
\possindeo Οι διαγώνιες διχοτομούνται και είναι κάθετες: σχηματίζουν ορθογώνιο τρίγωνο με κάθετες $5$ και $12$. Η πλευρά είναι $\sqrt{5^2+12^2}=13$ cm.\par
\apantisi $13$ cm (περίμετρος $52$ cm).\par
\elegxos $5^2+12^2=25+144=169=13^2$.

\paradeigma{Οι γωνίες του παραλληλογράμμου}
\tixero Παραλληλόγραμμο με $\hat{\mathrm{Α}}=65^\circ$.\par
\tithelo Οι άλλες γωνίες.\par
\possindeo Οι απέναντι είναι ίσες: $\hat{\Gamma}=65^\circ$. Οι διαδοχικές είναι παραπληρωματικές (εντός επί τα αυτά): $\hat{\mathrm{Β}}=180^\circ-65^\circ=115^\circ=\hat{\mathrm{Δ}}$.\par
\apantisi $65^\circ,115^\circ,65^\circ,115^\circ$.\par
\elegxos Το άθροισμα είναι $360^\circ$.

\paradeigma{Το ισοσκελές τραπέζιο}
\tixero Ισοσκελές τραπέζιο με μια γωνία στη μεγάλη βάση $70^\circ$.\par
\tithelo Οι άλλες γωνίες.\par
\possindeo Οι προσκείμενες στη βάση γωνίες είναι ίσες: δύο γωνίες $70^\circ$. Οι γωνίες στη μικρή βάση είναι παραπληρωματικές των πλαγίων: $180^\circ-70^\circ=110^\circ$ η καθεμία.\par
\apantisi $70^\circ,70^\circ,110^\circ,110^\circ$.\par
\elegxos $70+70+110+110=360$.

\tmima{Ασκήσεις}

\zoni{Το κατάλαβα}{πλευρές, γωνίες, διαγώνιοι}

\ex Παραλληλόγραμμο έχει πλευρές $8$ cm και $5$ cm. Βρες τις άλλες δύο και την περίμετρο.

\ex Παραλληλόγραμμο έχει γωνία $50^\circ$. Βρες τις άλλες τρεις γωνίες.

\ex Οι διαγώνιες ορθογωνίου είναι $13$ cm. Πόσο απέχει το κέντρο από κάθε κορυφή;

\ex Ρόμβος έχει διαγώνιες $12$ και $16$ cm. Βρες την πλευρά και την περίμετρο.

\ex Πόσους άξονες συμμετρίας έχει το τετράγωνο; Σχεδίασέ τους.

\ex Σχεδίασε ορθογώνιο ΑΒΓΔ με διάμετρο μια διαγώνιο και δες ότι ο κύκλος περνά από όλες τις κορυφές.

\ex Σωστό ή λάθος; α) Οι διαγώνιες του ορθογωνίου είναι ίσες. \quad β) Οι διαγώνιες του ρόμβου είναι ίσες. \quad γ) Οι διαγώνιες του παραλληλογράμμου διχοτομούνται. \quad δ) Το ορθογώνιο έχει $4$ άξονες συμμετρίας.

\zoni{Το χρησιμοποιώ}{κατασκευές και μετρήσεις}

\ex \textbf{Το πλαίσιο.} Ορθογώνιο πλαίσιο είναι $60$ cm $\times$ $80$ cm. Πόση πρέπει να είναι η διαγώνιος για να είναι σωστό ορθογώνιο;

\ex \textbf{Ο χαρταετός.} Ένας χαρταετός σε σχήμα ρόμβου έχει διαγώνιες $60$ και $80$ cm. Πόση είναι η πλευρά του και πόση η περίμετρός του;

\ex \textbf{Το τραπεζάκι.} Ισοσκελές τραπέζιο έχει γωνία $65^\circ$ στη μεγάλη βάση. Βρες τις άλλες γωνίες του.

\ex \textbf{Το οικόπεδο.} Ένα παραλληλόγραμμο οικόπεδο έχει γωνία $72^\circ$. Βρες τις υπόλοιπες γωνίες.

\ex \textbf{Ο γνώμονας.} Ένας τεχνίτης ελέγχει ότι ένα τραπέζι είναι ορθογώνιο μετρώντας τις διαγώνιες. Αρκεί αυτό αν το τραπέζι είναι σίγουρα παραλληλόγραμμο;

\zoni{Το σκέφτομαι}{γιατί συμβαίνει}

\ex Εξήγησε γιατί οι απέναντι πλευρές παραλληλογράμμου είναι ίσες. (Βοήθεια: δες το κέντρο συμμετρίας.)

\ex Σε ποια παραλληλόγραμμα οι διαγώνιες είναι διχοτόμοι των γωνιών; Εξήγησε.

\ex Βρες τους άξονες συμμετρίας του ισοσκελούς τραπεζίου. Εξήγησε γιατί οι προσκείμενες γωνίες βάσης είναι ίσες.

\ex Ποιο σχήμα προκύπτει αν ενώσουμε διαδοχικά τα μέσα των πλευρών ενός ορθογωνίου; (Ένα πείραμα με χαρτί.)

\zoni{Για το σπίτι}{πείραμα με χαρτί}

\ex Κόψε από χαρτί ένα παραλληλόγραμμο και τσάκισέ το στις διαγώνιες. Τι διαπιστώνεις;

\ex Μέτρησε τις διαγώνιες μιας πόρτας ή ενός παραθύρου και έλεγξε αν είναι ίσες.

\paragrafos{Ιδιότητες ισοσκελούς και ισόπλευρου τριγώνου}{Β.3.5, «Ιδιότητες ισοσκελούς και ισόπλευρου τριγώνου»}

\begin{anagke}
Η Αγγελική σχεδιάζει στο τετράδιό της ένα έλατο με τρίγωνα: κάθε τρίγωνο έχει δύο ίσες πλευρές. Τα κόβει με το ψαλίδι και τα διπλώνει στη μέση, και παρατηρεί ότι οι δύο μισές πέφτουν ακριβώς η μία πάνω στην άλλη. «Και οι γωνίες στη βάση είναι ίσες;», ρωτά η Ζωή. «Γιατί να είναι;» Η Πολυξένη το γράφει στο τετράδιό της με κωδικούς, γιατί της αρέσουν τα μοτίβα: ισοσκελές $\to$ ίσες γωνίες βάσης. Ο Ηλίας το μετρά σε τρία διαφορετικά ισοσκελή: τρεις φορές βγαίνουν ίσες.

Η κυρία Αλεξάνδρα ρωτά: «Η πτυχή του χαρτιού είναι και ύψος του τριγώνου; Και διχοτόμος; Και διάμεσος;». Η παρέα δοκιμάζει με ισόπλευρο τρίγωνο και βρίσκει ότι δεν υπάρχει μία πτυχή, αλλά τρεις.

Σε αυτή την παράγραφο μαθαίνουμε γιατί το ισοσκελές τρίγωνο έχει άξονα συμμετρίας, ποιες γωνίες του είναι ίσες, και ότι σε αυτό η διάμεσος, το ύψος και η διχοτόμος προς τη βάση είναι το ίδιο τμήμα. Θα δούμε και τις αντίστοιχες ιδιότητες του ισόπλευρου τριγώνου.
\end{anagke}

\begin{lexi}{βάση, προσκείμενες γωνίες, ταυτίζονται}{ονόματα και ιδιότητες}
Η \lexiword{βάση} του ισοσκελούς είναι η πλευρά που δεν είναι από τις δύο ίσες· στηρίζεται πάνω της σαν στο έδαφος. \lexiword{Προσκείμενες} (προς + κείμαι) λέμε τις γωνίες που «κείτονται πάνω» στην ίδια πλευρά, στα άκρα της. \lexiword{Ταυτίζονται} σημαίνει ότι δύο πράγματα είναι ένα και το ίδιο: δύο τμήματα που ταυτίζονται δεν είναι δύο, είναι ένα.
\end{lexi}

\begin{metafrasi}
Το ισοσκελές τρίγωνο με ίσες πλευρές $\mathrm{ΑΒ}$ και $\mathrm{ΑΓ}$. & $\mathrm{ΑΒ}=\mathrm{ΑΓ}$, βάση η $\mathrm{ΒΓ}$\\
Οι γωνίες στη βάση είναι ίσες. & $\hat{\mathrm{Β}}=\hat{\Gamma}$\\
Η διάμεσος $\mathrm{ΑΜ}$ προς τη βάση είναι και ύψος. & $\mathrm{ΒΜ}=\mathrm{ΜΓ}$ και $\mathrm{ΑΜ}\perp\mathrm{ΒΓ}$\\
Η διάμεσος $\mathrm{ΑΜ}$ είναι και διχοτόμος. & $\hat{\mathrm{Α}}_1=\hat{\mathrm{Α}}_2$\\
Το ισόπλευρο τρίγωνο. & $\mathrm{ΑΒ}=\mathrm{ΒΓ}=\mathrm{ΓΑ}$\\
\end{metafrasi}

\tmima{Η θεωρία}

\tmimasmall{Ισοσκελές τρίγωνο}

Ένα ισοσκελές τρίγωνο $\mathrm{ΑΒΓ}$ έχει $\mathrm{ΑΒ}=\mathrm{ΑΓ}$. Το $\mathrm{Α}$ είναι η \textbf{κορυφή} και η $\mathrm{ΒΓ}$ η \textbf{βάση} του. Αν πάρουμε το μέσο $\mathrm{Μ}$ της βάσης και φέρουμε την $\mathrm{ΑΜ}$, διπλώνοντας το σχήμα κατά την $\mathrm{ΑΜ}$ το $\mathrm{Β}$ πέφτει πάνω στο $\Gamma$.

\begin{center}
\begin{tikzpicture}[scale=0.95]
\draw[thick] (0,0)--(4,0)--(2,3)--cycle;
\draw[gblue,thick,dashed] (2,3.4)--(2,-0.4);
\node[below left] at (0,0) {$\mathrm{Β}$}; \node[below right] at (4,0) {$\Gamma$}; \node[above right] at (2,3) {$\mathrm{Α}$}; \node[below right] at (2,0) {$\mathrm{Μ}$};
\draw (2,0.3)--(2.3,0.3)--(2.3,0);
\draw (0.9,1.5)--(1.1,1.4); \draw (2.9,1.4)--(3.1,1.5); \draw (1,-0.1)--(1,0.1); \draw (3,-0.1)--(3,0.1);
\draw (0.6,0) arc(0:56:0.6); \draw (3.4,0) arc(180:124:0.6);
\end{tikzpicture}
\end{center}

\begin{kanonas}{ιδιότητες ισοσκελούς τριγώνου}
Σε ισοσκελές τρίγωνο $\mathrm{ΑΒΓ}$ με $\mathrm{ΑΒ}=\mathrm{ΑΓ}$:
\begin{itemize}[leftmargin=1.2em,itemsep=0.3ex,topsep=2pt]
\item Η ευθεία της διαμέσου $\mathrm{ΑΜ}$ προς τη βάση είναι \textbf{άξονας συμμετρίας} του τριγώνου.
\item Οι προσκείμενες στη βάση γωνίες είναι \textbf{ίσες}: $\hat{\mathrm{Β}}=\hat{\Gamma}$.
\item Η διάμεσος, το ύψος και η διχοτόμος προς τη βάση \textbf{ταυτίζονται}: είναι το ίδιο τμήμα $\mathrm{ΑΜ}$.
\end{itemize}
\end{kanonas}

Γιατί ισχύουν. Στη συμμετρία ως προς τον άξονα $\mathrm{ΑΜ}$ το $\mathrm{Β}$ γίνεται $\Gamma$, και το $\mathrm{Α}$ μένει στη θέση του. Έτσι η γωνία $\hat{\mathrm{Β}}$ πέφτει πάνω στη $\hat{\Gamma}$ και είναι ίσες. Η $\hat{\mathrm{Α}}_1$ πέφτει πάνω στην $\hat{\mathrm{Α}}_2$, άρα η $\mathrm{ΑΜ}$ είναι διχοτόμος. Και η γωνία $\mathrm{ΑΜ}\hat{\mathrm{Β}}$ πέφτει πάνω στην $\mathrm{ΑΜ}\hat{\Gamma}$, άρα είναι ίσες και μαζί κάνουν ευθεία γωνία: καθεμία είναι $90^\circ$, δηλαδή η $\mathrm{ΑΜ}$ είναι και ύψος.

Η ιδιότητα αυτή συνδέεται με τη μεσοκάθετο: αφού $\mathrm{ΑΒ}=\mathrm{ΑΓ}$, το $\mathrm{Α}$ ισαπέχει από τα $\mathrm{Β}$ και $\Gamma$, άρα βρίσκεται στη μεσοκάθετο της $\mathrm{ΒΓ}$. Και η $\mathrm{ΑΜ}$ είναι ακριβώς αυτή η μεσοκάθετος.

\tmimasmall{Ισόπλευρο τρίγωνο}

Το ισόπλευρο τρίγωνο έχει και τις τρεις πλευρές του ίσες. Αφού κάθε πλευρά του μπορεί να παίξει τον ρόλο της βάσης, όσα ισχύουν για το ισοσκελές ισχύουν και για καθεμία από τις τρεις πλευρές.

\begin{center}
\begin{tikzpicture}[scale=0.95]
\draw[thick] (0,0)--(4,0)--(2,3.464)--cycle;
\draw[gblue,dashed] (2,3.9)--(2,-0.4); \draw[gblue,dashed] (4,0)--(0.5,2.598)--(-0.3,3.17); \draw[gblue,dashed] (0,0)--(3.5,2.598)--(4.3,3.17);
\node[below left] at (0,0) {$\mathrm{Β}$}; \node[below right] at (4,0) {$\Gamma$}; \node[above] at (2,3.9) {$\mathrm{Α}$};
\end{tikzpicture}
\end{center}

\begin{kanonas}{ιδιότητες ισόπλευρου τριγώνου}
\begin{itemize}[leftmargin=1.2em,itemsep=0.3ex,topsep=2pt]
\item Οι ευθείες των τριών διαμέσων είναι \textbf{άξονες συμμετρίας} του τριγώνου: έχει τρεις άξονες συμμετρίας.
\item Όλες οι γωνίες του είναι \textbf{ίσες}. Θα αποδείξουμε στην 3.9 ότι κάθε μία είναι $60^\circ$.
\item Η διάμεσος, το ύψος και η διχοτόμος που αντιστοιχούν σε οποιαδήποτε πλευρά \textbf{ταυτίζονται}.
\end{itemize}
\end{kanonas}

\begin{prosoxi}
\begin{itemize}[leftmargin=1.2em,itemsep=0.3ex,topsep=2pt]
\item Οι ίσες γωνίες του ισοσκελούς είναι οι δύο γωνίες της \emph{βάσης}. Η γωνία της κορυφής είναι γενικά διαφορετική από αυτές.
\item Στο ισοσκελές η διάμεσος προς τη βάση ταυτίζεται με το ύψος και τη διχοτόμο. Η διάμεσος προς μία από τις ίσες πλευρές δεν είναι κατά κανόνα ύψος ούτε διχοτόμος.
\item Το ισοσκελές τρίγωνο που δεν είναι ισόπλευρο έχει έναν μόνο άξονα συμμετρίας, ενώ το ισόπλευρο έχει τρεις.
\item Η ιδιότητα $\hat{\mathrm{Β}}=\hat{\Gamma}$ ισχύει όταν $\mathrm{ΑΒ}=\mathrm{ΑΓ}$. Σε ένα σκαληνό τρίγωνο οι γωνίες δεν είναι ίσες, όσο κι αν στο σχέδιο μοιάζουν.
\end{itemize}
\end{prosoxi}

\tmima{Λυμένα παραδείγματα}

\paradeigma{Οι γωνίες της βάσης}
\tixero Το τρίγωνο $\mathrm{ΑΒΓ}$ είναι ισοσκελές με $\mathrm{ΑΒ}=\mathrm{ΑΓ}$ και $\hat{\mathrm{Β}}=70^\circ$.\par
\tithelo Η γωνία $\hat{\Gamma}$.\par
\possindeo Οι προσκείμενες στη βάση $\mathrm{ΒΓ}$ γωνίες του ισοσκελούς είναι ίσες.\par
\apantisi $\hat{\Gamma}=\hat{\mathrm{Β}}=70^\circ$.\par
\elegxos Σε σχέδιο με μοιρογνωμόνιο μετράμε και τις δύο γωνίες της βάσης: βγαίνουν $70^\circ$.

\paradeigma{Η διάμεσος του ισοσκελούς}
\tixero Το τρίγωνο $\mathrm{ΑΒΓ}$ είναι ισοσκελές με $\mathrm{ΑΒ}=\mathrm{ΑΓ}$. Η $\mathrm{ΑΜ}$ είναι διάμεσος, $\mathrm{ΒΓ}=10$ cm και $\hat{\mathrm{Α}}=50^\circ$.\par
\tithelo Τα $\mathrm{ΒΜ}$, $\mathrm{ΑΜ}\hat{\mathrm{Β}}$ και $\mathrm{ΒΑΜ}$.\par
\possindeo Η διάμεσος προς τη βάση είναι και ύψος και διχοτόμος. Το $\mathrm{Μ}$ είναι μέσο της $\mathrm{ΒΓ}$.\par
\apantisi $\mathrm{ΒΜ}=10:2=5$ cm. Είναι ύψος: $\mathrm{ΑΜ}\hat{\mathrm{Β}}=90^\circ$. Είναι διχοτόμος: $\mathrm{ΒΑΜ}=50^\circ:2=25^\circ$.\par
\elegxos Οι δύο μισές γωνίες $25^\circ+25^\circ=50^\circ$.

\paradeigma{Οι άξονες}
\tixero Ένα ισόπλευρο τρίγωνο πλευράς $6$ cm.\par
\tithelo Πόσους άξονες συμμετρίας έχει και ποιοι είναι;\par
\possindeo Κάθε διάμεσος είναι άξονας συμμετρίας και έχει τρεις διαμέσους.\par
\apantisi Τρεις άξονες: οι ευθείες των τριών διαμέσων, που περνούν από το ίδιο σημείο.\par
\elegxos Αν διπλώσω το τρίγωνο κατά οποιονδήποτε από τους τρεις άξονες, οι δύο μισές συμπίπτουν.

\paradeigma{Ποιο τρίγωνο διπλώνεται;}
\tixero Η Ρήνη διπλώνει ένα χάρτινο τρίγωνο ώστε το $\mathrm{Β}$ να πέσει πάνω στο $\Gamma$, και οι δύο μισές συμπίπτουν ακριβώς.\par
\tithelo Τι είδους τρίγωνο είναι;\par
\possindeo Η πτυχή είναι άξονας συμμετρίας στον οποίο το $\mathrm{Β}$ γίνεται $\Gamma$ και το $\mathrm{Α}$ μένει στη θέση του, άρα $\mathrm{ΑΒ}=\mathrm{ΑΓ}$.\par
\apantisi Ισοσκελές, με κορυφή το $\mathrm{Α}$.\par
\elegxos Η πτυχή περνά από το $\mathrm{Α}$ και από το μέσο της $\mathrm{ΒΓ}$.

\tmima{Ασκήσεις}

\zoni{Το κατάλαβα}{ισοσκελές και ισόπλευρο}

\ex Σωστό ή λάθος; α) Οι προσκείμενες στη βάση γωνίες ισοσκελούς τριγώνου είναι ίσες. \quad β) Σε κάθε ισοσκελές η διάμεσος προς τη βάση είναι και διχοτόμος. \quad γ) Σε κάθε ισόπλευρο τρίγωνο οι ευθείες των υψών είναι και μεσοκάθετοι των πλευρών του. \quad δ) Το ισοσκελές τρίγωνο έχει πάντα τρεις άξονες συμμετρίας. \quad ε) Σε κάθε ισοσκελές η διάμεσος προς τη βάση είναι άξονας συμμετρίας.

\ex Το τρίγωνο $\mathrm{ΑΒΓ}$ είναι ισοσκελές με $\mathrm{ΑΒ}=\mathrm{ΑΓ}$ και η $\mathrm{ΑΜ}$ είναι ύψος. Αν $\mathrm{ΒΓ}=8$ cm και $\hat{\mathrm{Α}}=50^\circ$, βρες τα $\mathrm{ΒΜ}$, $\mathrm{ΑΜ}\hat{\mathrm{Β}}$ και $\mathrm{ΒΑΜ}$.

\ex Στο ισοσκελές τρίγωνο $\mathrm{ΑΒΓ}$ με κορυφή το $\mathrm{Α}$ δίνονται: α) $\hat{\mathrm{Β}}=70^\circ$, β) $\hat{\Gamma}=48^\circ$. Βρες σε κάθε περίπτωση την άλλη γωνία της βάσης και δικαιολόγησε.

\ex Στο ισόπλευρο τρίγωνο $\mathrm{ΑΒΓ}$ η $\mathrm{ΑΔ}$ είναι διάμεσος και $\mathrm{ΑΒ}=6$ cm. Βρες το $\mathrm{ΒΓ}$, το $\mathrm{ΒΔ}$ και τη γωνία $\mathrm{ΑΔ}\hat{\mathrm{Β}}$. Πόσο είναι η γωνία $\mathrm{ΒΑΔ}$ αν ξέρουμε ότι $\hat{\mathrm{Α}}=60^\circ$;

\ex Κατασκεύασε με κανόνα και διαβήτη ένα ισοσκελές τρίγωνο με $\mathrm{ΑΒ}=\mathrm{ΑΓ}=5$ cm και $\mathrm{ΒΓ}=6$ cm. Μέτρησε τις γωνίες της βάσης και βρες τον άξονα συμμετρίας με δίπλωση.

\ex Κατασκεύασε με κανόνα και διαβήτη ένα ισόπλευρο τρίγωνο με πλευρά $5$ cm. Φέρε τους τρεις άξονες συμμετρίας και μέτρησε τις γωνίες του.

\ex Πόσους άξονες συμμετρίας έχει: α) ένα ισοσκελές τρίγωνο που δεν είναι ισόπλευρο, β) ένα ισόπλευρο, γ) ένα σκαληνό;

\zoni{Το χρησιμοποιώ}{στέγες, πανιά, πινακίδες}

\ex \textbf{Η στέγη.} Η στέγη του σπιτιού της γιαγιάς Μυρτώς έχει σχήμα ισοσκελούς τριγώνου με γωνία βάσης $35^\circ$. Πόση είναι η άλλη γωνία της βάσης;

\ex \textbf{Το πανί.} Ένα ισοσκελές τριγωνικό πανί έχει βάση $12$ cm. Η διάμεσος από την κορυφή προς τη βάση έχει μήκος $8$ cm. Πόσο απέχει το πόδι της από κάθε άκρο της βάσης; Τι γωνία σχηματίζει με τη βάση;

\ex \textbf{Η σημαία.} Στο ισοσκελές τρίγωνο $\mathrm{ΑΒΓ}$ με $\mathrm{ΑΒ}=\mathrm{ΑΓ}$ η $\mathrm{ΑΔ}$ είναι η διχοτόμος της γωνίας $\hat{\mathrm{Α}}=64^\circ$. Βρες τη γωνία $\mathrm{ΒΑΔ}$ και τη γωνία $\mathrm{ΑΔ}\hat{\mathrm{Β}}$.

\ex \textbf{Η πινακίδα.} Η πινακίδα «Δώσε προτεραιότητα» έχει σχήμα ισόπλευρου τριγώνου πλευράς $40$ cm. Βρες την περίμετρο και πόσους άξονες συμμετρίας έχει.

\ex \textbf{Η δίπλωση του Στέλιου.} Ο Στέλιος διπλώνει ένα τρίγωνο και οι δύο μισές συμπίπτουν. Λέει: «Άρα το τρίγωνο είναι ισόπλευρο!». Έχει δίκιο; Τι ακριβώς είναι;

\zoni{Το σκέφτομαι}{γιατί συμβαίνει}

\ex Εξήγησε γιατί οι γωνίες στη βάση ενός ισοσκελούς είναι ίσες, με βάση τη δίπλωση.

\ex Γιατί η ευθεία $\mathrm{ΑΜ}$ στο ισοσκελές τρίγωνο είναι η μεσοκάθετος της βάσης $\mathrm{ΒΓ}$; (Σκέψου το $\mathrm{Α}$ και τις αποστάσεις του από τα $\mathrm{Β}$, $\Gamma$.)

\ex Σε ισοσκελές τρίγωνο $\mathrm{ΑΒΓ}$ ($\mathrm{ΑΒ}=\mathrm{ΑΓ}$) φέρε τη διάμεσο προς την πλευρά $\mathrm{ΑΓ}$. Είναι και ύψος; Είναι και διχοτόμος; Έλεγξε με σχέδιο.

\ex Σχεδίασε ένα ορθογώνιο ισοσκελές τρίγωνο με κάθετες πλευρές $4$ cm. Μέτρησε τις δύο άλλες γωνίες. Τι παρατηρείς;

\zoni{Για το σπίτι}{με χαρτί}

\ex Κόψε ένα ισοσκελές τρίγωνο από χαρτί, δίπλωσέ το στη μέση και μέτρησε τις γωνίες της βάσης.

\ex Κόψε ένα ισόπλευρο τρίγωνο από χαρτί και βρες με δίπλωση τους τρεις άξονες συμμετρίας του. Πού τέμνονται;

\paragrafos{Κατακορυφήν, συμπληρωματικές και παραπληρωματικές γωνίες}{Β.3.6, «Κατακορυφήν, συμπληρωματικές και παραπληρωματικές γωνίες»}

\begin{anagke}
Δύο μονοπάτια στον κήπο της θείας Μαρίας διασταυρώνονται σαν ένα μεγάλο «Χ». Η Αγγελική μετράει με το μοιρογνωμόνιο τη μία γωνία και βρίσκει $48^\circ$. «Πόσες είναι οι άλλες τρεις;», ρωτά. Ο Νίκος μετράει όλες: $48^\circ$, $132^\circ$, $48^\circ$, $132^\circ$. «Οι απέναντι είναι ίσες! Και οι διπλανές κάνουν μαζί $180^\circ$». Η Πολυξένη ρωτά: «Είναι τυχαίο; Αν στρέψω τη μία ευθεία, θα αλλάξει;». Δοκιμάζουν με άλλη γωνία, $70^\circ$, και βρίσκουν $70^\circ$, $110^\circ$, $70^\circ$, $110^\circ$.

Η Ρήνη θέλει να φυτέψει τριαντάφυλλα σε μια γωνιά του φράχτη που είναι ορθή. Ένα παρτέρι με λουλούδια παίρνει γωνία $35^\circ$ και ρωτά πόση γωνία μένει για τα τριαντάφυλλα. Ο Στέλιος υπολογίζει: «Εδώ το άθροισμα είναι $90^\circ$, όχι $180^\circ$».

Έτσι η παρέα συναντά τρεις έννοιες: τις γωνίες που αθροίζουν $90^\circ$ (συμπληρωματικές), εκείνες που αθροίζουν $180^\circ$ (παραπληρωματικές) και αυτές που βρίσκονται απέναντι όταν δύο ευθείες τέμνονται (κατακορυφήν). Σε αυτή την παράγραφο θα δούμε και γιατί οι κατακορυφήν είναι ίσες, χωρίς μοιρογνωμόνιο.
\end{anagke}

\begin{lexi}{εφεξής, συμπληρωματική, παραπληρωματική, κατακορυφήν}{γωνίες που σχετίζονται}
Δύο γωνίες είναι \lexiword{εφεξής} όταν η μία έρχεται «από κοντά» στην άλλη. Η \lexiword{συμπληρωματική} γωνία «συμπληρώνει» μια γωνία μέχρι την ορθή. Η \lexiword{παραπληρωματική} «παρα-πληρώνει», δηλαδή συμπληρώνει ως τη μισή στροφή, τις $180^\circ$. Οι \lexiword{κατακορυφήν} γωνίες είναι αυτές που βρίσκονται «κατά κορυφήν», δηλαδή απέναντι, με κοινή κορυφή.
\end{lexi}

\begin{metafrasi}
Δύο γωνίες με κοινή κορυφή και κοινή πλευρά, χωρίς άλλο κοινό σημείο. & εφεξής\\
Δύο γωνίες με άθροισμα $90^\circ$. & συμπληρωματικές: $\hat{\alpha}+\hat{\beta}=90^\circ$\\
Δύο γωνίες με άθροισμα $180^\circ$. & παραπληρωματικές: $\hat{\alpha}+\hat{\beta}=180^\circ$\\
Οι απέναντι γωνίες δύο τεμνόμενων ευθειών. & κατακορυφήν: $\hat{\alpha}=\hat{\gamma}$\\
Η συμπληρωματική της $35^\circ$. & $90^\circ-35^\circ=55^\circ$\\
Η παραπληρωματική της $130^\circ$. & $180^\circ-130^\circ=50^\circ$\\
\end{metafrasi}

\tmima{Η θεωρία}

\tmimasmall{Εφεξής και διαδοχικές γωνίες}

\begin{orismos}{εφεξής γωνίες}
Δύο γωνίες λέγονται \textbf{εφεξής} όταν έχουν κοινή κορυφή, μία κοινή πλευρά και κανένα άλλο κοινό σημείο.
\end{orismos}

Μια σειρά από γωνίες που η καθεμία είναι εφεξής με την επόμενη λέγονται \textbf{διαδοχικές}. Για παράδειγμα στο σχήμα οι γωνίες $\mathrm{ΑΟΒ}$, $\mathrm{ΒΟΓ}$ και $\mathrm{ΓΟΔ}$ είναι διαδοχικές. Το άθροισμά τους είναι η γωνία $\mathrm{ΑΟΔ}$.

\begin{center}
\begin{tikzpicture}[scale=0.9]
\draw[thick] (0,0)--(4.5,0) node[right] {$\mathrm{Α}$};
\draw[thick] (0,0)--({4.2*cos(35)},{4.2*sin(35)}) node[above right] {$\mathrm{Β}$};
\draw[thick] (0,0)--({4.2*cos(80)},{4.2*sin(80)}) node[above] {$\Gamma$};
\draw[thick] (0,0)--({4.2*cos(130)},{4.2*sin(130)}) node[above left] {$\Delta$};
\fill (0,0) circle (2pt) node[below left] {$\mathrm{Ο}$};
\draw[gblue] (1.0,0) arc(0:35:1.0); \draw[gblue] (1.2,0.55) arc(25:80:1.2); \draw[gblue] (0.7,0.5) arc(35:130:0.7);
\end{tikzpicture}
\end{center}

\tmimasmall{Συμπληρωματικές και παραπληρωματικές}

\begin{orismos}{συμπληρωματικές και παραπληρωματικές}
Δύο γωνίες λέγονται \textbf{συμπληρωματικές} όταν το άθροισμά τους είναι $90^\circ$ και \textbf{παραπληρωματικές} όταν το άθροισμά τους είναι $180^\circ$. Κάθε μία από τις δύο λέγεται συμπληρωματική (ή παραπληρωματική) της άλλης.
\end{orismos}

Για παράδειγμα η συμπληρωματική της $35^\circ$ είναι η $55^\circ$ και η παραπληρωματική της $130^\circ$ είναι η $50^\circ$. Δύο γωνίες δεν χρειάζεται να είναι εφεξής για να είναι συμπληρωματικές ή παραπληρωματικές: αρκεί το άθροισμα των μέτρων τους.

\begin{center}
\begin{tikzpicture}[scale=0.8]
\draw[thick] (-0.3,0)--(3.5,0); \draw[thick] (0,-0.3)--(0,3.5); \draw[thick] (0,0)--({3.3*cos(35)},{3.3*sin(35)});
\draw[gblue] (0.8,0) arc(0:35:0.8); \draw[gwarmD] (0.6,0.4) arc(35:90:0.6); \draw (0.3,0)--(0.3,0.3)--(0,0.3);
\node[gblue] at (1.2,0.3) {\small $\hat{\alpha}$}; \node[gwarmD] at (0.4,1.1) {\small $\hat{\beta}$}; \node at (1.7,-0.7) {\small $\hat{\alpha}+\hat{\beta}=90^\circ$};
\begin{scope}[xshift=6.2cm]
\draw[thick] (-2.2,0)--(2.6,0); \draw[thick] (0,0)--({2.3*cos(60)},{2.3*sin(60)});
\draw[gblue] (0.7,0) arc(0:60:0.7); \draw[gwarmD] (-0.55,0) arc(180:60:0.55);
\node[gblue] at (1.05,0.3) {\small $\hat{\alpha}$}; \node[gwarmD] at (-0.3,0.75) {\small $\hat{\beta}$}; \node at (0.3,-0.7) {\small $\hat{\alpha}+\hat{\beta}=180^\circ$};
\end{scope}
\end{tikzpicture}
\end{center}

\tmimasmall{Κατακορυφήν γωνίες}

\begin{orismos}{κατακορυφήν γωνίες}
Δύο γωνίες λέγονται \textbf{κατακορυφήν} όταν έχουν κοινή κορυφή και οι πλευρές τους είναι αντικείμενες ημιευθείες.
\end{orismos}

Όταν δύο ευθείες τέμνονται σχηματίζονται τέσσερις γωνίες, και οι δύο ζεύγη απέναντι γωνιών είναι κατακορυφήν.

\begin{center}
\begin{tikzpicture}[scale=1.1]
\draw[thick,{Stealth}-{Stealth}] (-2.5,-1)--(2.5,1); \draw[thick,{Stealth}-{Stealth}] (-1.2,-1.8)--(1.2,1.8);
\fill (0,0) circle (2pt); \node[below right] at (0.05,0) {$\mathrm{Ο}$};
\node[gblue] at (1.2,0.15) {$\hat{\alpha}$}; \node[gblue] at (-1.2,-0.15) {$\hat{\gamma}$}; \node at (0.2,0.9) {$\hat{\beta}$}; \node at (-0.2,-0.9) {$\hat{\delta}$};
\end{tikzpicture}
\end{center}

\begin{idiotita}{κατακορυφήν γωνίες}
Οι κατακορυφήν γωνίες είναι \textbf{ίσες}.
\end{idiotita}

Γιατί: Η $\hat{\alpha}$ και η $\hat{\beta}$ είναι παραπληρωματικές ($\hat{\alpha}+\hat{\beta}=180^\circ$), και η $\hat{\gamma}$ και η $\hat{\beta}$ επίσης ($\hat{\gamma}+\hat{\beta}=180^\circ$). Άρα $\hat{\alpha}=180^\circ-\hat{\beta}=\hat{\gamma}$.

\begin{prosoxi}
\begin{itemize}[leftmargin=1.2em,itemsep=0.3ex,topsep=2pt]
\item Συμπληρωματικές: $90^\circ$. Παραπληρωματικές: $180^\circ$. Η «παραπληρωματική» είναι η μεγαλύτερη λέξη και το μεγαλύτερο άθροισμα.
\item Δεν υπάρχει συμπληρωματική γωνία για γωνία ίση ή μεγαλύτερη από $90^\circ$ (θα βγαίνε μηδέν ή αρνητικό).
\item Οι διπλανές γωνίες δύο τεμνόμενων ευθειών δεν είναι ίσες (εκτός αν είναι ορθές): είναι παραπληρωματικές.
\item Δύο ίσες γωνίες δεν είναι πάντα κατακορυφήν: αρκεί να έχουν ίδιο μέτρο.
\end{itemize}
\end{prosoxi}

\tmima{Λυμένα παραδείγματα}

\paradeigma{Συμπλήρωμα και παραπλήρωμα}
\tixero Οι γωνίες $35^\circ$ και $130^\circ$.\par
\tithelo Η συμπληρωματική της πρώτης και η παραπληρωματική της δεύτερης.\par
\possindeo Αφαιρώ από το $90^\circ$ και από το $180^\circ$ αντίστοιχα.\par
\apantisi $90^\circ-35^\circ=55^\circ$ και $180^\circ-130^\circ=50^\circ$.\par
\elegxos $35^\circ+55^\circ=90^\circ$ και $130^\circ+50^\circ=180^\circ$.

\paradeigma{Δύο τεμνόμενες ευθείες}
\tixero Δύο ευθείες τέμνονται και μία από τις γωνίες είναι $50^\circ$.\par
\tithelo Οι άλλες τρεις γωνίες.\par
\possindeo Η απέναντί της είναι κατακορυφήν, άρα $50^\circ$. Οι διπλανές της είναι παραπληρωματικές της, δηλαδή $180^\circ-50^\circ=130^\circ$, και η τέταρτη είναι κατακορυφήν της μίας από αυτές, άρα επίσης $130^\circ$.\par
\apantisi $50^\circ$, $130^\circ$ και $130^\circ$.\par
\elegxos Το άθροισμα και των τεσσάρων: $50^\circ+130^\circ+50^\circ+130^\circ=360^\circ$.

\paradeigma{Η συμπληρωματική με το διπλάσιο}
\tixero Μια γωνία είναι συμπληρωματική με μια άλλη που είναι διπλάσια από αυτήν.\par
\tithelo Οι δύο γωνίες.\par
\possindeo Αν η μία είναι $\omega$, η άλλη είναι $2\omega$ και το άθροισμά τους είναι $90^\circ$: $3\omega=90^\circ$, άρα $\omega=30^\circ$.\par
\apantisi $30^\circ$ και $60^\circ$.\par
\elegxos $30^\circ+60^\circ=90^\circ$ και η δεύτερη είναι διπλάσια της πρώτης.

\paradeigma{Όλες οι κατακορυφήν}
\tixero Δύο ευθείες τέμνονται και μία γωνία είναι $\hat{\alpha}=2x+10^\circ$ ... \par
\tithelo ...\par
\possindeo ...\par
\apantisi ...\par
\elegxos ...

\tmima{Ασκήσεις}

\zoni{Το κατάλαβα}{συμπληρωματικές, παραπληρωματικές, κατακορυφήν}

\ex Σωστό ή λάθος; α) Δύο γωνίες με άθροισμα $90^\circ$ λέγονται παραπληρωματικές. \quad β) Δύο κατακορυφήν γωνίες είναι πάντα ίσες. \quad γ) Η συμπληρωματική της μηδενικής γωνίας είναι η ορθή. \quad δ) Η παραπληρωματική μιας αμβλείας είναι πάντα οξεία. \quad ε) Δύο ίσες γωνίες είναι πάντα κατακορυφήν. \quad στ) Δύο εφεξής γωνίες έχουν κοινή κορυφή και μία κοινή πλευρά.

\ex Βρες τη συμπληρωματική των γωνιών: $35^\circ$, $62^\circ$, $8^\circ$, $89^\circ$.

\ex Βρες την παραπληρωματική των γωνιών: $130^\circ$, $90^\circ$, $5^\circ$, $111^\circ$.

\ex Δύο ευθείες τέμνονται και μία γωνία είναι $48^\circ$. Βρες τις άλλες τρεις γωνίες.

\ex Δύο ευθείες $x\mathrm{Ο}x'$ και $y\mathrm{Ο}y'$ τέμνονται στο $\mathrm{Ο}$. Γράψε όλα τα ζεύγη κατακορυφήν γωνιών και όλα τα ζεύγη παραπληρωματικών γωνιών που σχηματίζονται.

\ex Σχεδίασε μια γωνία $40^\circ$ και μετά τη συμπληρωματική της και την παραπληρωματική της.

\ex Τι είδους γωνία είναι: α) η παραπληρωματική μιας οξείας, β) η παραπληρωματική μιας ορθής, γ) η παραπληρωματική μιας αμβλείας, δ) η συμπληρωματική μιας οξείας;

\ex Οι γωνίες $\mathrm{ΑΟΒ}=30^\circ$, $\mathrm{ΒΟΓ}=45^\circ$ και $\mathrm{ΓΟΔ}=50^\circ$ είναι διαδοχικές. Πόση είναι η γωνία $\mathrm{ΑΟΔ}$; Είναι οξεία, ορθή ή αμβλεία;

\zoni{Το χρησιμοποιώ}{διασταυρώσεις, σκάλες, ψαλίδια}

\ex \textbf{Η διασταύρωση.} Δύο μονοπάτια του κήπου τέμνονται και μία γωνία είναι $70^\circ$. Πόσες μοίρες έχουν οι άλλες γωνίες; Πόσες διαφορετικές τιμές εμφανίζονται;

\ex \textbf{Τα τριαντάφυλλα.} Στη γωνιά του φράχτη, που είναι ορθή, φυτεύεται ένα παρτέρι με γωνία $35^\circ$ και το υπόλοιπο φυτεύεται με τριαντάφυλλα. Πόση γωνία μένει για τα τριαντάφυλλα;

\ex \textbf{Η σκάλα.} Μια σκάλα ακουμπά σε έναν τοίχο και σχηματίζει με το έδαφος γωνία $72^\circ$. Πόση γωνία σχηματίζει με τον τοίχο; (Ο τοίχος είναι κάθετος στο έδαφος.)

\ex \textbf{Το ψαλίδι.} Ο παππούς Γιάννης ανοίγει ένα ψαλίδι ώστε οι δύο λεπίδες να σχηματίζουν γωνία $40^\circ$. Πόση γωνία σχηματίζουν οι λαβές του;

\ex \textbf{Η πίτσα.} Μια γωνία είναι ίση με τα $\frac{2}{5}$ της ορθής. Πόσες μοίρες είναι; Βρες τη συμπληρωματική και την παραπληρωματική της.

\zoni{Το σκέφτομαι}{γιατί και αν}

\ex Εξήγησε γιατί οι κατακορυφήν γωνίες είναι ίσες, χρησιμοποιώντας τις παραπληρωματικές γωνίες.

\ex Δύο ίσες γωνίες είναι απαραίτητα κατακορυφήν; Δώσε ένα αντιπαράδειγμα.

\ex Βρες το μέτρο της γωνίας για την οποία ισχύει ότι: α) η συμπληρωματική της είναι κατά $20^\circ$ μεγαλύτερη από αυτήν, β) η παραπληρωματική της είναι κατά $40^\circ$ μικρότερη από αυτήν, γ) η συμπληρωματική της είναι διπλάσια από αυτήν, δ) η παραπληρωματική της είναι τριπλάσια από αυτήν.

\ex Υπάρχει γωνία ίση με τη συμπληρωματική της; Και γωνία ίση με την παραπληρωματική της; Βρες τες.

\zoni{Για το σπίτι}{ψάξε Χ}

\ex Σχεδίασε δύο τεμνόμενες ευθείες, μέτρησε και τις τέσσερις γωνίες και έλεγξε ότι οι κατακορυφήν είναι ίσες.

\ex Βρες στο σπίτι ένα ζευγάρι τεμνόμενων ευθειών (ψαλίδι, παράθυρο, αναδιπλούμενη καρέκλα) και μέτρησε τις γωνίες που σχηματίζουν.

\paragrafos{Σχετικές θέσεις ευθειών στο επίπεδο}{Β.3.7, «Σχετικές θέσεις ευθειών στο επίπεδο»}

\begin{anagke}
Ο παππούς Γιάννης δείχνει πώς θα φυτευτούν οι σειρές με τα καρότα στον κήπο: «Κάθε σειρά θα είναι ευθεία, όλες με την ίδια κατεύθυνση. Καμία δεν θα ακουμπά την άλλη, όσο κι αν τις συνεχίσεις». Η Αγγελική σκέφτεται τις γραμμές στην εθνική οδό και τις ράγες του τρένου: «Δεν συναντιούνται ποτέ, ούτε αν προχωρήσουμε χιλιόμετρα». Ο Ηλίας δείχνει δύο μονοπάτια του κήπου που διασταυρώνονται: «Αυτά ακουμπούν σε ένα σημείο».

«Μα αν συνεχίσεις ένα μονοπάτι που φαινομενικά δεν διασταυρώνεται, μπορεί να διασταυρωθεί αργότερα!», διαμαρτύρεται ο Στέλιος. Η Πολυξένη προτείνει ένα τεστ: «Ας φανταστούμε ότι οι ευθείες δεν τελειώνουν ποτέ, ούτε στο χαρτί. Τότε δύο ευθείες ή έχουν ένα κοινό σημείο ή δεν έχουν κανένα». Η κυρία Αλεξάνδρα προσθέτει μια ερώτηση: «Και πόσες παράλληλες στην ευθεία σας περνούν από ένα σημείο που βρίσκεται δίπλα της;».

Σε αυτή την παράγραφο ονομάζουμε τις δύο δυνατές σχετικές θέσεις δύο ευθειών του επιπέδου, τις τεμνόμενες και τις παράλληλες. Θα μάθουμε και το αίτημα του Ευκλείδη για τις παράλληλες, καθώς και πώς φέρνουμε παράλληλη από ένα σημείο.
\end{anagke}

\begin{lexi}{παράλληλες, τεμνόμενες}{παράλληλες, τεμνόμενες}
Οι \lexiword{παράλληλες} (παρ' άλληλες) πηγαίνουν η μία «δίπλα στην άλλη» χωρίς να συναντιούνται ποτέ. Οι \lexiword{τεμνόμενες} είναι εκείνες που «τέμνονται», δηλαδή «σταυρώνουν». Το κοινό σημείο τους λέγεται \lexiword{σημείο τομής}.
\end{lexi}

\begin{metafrasi}
Δύο ευθείες χωρίς κανένα κοινό σημείο. & παράλληλες: $\varepsilon_1\parallel\varepsilon_2$\\
Δύο ευθείες με ένα κοινό σημείο. & τεμνόμενες\\
Το κοινό τους σημείο. & σημείο τομής $\mathrm{Α}$\\
Τμήματα πάνω σε παράλληλες ευθείες. & $\mathrm{ΑΒ}\parallel\Gamma\Delta$\\
Δύο ευθείες που τέμνονται σε ορθές γωνίες. & κάθετες: $\varepsilon_1\perp\varepsilon_2$\\
\end{metafrasi}

\tmima{Η θεωρία}

\begin{orismos}{παράλληλες και τεμνόμενες ευθείες}
Δύο ευθείες του επιπέδου λέγονται \textbf{παράλληλες} όταν δεν έχουν κανένα κοινό σημείο, όσο κι αν προεκταθούν. Γράφουμε $\varepsilon_1\parallel\varepsilon_2$.\\
Δύο ευθείες του επιπέδου λέγονται \textbf{τεμνόμενες} όταν έχουν ένα μόνο κοινό σημείο. Το σημείο αυτό λέγεται \textbf{σημείο τομής} τους.
\end{orismos}

\begin{center}
\begin{tikzpicture}[scale=0.9]
\draw[thick,{Stealth}-{Stealth}] (0,0)--(4,0); \node[right] at (4,0) {$\varepsilon_2$}; \draw[thick,{Stealth}-{Stealth}] (0,1.3)--(4,1.3); \node[right] at (4,1.3) {$\varepsilon_1$};
\node at (2,-0.6) {$\varepsilon_1\parallel\varepsilon_2$};
\begin{scope}[xshift=7cm]
\draw[thick,{Stealth}-{Stealth}] (0,0)--(4,1.6); \draw[thick,{Stealth}-{Stealth}] (0,1.5)--(4,-0.1); \fill (2,0.77) circle (2pt) node[below] {$\mathrm{Α}$};
\node at (2,-0.6) {τεμνόμενες};
\end{scope}
\end{tikzpicture}
\end{center}

Δύο διαφορετικές ευθείες του ίδιου επιπέδου έχουν δύο μόνο δυνατότητες: ή είναι \emph{τεμνόμενες} ή είναι \emph{παράλληλες}. (Αν είχαν δύο κοινά σημεία, θα ταυτίζονταν, γιατί από δύο σημεία περνά μία μόνο ευθεία.) Δύο τμήματα λέγονται παράλληλα όταν βρίσκονται σε παράλληλες ευθείες.

Οι κάθετες ευθείες είναι μια ειδική περίπτωση τεμνόμενων ευθειών: τέμνονται και σχηματίζουν ορθές γωνίες.

\begin{kanonas}{δύο κάθετες στην ίδια ευθεία}
Δύο ευθείες που είναι κάθετες στην ίδια ευθεία είναι μεταξύ τους \textbf{παράλληλες}.
\end{kanonas}

\begin{center}
\begin{tikzpicture}[scale=0.9]
\draw[thick] (0,0.2)--(5,0.2) node[right] {$\varepsilon$};
\draw[thick] (1,-0.8)--(1,3) node[above] {$\varepsilon_1$}; \draw[thick] (3.6,-0.8)--(3.6,3) node[above] {$\varepsilon_2$};
\draw (1,0.5)--(1.3,0.5)--(1.3,0.2); \draw (3.6,0.5)--(3.9,0.5)--(3.9,0.2);
\node at (2.3,-1.2) {$\varepsilon_1\perp\varepsilon$, $\varepsilon_2\perp\varepsilon$, άρα $\varepsilon_1\parallel\varepsilon_2$};
\end{tikzpicture}
\end{center}

Έτσι φέρνουμε παράλληλη προς μια ευθεία $\zeta$ από ένα σημείο $\mathrm{Α}$ εκτός αυτής: φέρνουμε από το $\mathrm{Α}$ την κάθετη $\mathrm{ΑΒ}$ προς τη $\zeta$ και μετά, στο $\mathrm{Α}$, την κάθετη προς την $\mathrm{ΑΒ}$. Στην πράξη το κάνουμε με γνώμονα και χάρακα: ακουμπάμε τον χάρακα στη δοσμένη ευθεία, τον γνώμονα πάνω στον χάρακα και τον σύρουμε μέχρι να περάσει από το σημείο.

\begin{idiotita}{αίτημα του Ευκλείδη}
Από ένα σημείο εκτός ευθείας διέρχεται \textbf{μία και μόνο μία} ευθεία παράλληλη προς αυτή.
\end{idiotita}

Το αίτημα αυτό το δεχόμαστε χωρίς απόδειξη. Ανήκει στα «Στοιχεία» του Ευκλείδη, που γράφτηκαν στην Αλεξάνδρεια γύρω στο $300$ π.Χ. Από αυτό προκύπτει ότι αν δύο ευθείες είναι παράλληλες προς μια τρίτη, τότε είναι παράλληλες μεταξύ τους: αλλιώς θα είχαν σημείο τομής και από αυτό θα περνούσαν δύο παράλληλες προς την τρίτη.

\tmimasmall{Απόσταση δύο παράλληλων ευθειών}

Αν από ένα σημείο της μίας από δύο παράλληλες ευθείες φέρουμε το κάθετο τμήμα προς την άλλη, το μήκος του είναι η \textbf{απόσταση των παραλλήλων}. Όπου κι αν πάρουμε το σημείο, το κάθετο τμήμα έχει το ίδιο μήκος: οι παράλληλες είναι «ισαπέχουσες».

\begin{center}
\begin{tikzpicture}[scale=0.9]
\draw[thick] (0,0)--(5,0) node[right] {$\varepsilon_2$}; \draw[thick] (0,1.5)--(5,1.5) node[right] {$\varepsilon_1$};
\draw[gblue,thick] (1,0)--(1,1.5) node[midway,left] {\small $\delta$}; \draw[gblue,thick] (3.2,0)--(3.2,1.5) node[midway,right] {\small $\delta$};
\draw (1,0.25)--(1.25,0.25)--(1.25,0); \draw (3.2,0.25)--(3.45,0.25)--(3.45,0);
\end{tikzpicture}
\end{center}

\begin{prosoxi}
\begin{itemize}[leftmargin=1.2em,itemsep=0.3ex,topsep=2pt]
\item Για να είναι παράλληλες δύο ευθείες δεν αρκεί να μην συναντιούνται στο σχέδιο: πρέπει να μην έχουν κοινό σημείο όσο κι αν προεκταθούν. Δύο τμήματα που δεν ακουμπούν μπορεί να βρίσκονται σε τεμνόμενες ευθείες.
\item Δύο παράλληλα τμήματα δεν χρειάζεται να έχουν το ίδιο μήκος.
\item Οι κάθετες είναι τεμνόμενες, αλλά οι τεμνόμενες δεν είναι πάντα κάθετες.
\end{itemize}
\end{prosoxi}

\tmima{Λυμένα παραδείγματα}

\paradeigma{Παράλληλες ή τεμνόμενες;}
\tixero Οι ευθείες: οι ράγες του τρένου, δύο μονοπάτια που διασταυρώνονται, οι γραμμές του τετραδίου.\par
\tithelo Να πω ποιες ευθείες είναι παράλληλες και ποιες τεμνόμενες.\par
\possindeo Παράλληλες είναι όσες δεν έχουν κοινό σημείο, τεμνόμενες όσες έχουν ένα.\par
\apantisi Ράγες και γραμμές του τετραδίου: παράλληλες. Μονοπάτια: τεμνόμενες.\par
\elegxos Οι ράγες δεν θα συναντηθούν ούτε αν προεκταθούν.

\paradeigma{Φέρνω παράλληλη}
\tixero Μια ευθεία $\zeta$ και ένα σημείο $\mathrm{Α}$ εκτός της.\par
\tithelo Η ευθεία $\varepsilon$ από το $\mathrm{Α}$ παράλληλη προς τη $\zeta$.\par
\possindeo Φέρνω από το $\mathrm{Α}$ το κάθετο τμήμα $\mathrm{ΑΒ}$ προς τη $\zeta$. Στο $\mathrm{Α}$ φέρνω την κάθετη στην $\mathrm{ΑΒ}$, δηλαδή την $\varepsilon$. Οι $\zeta$ και $\varepsilon$ είναι κάθετες στην ίδια ευθεία.\par
\apantisi Η $\varepsilon$ είναι παράλληλη προς τη $\zeta$ και είναι η μοναδική που περνά από το $\mathrm{Α}$.\par
\elegxos Η απόσταση των $\varepsilon$ και $\zeta$ είναι $\mathrm{ΑΒ}$, και μετρώντας σε δεύτερο σημείο βρίσκουμε το ίδιο μήκος.

\paradeigma{Δύο παράλληλες σε απόσταση}
\tixero Μια ευθεία $\varepsilon$.\par
\tithelo Ευθείες παράλληλες προς την $\varepsilon$ που απέχουν $3$ cm από αυτήν.\par
\possindeo Παίρνω ένα σημείο $\mathrm{Α}$ της $\varepsilon$ και φέρνω το κάθετο τμήμα $\mathrm{ΑΒ}=3$ cm προς τη μία μεριά και $\mathrm{ΑΒ}'=3$ cm προς την άλλη. Από τα $\mathrm{Β}$, $\mathrm{Β}'$ φέρνω παράλληλες προς την $\varepsilon$.\par
\apantisi Υπάρχουν δύο τέτοιες ευθείες, μία σε κάθε μεριά της $\varepsilon$.\par
\elegxos Και οι δύο έχουν απόσταση $3$ cm από την $\varepsilon$.

\paradeigma{Ευθεία ή τμήμα;}
\tixero Δύο τμήματα $\mathrm{ΑΒ}$ και $\Gamma\Delta$ που δεν ακουμπούν στο χαρτί.\par
\tithelo Είναι παράλληλα;\par
\possindeo Προεκτείνω και τα δύο: αν οι ευθείες τους τέμνονται κάπου εκτός χαρτιού, τα τμήματα δεν είναι παράλληλα.\par
\apantisi Μπορεί να μην είναι: αποφασίζω από τις ευθείες, όχι από τα τμήματα.\par
\elegxos Αν οι προεκτάσεις δεν συναντιούνται, είναι παράλληλα.

\tmima{Ασκήσεις}

\zoni{Το κατάλαβα}{παράλληλες και τεμνόμενες}

\ex Σωστό ή λάθος; α) Δύο ευθείες του επιπέδου που δεν έχουν κανένα κοινό σημείο είναι παράλληλες. \quad β) Δύο τμήματα που δεν έχουν κανένα κοινό σημείο είναι πάντα παράλληλα. \quad γ) Αν δύο ευθείες είναι κάθετες μεταξύ τους, τότε είναι τεμνόμενες. \quad δ) Από ένα σημείο εκτός ευθείας περνά μία μόνο παράλληλη προς αυτήν. \quad ε) Δύο ευθείες κάθετες στην ίδια ευθεία είναι παράλληλες.

\ex Πες αν οι ευθείες είναι παράλληλες ή τεμνόμενες: α) οι δύο γραμμές διαχωρισμού μιας λωρίδας του δρόμου, β) οι ράγες του τρένου, γ) δύο δρόμοι που διασταυρώνονται, δ) οι απέναντι πλευρές ενός θρανίου, ε) τα δύο χέρια του ρολογιού στις $3{:}00$.

\ex Σχεδίασε μια ευθεία $\varepsilon$ και δύο σημεία $\mathrm{Α}$, $\mathrm{Β}$ εκτός της. Από τα $\mathrm{Α}$ και $\mathrm{Β}$ φέρε παράλληλες προς την $\varepsilon$. Τι παρατηρείς; Πότε ταυτίζονται οι δύο ευθείες;

\ex Σχεδίασε μια ευθεία $\varepsilon$ και ένα σημείο $\mathrm{Α}$ εκτός της. Φέρε από το $\mathrm{Α}$ την κάθετη προς την $\varepsilon$ και μετά, στο $\mathrm{Α}$, την κάθετη προς αυτήν. Τι σχέση έχει η δεύτερη κάθετη με την $\varepsilon$;

\ex Σχεδίασε μια γωνία $x\hat{\mathrm{Ο}}y=50^\circ$. Από ένα σημείο $\mathrm{Α}$ της $\mathrm{Ο}x$ φέρε παράλληλη προς την $\mathrm{Ο}y$ και από ένα σημείο $\mathrm{Β}$ της $\mathrm{Ο}y$ παράλληλη προς την $\mathrm{Ο}x$. Αν $\mathrm{Κ}$ είναι το σημείο τομής τους, μέτρησε τη γωνία $\mathrm{ΑΚΒ}$.

\ex Σχεδίασε μια ευθεία $\varepsilon$ και δύο ευθείες που απέχουν $4$ cm από αυτήν και είναι παράλληλες προς αυτήν. Πόσες τέτοιες ευθείες υπάρχουν;

\ex Σχεδίασε δύο τμήματα που δεν τέμνονται πάνω στο χαρτί, αλλά οι ευθείες τους τέμνονται εκτός χαρτιού. Είναι παράλληλα τα τμήματα;

\zoni{Το χρησιμοποιώ}{κήποι, ράγες, πλακάκια}

\ex \textbf{Οι σειρές των καρότων.} Ο παππούς Γιάννης φυτεύει πέντε παράλληλες σειρές καρότα και κάθε δύο διαδοχικές απέχουν $40$ cm. Πόση είναι η απόσταση ανάμεσα στην πρώτη και στην τελευταία σειρά;

\ex \textbf{Οι ράγες.} Οι ράγες του τρένου απέχουν $1{,}5$ μ. σε ένα σημείο. Πόσο απέχουν $2$ χιλιόμετρα πιο κάτω; Εξήγησε.

\ex \textbf{Οι δρόμοι.} Οι οδοί $\varepsilon_1$ και $\varepsilon_2$ είναι κάθετες στην οδό $\varepsilon_3$. Πώς είναι οι $\varepsilon_1$ και $\varepsilon_2$ μεταξύ τους;

\ex \textbf{Τα πλακάκια.} Η θεία Μαρία στρώνει πλακάκια στην ταράτσα. Οι αρμοί προς τη μία κατεύθυνση είναι κάθετοι στους αρμούς της άλλης κατεύθυνσης. Πώς είναι οι αρμοί της ίδιας κατεύθυνσης μεταξύ τους; Εξήγησε.

\ex \textbf{Το τετράδιο.} Πώς θα σχεδιάσεις με χάρακα και γνώμονα μια ευθεία παράλληλη προς μια δοσμένη, που να περνά από ένα δοσμένο σημείο; Γράψε τα βήματα.

\zoni{Το σκέφτομαι}{θέσεις και αιτήματα}

\ex Γιατί δύο διαφορετικές ευθείες του επιπέδου δεν μπορούν να έχουν δύο κοινά σημεία;

\ex Αν $\varepsilon_1\parallel\varepsilon_2$ και $\varepsilon_2\parallel\varepsilon_3$, τι σχέση έχουν οι $\varepsilon_1$ και $\varepsilon_3$; Εξήγησε με το αίτημα του Ευκλείδη.

\ex Από ένα σημείο εκτός ευθείας $\varepsilon$ πόσες ευθείες περνούν που δεν τέμνουν την $\varepsilon$; Και πόσες που την τέμνουν;

\ex Είναι κάθετες οι τεμνόμενες ευθείες; Είναι τεμνόμενες οι κάθετες ευθείες; Δώσε παραδείγματα.

\zoni{Για το σπίτι}{ψάξε παράλληλες}

\ex Βρες στο σπίτι σου πέντε ζεύγη παράλληλων ευθειών (ράφια, πλευρές παραθύρου κ.λπ.) και δύο ζεύγη κάθετων.

\ex Δίπλωσε ένα φύλλο στη μέση. Δίπλωσέ το ξανά ώστε η πτυχή να πέσει πάνω στον εαυτό της. Έτσι έχεις φτιάξει κάθετες πτυχές. Φτιάξε με τον ίδιο τρόπο δεύτερη κάθετη στην ίδια πτυχή. Πώς είναι οι δύο μεταξύ τους;

\paragrafos{Γωνίες σε παράλληλες ευθείες με τέμνουσα}{Β.3.8, «Γωνίες σε παράλληλες ευθείες με τέμνουσα»}

\begin{anagke}
Η θεία Μαρία φυτεύει ντομάτες σε δύο ευθείες σειρές, που είναι παράλληλες, και ένα λοξό μονοπάτι τις διασχίζει και τις δύο. Η Ρήνη, που βοηθά στον κήπο, μετρά με τον μοιρογνωμόνιο τη γωνία ανάμεσα στο μονοπάτι και στην πρώτη σειρά: $62^\circ$. «Να μετρήσω και την άλλη σειρά;», ρωτά. «Δεν χρειάζεται», λέει η κυρία Αλεξάνδρα, που περνούσε τυχαία, «αν ξέρεις τι σημαίνει να είναι παράλληλες, τη βρίσκεις με το μυαλό». Ο Ηλίας μετρά, για να βεβαιωθεί: $62^\circ$. Και στις υπόλοιπες θέσεις; Βγαίνουν συνολικά οκτώ γωνίες και, όπως φαίνεται, μόνο δύο διαφορετικά μέτρα.

Για να μιλήσουμε με ακρίβεια για το ποια γωνία είναι ίση με ποια, χρειαζόμαστε ονόματα: «εντός», «εκτός», «εναλλάξ», «επί τα αυτά».

Σε αυτή την παράγραφο μαθαίνουμε να ονομάζουμε τις οκτώ γωνίες που σχηματίζονται όταν δύο παράλληλες ευθείες τέμνονται από μια τρίτη, και να βρίσκουμε όλες τις γωνίες από μία μόνο.
\end{anagke}

\begin{lexi}{τέμνουσα, εντός, εκτός, εναλλάξ, επί τα αυτά}{θέσεις γωνιών}
Η \lexiword{τέμνουσα} είναι η ευθεία που «κόβει» τις άλλες δύο. Το \lexiword{εντός} σημαίνει «μέσα» και το \lexiword{εκτός} «έξω». Το \lexiword{εναλλάξ} είναι οι θέσεις που «αλλάζουν» πλευρά, δηλαδή αντίθετες. Το \lexiword{επί τα αυτά} σημαίνει «προς τα ίδια μέρη», δηλαδή προς την ίδια πλευρά.
\end{lexi}

\begin{metafrasi}
Η ευθεία που κόβει δύο άλλες ευθείες. & τέμνουσα\\
Μέσα στη λωρίδα ανάμεσα στις παράλληλες. & εντός\\
Έξω από τη λωρίδα των παραλλήλων. & εκτός\\
Στην ίδια πλευρά της τέμνουσας. & επί τα αυτά\\
Σε αντίθετες πλευρές της τέμνουσας. & εναλλάξ\\
\end{metafrasi}

\tmima{Η θεωρία}

Δύο ευθείες $\varepsilon_1$ και $\varepsilon_2$ του ίδιου επιπέδου είναι \textbf{παράλληλες} ($\varepsilon_1\parallel\varepsilon_2$) όταν δεν έχουν κανένα κοινό σημείο, όσο κι αν προεκταθούν. Αν μια τρίτη ευθεία $\varepsilon_3$ τις τέμνει και τις δύο, λέγεται \textbf{τέμνουσά} τους. Τότε σχηματίζονται οκτώ γωνίες, τις οποίες ονομάζουμε $\alpha,\beta,\gamma,\delta$ γύρω από το πρώτο σημείο τομής και $\varepsilon,\zeta,\eta,\theta$ γύρω από το δεύτερο.

\begin{center}
\begin{tikzpicture}[scale=1.05]
\draw[thick] (-3,1.2)--(3,1.2) node[right] {$\varepsilon_1$};
\draw[thick] (-3,-1.2)--(3,-1.2) node[right] {$\varepsilon_2$};
\draw[thick,gblue] (1.4,2.8)--(-1.4,-2.8) node[below left] {$\varepsilon_3$};
\fill (0.6,1.2) circle (1.6pt); \fill (-0.6,-1.2) circle (1.6pt);
\node at (1.05,1.47) {\small $\alpha$}; \node at (0.33,1.65) {\small $\beta$};
\node at (0.15,0.75) {\small $\gamma$}; \node at (1.0,0.75) {\small $\delta$};
\node at (-0.15,-0.93) {\small $\varepsilon$}; \node at (-0.87,-0.75) {\small $\zeta$};
\node at (-1.05,-1.47) {\small $\eta$}; \node at (-0.15,-1.65) {\small $\theta$};
\end{tikzpicture}
\end{center}

\textbf{Ως προς τις παράλληλες:} οι γωνίες $\gamma,\delta,\varepsilon,\zeta$ βρίσκονται ανάμεσα στις $\varepsilon_1,\varepsilon_2$ και λέγονται \textbf{εντός}. Οι $\alpha,\beta,\eta,\theta$ βρίσκονται έξω από τη λωρίδα και λέγονται \textbf{εκτός}.

\textbf{Ως προς την τέμνουσα:} δύο γωνίες που βρίσκονται προς το ίδιο μέρος της $\varepsilon_3$ λέγονται \textbf{επί τα αυτά}. Αν βρίσκονται σε διαφορετικά μέρη, λέγονται \textbf{εναλλάξ}.

Συνδυάζοντας τα δύο, παίρνουμε τα ζεύγη του πίνακα.

\begin{center}
\renewcommand{\arraystretch}{1.25}
\begin{tabular}{@{}p{0.37\linewidth}p{0.29\linewidth}p{0.28\linewidth}@{}}
\toprule
Ζεύγος & Ονομασία & Σχέση\\
\midrule
$\gamma,\varepsilon$ και $\delta,\zeta$ & εντός εναλλάξ & ίσες\\
$\alpha,\eta$ και $\beta,\theta$ & εκτός εναλλάξ & ίσες\\
$\alpha,\varepsilon$ και $\beta,\zeta$ και $\gamma,\eta$ και $\delta,\theta$ & εντός--εκτός και επί τα αυτά & ίσες\\
$\delta,\varepsilon$ και $\gamma,\zeta$ & εντός και επί τα αυτά & παραπληρωματικές\\
$\alpha,\theta$ και $\beta,\eta$ & εκτός και επί τα αυτά & παραπληρωματικές\\
\bottomrule
\end{tabular}
\end{center}

\begin{kanonas}{παράλληλες ευθείες και τέμνουσα}
Αν $\varepsilon_1\parallel\varepsilon_2$, τότε:
\begin{itemize}[leftmargin=1.2em,itemsep=0.3ex,topsep=2pt]
\item οι \textbf{εντός εναλλάξ} γωνίες είναι ίσες,
\item οι \textbf{εκτός εναλλάξ} γωνίες είναι ίσες,
\item οι \textbf{εντός και εκτός επί τα αυτά} γωνίες είναι ίσες,
\item οι \textbf{εντός και επί τα αυτά} γωνίες έχουν άθροισμα $180^\circ$ (είναι παραπληρωματικές), και το ίδιο οι \textbf{εκτός και επί τα αυτά}.
\end{itemize}
Με απλά λόγια: όλες οι \textbf{οξείες} γωνίες του σχήματος είναι ίσες μεταξύ τους, όλες οι \textbf{αμβλείες} είναι ίσες μεταξύ τους, και μια οξεία με μια αμβλεία έχουν άθροισμα $180^\circ$.
\end{kanonas}

Γιατί συμβαίνει αυτό; Όταν δύο ευθείες είναι παράλληλες, η τέμνουσα «γέρνει» προς αυτές με την ίδια κλίση και στις δύο, άρα οι αντίστοιχες γωνίες είναι ίδιες. Οι υπόλοιπες προκύπτουν από τις κατακορυφήν γωνίες (ίσες) και τις παραπληρωματικές (άθροισμα $180^\circ$).

Ισχύει και το \textbf{αντίστροφο}, που μας βοηθά να ελέγχουμε αν δύο ευθείες είναι παράλληλες: αν δύο ευθείες τεμνόμενες από τρίτη σχηματίζουν \emph{εντός εναλλάξ} γωνίες ίσες (ή εντός και επί τα αυτά γωνίες με άθροισμα $180^\circ$), τότε οι ευθείες είναι παράλληλες.

\begin{prosoxi}
\begin{itemize}[leftmargin=1.2em,itemsep=0.3ex,topsep=2pt]
\item Οι ισότητες ισχύουν \emph{μόνο όταν οι ευθείες είναι παράλληλες}. Αν δεν είναι, οι εντός εναλλάξ γωνίες δεν είναι ίσες.
\item Οι εντός και επί τα αυτά γωνίες είναι \emph{παραπληρωματικές}, όχι ίσες (εκτός αν είναι και οι δύο $90^\circ$).
\item Μην μπερδεύεις το «εναλλάξ» με το «επί τα αυτά»: το πρώτο σημαίνει διαφορετικές πλευρές της τέμνουσας, το δεύτερο την ίδια.
\end{itemize}
\end{prosoxi}

\tmima{Λυμένα παραδείγματα}

\paradeigma{Όλες οι γωνίες από μία}
\tixero $\varepsilon_1\parallel\varepsilon_2$ και στο σχήμα της θεωρίας $\gamma=52^\circ$.\par
\tithelo Οι υπόλοιπες επτά γωνίες.\par
\possindeo Η $\alpha$ είναι κατακορυφήν με τη $\gamma$, άρα $\alpha=52^\circ$. Η $\beta$ είναι παραπληρωματική της $\gamma$ (σχηματίζουν ευθεία), άρα $\beta=180^\circ-52^\circ=128^\circ$, και η $\delta=\beta=128^\circ$ (κατακορυφήν). Η $\varepsilon$ είναι εντός εναλλάξ με τη $\gamma$, άρα $\varepsilon=52^\circ$. Η $\zeta$ είναι εντός και επί τα αυτά με τη $\gamma$, άρα $\zeta=180^\circ-52^\circ=128^\circ$. Τέλος $\eta=\varepsilon=52^\circ$ και $\theta=\zeta=128^\circ$ (κατακορυφήν).\par
\apantisi $\alpha=\gamma=\varepsilon=\eta=52^\circ$ και $\beta=\delta=\zeta=\theta=128^\circ$.\par
\elegxos Τέσσερις οξείες ίσες, τέσσερις αμβλείες ίσες, και $52^\circ+128^\circ=180^\circ$.

\paradeigma{Εντός εναλλάξ με άγνωστο}
\tixero $\varepsilon_1\parallel\varepsilon_2$ και δύο εντός εναλλάξ γωνίες είναι $x+20^\circ$ και $2x-40^\circ$.\par
\tithelo Το $x$ και οι δύο γωνίες.\par
\possindeo Οι εντός εναλλάξ γωνίες είναι ίσες, άρα $x+20^\circ=2x-40^\circ$. Προσθέτω $40^\circ$ και στα δύο μέλη: $x+60^\circ=2x$, άρα $x=60^\circ$.\par
\apantisi $x=60^\circ$ και κάθε γωνία είναι $80^\circ$.\par
\elegxos $60^\circ+20^\circ=80^\circ$ και $2\cdot60^\circ-40^\circ=80^\circ$.

\paradeigma{Εντός και επί τα αυτά με άγνωστο}
\tixero $\varepsilon_1\parallel\varepsilon_2$ και δύο εντός και επί τα αυτά γωνίες είναι $4x$ και $x+20^\circ$.\par
\tithelo Οι δύο γωνίες.\par
\possindeo Είναι παραπληρωματικές: $4x+x+20^\circ=180^\circ$, άρα $5x=160^\circ$ και $x=32^\circ$.\par
\apantisi $4x=128^\circ$ και $x+20^\circ=52^\circ$.\par
\elegxos $128^\circ+52^\circ=180^\circ$.

\paradeigma{Παράλληλες ή όχι;}
\tixero Δύο ευθείες τέμνονται από τρίτη και οι εντός εναλλάξ γωνίες είναι $72^\circ$ και $75^\circ$.\par
\tithelo Είναι παράλληλες;\par
\possindeo Αν ήταν παράλληλες, οι εντός εναλλάξ γωνίες θα ήταν ίσες.\par
\apantisi Όχι, γιατί $72^\circ\neq75^\circ$.\par
\elegxos Αν και οι δύο ήταν $72^\circ$, θα ήταν παράλληλες.

\tmima{Ασκήσεις}

\zoni{Το κατάλαβα}{ονόματα και ισότητες}

\ex Στο σχήμα της θεωρίας, πες πώς ονομάζεται κάθε ζεύγος γωνιών: α) $\alpha,\varepsilon$ β) $\gamma,\varepsilon$ γ) $\beta,\theta$ δ) $\delta,\varepsilon$ ε) $\gamma,\zeta$ στ) $\beta,\zeta$.

\ex Οι $\varepsilon_1,\varepsilon_2$ είναι παράλληλες και $\alpha=70^\circ$ στο σχήμα της θεωρίας. Βρες και τις επτά άλλες γωνίες.

\ex Αν $\varepsilon_1\parallel\varepsilon_2$ και $\delta=118^\circ$, βρες τις $\beta,\zeta,\theta,\varepsilon,\alpha$ και εξήγησε με ποια ιδιότητα βρίσκεις την καθεμία.

\ex Δύο εντός εναλλάξ γωνίες παραλλήλων είναι $3x-10^\circ$ και $x+50^\circ$. Βρες το $x$ και τις γωνίες.

\ex Δύο εντός και επί τα αυτά γωνίες παραλλήλων είναι $x$ και $4x$. Βρες τες.

\ex Σωστό ή λάθος; α) Οι εντός εναλλάξ γωνίες είναι πάντα ίσες. β) Αν οι ευθείες είναι παράλληλες, οι εκτός εναλλάξ γωνίες είναι ίσες. γ) Οι εντός και επί τα αυτά γωνίες παραλλήλων είναι ίσες. δ) Δύο κατακορυφήν γωνίες είναι πάντα ίσες, είτε οι ευθείες είναι παράλληλες είτε όχι. ε) Αν οι εντός εναλλάξ γωνίες είναι ίσες, οι ευθείες είναι παράλληλες.

\ex Σχεδίασε δύο παράλληλες ευθείες που απέχουν $3$ cm και μια τέμνουσα που σχηματίζει $55^\circ$ με τη μία. Υπολόγισε όλες τις γωνίες και μέτρησε τες με τον μοιρογνωμόνιο για να ελέγξεις.

\zoni{Το χρησιμοποιώ}{κήπος, δρόμος, σχολείο}

\ex \textbf{Το μονοπάτι.} Στον κήπο της θείας Μαρίας οι δύο σειρές ντομάτας είναι παράλληλες. Μία από τις γωνίες ανάμεσα στο μονοπάτι και μια σειρά είναι $62^\circ$. Πόσες από τις οκτώ γωνίες είναι $62^\circ$ και πόσες είναι $118^\circ$;

\ex \textbf{Ο ποδηλατόδρομος.} Δύο παράλληλες λωρίδες ποδηλατόδρομου τέμνονται από μια διάβαση πεζών που σχηματίζει με την πρώτη λωρίδα μια γωνία $48^\circ$, προς το εσωτερικό του δρόμου. Πόση είναι η εντός εναλλάξ γωνία της στη δεύτερη λωρίδα; Και πόση η εντός και επί τα αυτά;

\ex \textbf{Το κιγκλίδωμα.} Σε ένα κιγκλίδωμα δύο οριζόντιες ράβδοι είναι παράλληλες και μια λοξή ράβδος τις διασχίζει. Μια γωνία της λοξής με την πάνω ράβδο, προς το εσωτερικό, είναι $125^\circ$. Πόση είναι η γωνία της με την κάτω ράβδο, επί τα αυτά μέρη και εντός;

\ex \textbf{Ο φράχτης του Βασίλη.} Ο Βασίλης λέει ότι τα δύο δοκάρια του φράχτη του σχολείου είναι παράλληλα. Μια λοξή σανίδα σχηματίζει με αυτά εντός εναλλάξ γωνίες $74^\circ$ και $76^\circ$. Έχει δίκιο;

\ex \textbf{Ο πίνακας της Φανής.} Η Φανή ελέγχει αν οι δύο απέναντι πλευρές ενός πίνακα είναι παράλληλες. Με μια λοξή ράβδο μετρά εντός και επί τα αυτά γωνίες $98^\circ$ και $82^\circ$. Είναι παράλληλες; Δικαιολόγησε.

\zoni{Το σκέφτομαι}{γιατί συμβαίνει}

\ex Εξήγησε γιατί, αν οι εντός εναλλάξ γωνίες δύο παραλλήλων είναι ίσες, τότε οι εντός και επί τα αυτά είναι παραπληρωματικές.

\ex Στο σχήμα, $\mathrm{ΒΕ}$ και $\mathrm{ΕΓ}$ είναι δύο τμήματα που ξεκινούν από το $\mathrm{Ε}$, με $\mathrm{Β}$ στην $\varepsilon_1$ και $\mathrm{Γ}$ στην $\varepsilon_2$, και $\varepsilon_1\parallel\varepsilon_2$. Η $\mathrm{ΕΒ}$ σχηματίζει με την $\varepsilon_1$ γωνία $40^\circ$ και η $\mathrm{ΕΓ}$ με την $\varepsilon_2$ γωνία $30^\circ$. Βρες τη γωνία $\widehat{\mathrm{ΒΕΓ}}$. (Βοήθεια: φέρε από το $\mathrm{Ε}$ παράλληλη προς τις δύο ευθείες.)

\begin{center}
\begin{tikzpicture}[scale=0.9]
\draw[thick] (-0.8,1.4)--(5.2,1.4) node[right] {$\varepsilon_1$};
\draw[thick] (-0.8,-1.4)--(5.2,-1.4) node[right] {$\varepsilon_2$};
\draw[thick] (3.4,1.4)--(1.73,0)--(4.15,-1.4);
\fill (3.4,1.4) circle (1.6pt) node[above] {$\mathrm{Β}$};
\fill (1.73,0) circle (1.6pt) node[left] {$\mathrm{Ε}$};
\fill (4.15,-1.4) circle (1.6pt) node[below] {$\mathrm{Γ}$};
\node at (2.5,1.15) {\small $40^\circ$};
\node at (3.2,-1.15) {\small $30^\circ$};
\end{tikzpicture}
\end{center}

\ex Στο παραλληλόγραμμο που σχηματίζουν οι παράλληλες $\varepsilon_1\parallel\varepsilon_2$ και $\varepsilon_3\parallel\varepsilon_4$, μια γωνία είναι $70^\circ$. Βρες τις τέσσερις γωνίες του τετραπλεύρου που περικλείεται.

\ex Μια οξεία και μια αμβλεία γωνία που σχηματίζονται από παράλληλες και τέμνουσα είναι εντός και επί τα αυτά, και η αμβλεία είναι τριπλάσια της οξείας. Βρες και τις δύο.

\zoni{Για το σπίτι}{ρίξε μια ματιά}

\ex Βρες στο σπίτι ή στον δρόμο δύο παράλληλες ευθείες (σιδηροτροχιά, κάγκελα, ράφια) που τέμνονται από μια τρίτη. Μέτρησε με τον μοιρογνωμόνιο δύο εντός εναλλάξ γωνίες. Είναι ίσες;

\ex Σχεδίασε με χάρακα και γνώμονα δύο παράλληλες ευθείες και μια τέμνουσα. Βάψε με το ίδιο χρώμα τις γωνίες που είναι ίσες και με άλλο χρώμα τις παραπληρωματικές τους.

\paragrafos{Άθροισμα γωνιών τριγώνου}{Β.3.9, «Άθροισμα γωνιών τριγώνου»}

\begin{anagke}
Η κυρία Αλεξάνδρα ζωγραφίζει στον πίνακα ένα τρίγωνο και ζητά από τον Στέλιο να μετρήσει τις γωνίες του. Εκείνος μετρά δύο γωνίες, $62^\circ$ και $48^\circ$, αλλά το μοιρογνωμόνιο γλιστράει και η τρίτη μέτρηση χάνεται. «Δεν πειράζει», λέει η Φανή, «δεν χρειάζεται να τη μετρήσουμε ξανά». Ο Στέλιος δεν πείθεται: «Και πώς θα τη βρούμε χωρίς να τη μετρήσουμε;». Η Ζωή ρωτά: «Μήπως οι τρεις γωνίες ενός τριγώνου έχουν πάντα το ίδιο άθροισμα, όποιο τρίγωνο κι αν σχεδιάσεις;». Η κυρία Αλεξάνδρα χαμογελά: «Κόψτε ένα τρίγωνο από χαρτί, σκίστε τις τρεις γωνίες και βάλτε τες δίπλα δίπλα με τις κορυφές στο ίδιο σημείο. Τι σχηματίζεται;».

Σε αυτή την παράγραφο μαθαίνουμε ότι το άθροισμα των γωνιών κάθε τριγώνου είναι πάντα το ίδιο, και το χρησιμοποιούμε για να βρίσκουμε γωνίες που δεν μπορούμε ή δεν θέλουμε να μετρήσουμε.
\end{anagke}

\begin{lexi}{άθροισμα, προσκείμενες γωνίες, συμπληρωματικές}{άθροισμα γωνιών}
Το \lexiword{άθροισμα} είναι το «ύψος» που φτάνει ένας σωρός όταν βάζεις τα πράγματα ένα πάνω στο άλλο, δηλαδή το αποτέλεσμα της πρόσθεσης. Οι \lexiword{προσκείμενες γωνίες} μιας πλευράς είναι όσες «κάθονται» πάνω της (έχουν την πλευρά αυτή κοινή). Οι \lexiword{συμπληρωματικές} γωνίες «συμπληρώνουν» η μία την άλλη μέχρι την ορθή γωνία.
\end{lexi}

\begin{metafrasi}
Οι τρεις γωνίες του τριγώνου μαζί. & $\hat{\mathrm{Α}}+\hat{\mathrm{Β}}+\hat{\mathrm{Γ}}$\\
Μια ευθεία γωνία έχει μέτρο. & $180^\circ$\\
Οι γωνίες που κάθονται στη βάση του ισοσκελούς. & προσκείμενες στη βάση, $\hat{\mathrm{Β}}=\hat{\mathrm{Γ}}$\\
Δύο γωνίες που μαζί κάνουν ορθή. & συμπληρωματικές ($90^\circ$)\\
\end{metafrasi}

\tmima{Η θεωρία}

Το πείραμα με το χαρτί δείχνει ότι οι τρεις γωνίες ενός τριγώνου σχηματίζουν μαζί μια ευθεία γωνία. Μπορούμε να το αποδείξουμε με όσα ξέρουμε για τις παράλληλες ευθείες.

\begin{kanonas}{άθροισμα γωνιών τριγώνου}
Το άθροισμα των γωνιών κάθε τριγώνου είναι $180^\circ$:
\[\hat{\mathrm{Α}}+\hat{\mathrm{Β}}+\hat{\mathrm{Γ}}=180^\circ.\]
\end{kanonas}

\begin{center}
\begin{tikzpicture}[scale=0.95]
\draw[thick] (0,0)--(5,0)--(1.6,2.4)--cycle;
\draw[gblue,thick] (-0.8,2.4)--(5.8,2.4) node[right] {$\varepsilon$};
\node[below left] at (0,0) {$\mathrm{Β}$}; \node[below right] at (5,0) {$\mathrm{Γ}$}; \node[above] at (1.6,2.4) {$\mathrm{Α}$};
\fill (1.6,2.4) circle (1.6pt);
\node at (0.8,0.2) {\small $\hat{\mathrm{Β}}$};
\node at (4.2,0.2) {\small $\hat{\mathrm{Γ}}$};
\node at (1.5,1.9) {\small $\hat{\mathrm{Α}}$};
\node at (0.85,2.15) {\small $\hat{\mathrm{Α}}_1$};
\node at (2.9,2.15) {\small $\hat{\mathrm{Α}}_2$};
\end{tikzpicture}
\end{center}

\textbf{Απόδειξη.} Από την κορυφή $\mathrm{Α}$ φέρνουμε την ευθεία $\varepsilon$ παράλληλη προς τη $\mathrm{ΒΓ}$. Τότε $\hat{\mathrm{Α}}_1=\hat{\mathrm{Β}}$ (εντός εναλλάξ, με τέμνουσα την $\mathrm{ΑΒ}$) και $\hat{\mathrm{Α}}_2=\hat{\mathrm{Γ}}$ (εντός εναλλάξ, με τέμνουσα την $\mathrm{ΑΓ}$). Οι γωνίες $\hat{\mathrm{Α}}_1$, $\hat{\mathrm{Α}}$, $\hat{\mathrm{Α}}_2$ σχηματίζουν μαζί ευθεία γωνία, άρα $\hat{\mathrm{Α}}_1+\hat{\mathrm{Α}}+\hat{\mathrm{Α}}_2=180^\circ$. Αντικαθιστώντας, παίρνουμε $\hat{\mathrm{Β}}+\hat{\mathrm{Α}}+\hat{\mathrm{Γ}}=180^\circ$.

Από τον κανόνα προκύπτουν χρήσιμα συμπεράσματα.

\begin{idiotita}{συνέπειες}
\begin{itemize}[leftmargin=1.2em,itemsep=0.3ex,topsep=2pt]
\item Σε \textbf{ισόπλευρο} τρίγωνο όλες οι γωνίες είναι ίσες, άρα κάθε γωνία είναι $180^\circ:3=60^\circ$.
\item Σε \textbf{ορθογώνιο} τρίγωνο οι δύο οξείες γωνίες είναι \textbf{συμπληρωματικές}: $\hat{\mathrm{Β}}+\hat{\mathrm{Γ}}=90^\circ$.
\item Σε \textbf{ορθογώνιο και ισοσκελές} τρίγωνο οι οξείες γωνίες είναι $45^\circ$ και $45^\circ$.
\item Σε \textbf{ισοσκελές} τρίγωνο οι γωνίες της βάσης είναι ίσες, άρα
\[\text{κορυφή}=180^\circ-2\cdot(\text{βάση}),\]
\[\text{βάση}=\bigl(180^\circ-\text{κορυφή}\bigr):2.\]
\item Ένα τρίγωνο έχει \textbf{το πολύ μία} ορθή ή αμβλεία γωνία, και \textbf{τουλάχιστον δύο} οξείες.
\end{itemize}
\end{idiotita}

\begin{prosoxi}
\begin{itemize}[leftmargin=1.2em,itemsep=0.3ex,topsep=2pt]
\item Το άθροισμα είναι πάντα $180^\circ$, όσο μεγάλο ή μικρό κι αν είναι το τρίγωνο. Το μέγεθος δεν αλλάζει το άθροισμα.
\item Αν σε ισοσκελές τρίγωνο σου δίνουν μία γωνία, πρέπει να σκεφτείς αν είναι η γωνία της κορυφής ή της βάσης. Συχνά υπάρχουν \emph{δύο} δυνατές περιπτώσεις.
\item Το ισοσκελές τρίγωνο έχει ίσες μόνο τις γωνίες της βάσης. Και οι τρεις γωνίες είναι ίσες μόνο στο ισόπλευρο.
\end{itemize}
\end{prosoxi}

\tmima{Λυμένα παραδείγματα}

\paradeigma{Η τρίτη γωνία του Στέλιου}
\tixero Σε τρίγωνο δύο γωνίες είναι $62^\circ$ και $48^\circ$.\par
\tithelo Η τρίτη γωνία και το είδος του τριγώνου.\par
\possindeo Το άθροισμα των τριών γωνιών είναι $180^\circ$. Οι δύο γωνίες μαζί κάνουν $62^\circ+48^\circ=110^\circ$, άρα η τρίτη είναι $180^\circ-110^\circ=70^\circ$.\par
\apantisi $70^\circ$. Όλες οι γωνίες είναι οξείες, άρα το τρίγωνο είναι οξυγώνιο.\par
\elegxos $62^\circ+48^\circ+70^\circ=180^\circ$.

\paradeigma{Ισοσκελές με γνωστή κορυφή}
\tixero Ισοσκελές τρίγωνο $\mathrm{ΑΒΓ}$ ($\mathrm{ΑΒ}=\mathrm{ΑΓ}$) με $\hat{\mathrm{Α}}=50^\circ$.\par
\tithelo Οι γωνίες $\hat{\mathrm{Β}}$ και $\hat{\mathrm{Γ}}$.\par
\possindeo Οι γωνίες της βάσης είναι ίσες, $\hat{\mathrm{Β}}=\hat{\mathrm{Γ}}=x$. Τότε $50^\circ+x+x=180^\circ$, άρα $2x=130^\circ$ και $x=65^\circ$.\par
\apantisi $\hat{\mathrm{Β}}=\hat{\mathrm{Γ}}=65^\circ$.\par
\elegxos $50^\circ+65^\circ+65^\circ=180^\circ$.

\paradeigma{Ισοσκελές με μία γνωστή γωνία}
\tixero Ένα ισοσκελές τρίγωνο έχει μια γωνία $80^\circ$.\par
\tithelo Οι άλλες δύο γωνίες.\par
\possindeo Η γωνία των $80^\circ$ μπορεί να είναι η γωνία της κορυφής ή μία γωνία της βάσης.\par
\textbf{1η περίπτωση: κορυφή.} Οι γωνίες της βάσης είναι ίσες: $(180^\circ-80^\circ):2=50^\circ$ η καθεμία.\par
\textbf{2η περίπτωση: βάση.} Η άλλη γωνία της βάσης είναι επίσης $80^\circ$, και η κορυφή είναι $180^\circ-2\cdot80^\circ=20^\circ$.\par
\apantisi $(80^\circ,50^\circ,50^\circ)$ ή $(80^\circ,80^\circ,20^\circ)$.\par
\elegxos Και στις δύο περιπτώσεις το άθροισμα είναι $180^\circ$.

\paradeigma{Γωνίες με άγνωστο}
\tixero Οι γωνίες ενός τριγώνου είναι $x$, $x+20^\circ$ και $2x-20^\circ$.\par
\tithelo Το $x$ και οι γωνίες του τριγώνου.\par
\possindeo Το άθροισμα είναι $180^\circ$: $x+(x+20^\circ)+(2x-20^\circ)=180^\circ$. Στο άθροισμα, τα $+20^\circ$ και $-20^\circ$ αλληλοαναιρούνται και μένει $4x=180^\circ$, άρα $x=45^\circ$.\par
\apantisi Οι γωνίες είναι $45^\circ$, $65^\circ$ και $70^\circ$.\par
\elegxos $45^\circ+65^\circ+70^\circ=180^\circ$. Το τρίγωνο είναι οξυγώνιο.

\paradeigma{Παράλληλες και τρίγωνα}
\tixero Οι $\mathrm{ΑΒ}$ και $\mathrm{ΓΔ}$ είναι παράλληλες και τα τμήματα $\mathrm{ΑΔ}$ και $\mathrm{ΒΓ}$ τέμνονται στο $\mathrm{Ο}$. Είναι $\widehat{\mathrm{ΟΑΒ}}=40^\circ$ και $\widehat{\mathrm{ΟΓΔ}}=35^\circ$.\par
\tithelo Η γωνία $\widehat{\mathrm{ΑΟΒ}}$.\par
\possindeo Η $\widehat{\mathrm{ΟΒΑ}}$ είναι εντός εναλλάξ με τη $\widehat{\mathrm{ΟΓΔ}}$ (παράλληλες $\mathrm{ΑΒ}$, $\mathrm{ΓΔ}$ και τέμνουσα $\mathrm{ΒΓ}$), άρα $\widehat{\mathrm{ΟΒΑ}}=35^\circ$. Στο τρίγωνο $\mathrm{ΟΑΒ}$ είναι $40^\circ+35^\circ+\widehat{\mathrm{ΑΟΒ}}=180^\circ$, άρα $\widehat{\mathrm{ΑΟΒ}}=105^\circ$.\par
\apantisi $\widehat{\mathrm{ΑΟΒ}}=105^\circ$.\par
\elegxos Η $\widehat{\mathrm{ΓΟΔ}}$ είναι κατακορυφήν με την $\widehat{\mathrm{ΑΟΒ}}$, άρα $105^\circ$, και στο τρίγωνο $\mathrm{ΟΓΔ}$ έχουμε $35^\circ+40^\circ+105^\circ=180^\circ$.

\begin{center}
\begin{tikzpicture}[scale=0.85]
\draw[thick] (-1.6,1.34)--(1.92,1.34);
\draw[thick] (-1.92,-1.34)--(1.6,-1.34);
\draw[thick,gblue] (-1.6,1.34)--(1.6,-1.34);
\draw[thick,gblue] (1.92,1.34)--(-1.92,-1.34);
\fill (-1.6,1.34) circle (1.6pt) node[above left] {$\mathrm{Α}$};
\fill (1.92,1.34) circle (1.6pt) node[above right] {$\mathrm{Β}$};
\fill (-1.92,-1.34) circle (1.6pt) node[below left] {$\mathrm{Γ}$};
\fill (1.6,-1.34) circle (1.6pt) node[below right] {$\mathrm{Δ}$};
\fill (0,0) circle (1.6pt) node[above] {\small $\mathrm{Ο}$};
\node at (-0.65,1.1) {\small $40^\circ$};
\node at (-1.1,-1.1) {\small $35^\circ$};
\end{tikzpicture}
\end{center}

\tmima{Ασκήσεις}

\zoni{Το κατάλαβα}{άθροισμα $180^\circ$}

\ex Σωστό ή λάθος; α) Υπάρχει τρίγωνο με δύο ορθές γωνίες. β) Υπάρχει τρίγωνο με γωνίες $50^\circ$, $60^\circ$ και $70^\circ$. γ) Οι οξείες γωνίες ενός ορθογωνίου τριγώνου είναι συμπληρωματικές. δ) Υπάρχει ισόπλευρο τρίγωνο με γωνία $70^\circ$. ε) Αν η γωνία της κορυφής ισοσκελούς τριγώνου είναι $30^\circ$, κάθε γωνία της βάσης είναι $75^\circ$. στ) Ένα τρίγωνο μπορεί να έχει δύο αμβλείες γωνίες.

\ex Βρες την τρίτη γωνία τριγώνου αν οι δύο άλλες είναι: α) $45^\circ$ και $65^\circ$ β) $90^\circ$ και $34^\circ$ γ) $108^\circ$ και $25^\circ$ δ) $59^\circ$ και $61^\circ$.

\ex Σε ισοσκελές τρίγωνο η γωνία της κορυφής είναι: α) $40^\circ$ β) $100^\circ$ γ) $90^\circ$ δ) $20^\circ$. Βρες τις γωνίες της βάσης.

\ex Σε ισοσκελές τρίγωνο η γωνία της βάσης είναι: α) $65^\circ$ β) $30^\circ$ γ) $45^\circ$ δ) $80^\circ$. Βρες τη γωνία της κορυφής.

\ex Οι γωνίες ενός τριγώνου είναι $x$, $x+10^\circ$ και $x+20^\circ$. Βρες το $x$ και τις τρεις γωνίες.

\ex Ποιες από τις παρακάτω τριάδες μπορεί να είναι οι γωνίες ενός τριγώνου; α) $70^\circ,60^\circ,50^\circ$ β) $90^\circ,50^\circ,50^\circ$ γ) $100^\circ,40^\circ,40^\circ$ δ) $120^\circ,30^\circ,40^\circ$ ε) $60^\circ,60^\circ,60^\circ$ στ) $35^\circ,45^\circ,90^\circ$.

\ex Πες αν το τρίγωνο με τις παρακάτω γωνίες είναι οξυγώνιο, ορθογώνιο ή αμβλυγώνιο: α) $50^\circ,60^\circ,70^\circ$ β) $90^\circ,35^\circ,55^\circ$ γ) $105^\circ,40^\circ,35^\circ$ δ) $80^\circ,50^\circ,50^\circ$.

\zoni{Το χρησιμοποιώ}{στέγες, σκάλες, χωράφια}

\ex \textbf{Η στέγη.} Το αέτωμα του εξοχικού της θείας Μαρίας έχει σχήμα ισοσκελούς τριγώνου με γωνία στην κορυφή $110^\circ$. Πόση γωνία σχηματίζει κάθε πλευρά της στέγης με το οριζόντιο δοκάρι της βάσης;

\ex \textbf{Το σημαιάκι.} Η Ρήνη κόβει ένα ισοσκελές τρίγωνο για σημαιάκι και μία γωνία του είναι $50^\circ$. Βρες όλες τις δυνατές τριάδες γωνιών που μπορεί να έχει.

\ex \textbf{Η σκάλα.} Ο Ηλίας ακουμπά μια σκάλα σε κατακόρυφο τοίχο. Η σκάλα σχηματίζει με το οριζόντιο δάπεδο γωνία $72^\circ$. Πόση γωνία σχηματίζει με τον τοίχο;

\ex \textbf{Το χωράφι.} Ένα τριγωνικό χωράφι του παππού Γιάννη έχει γωνίες $2x$, $3x$ και $4x$. Βρες τις γωνίες του και το είδος του τριγώνου ως προς τις γωνίες.

\ex \textbf{Το σχέδιο της Αγγελικής.} Η Αγγελική ζωγραφίζει ισοσκελές τρίγωνο όπου η γωνία της κορυφής είναι τριπλάσια από τη γωνία της βάσης. Βρες τις γωνίες του.

\zoni{Το σκέφτομαι}{γιατί και πώς}

\ex Εξήγησε γιατί κάθε τρίγωνο έχει τουλάχιστον δύο οξείες γωνίες.

\ex Το τρίγωνο $\mathrm{ΑΒΓ}$ είναι ισοσκελές με $\mathrm{ΑΒ}=\mathrm{ΑΓ}$ και $\hat{\mathrm{Α}}=44^\circ$. Η ευθεία $\mathrm{ΔΕ}$ είναι παράλληλη προς τη $\mathrm{ΒΓ}$, με το $\mathrm{Δ}$ στην $\mathrm{ΑΒ}$ και το $\mathrm{Ε}$ στην $\mathrm{ΑΓ}$. Βρες τις γωνίες $\widehat{\mathrm{ΑΔΕ}}$, $\widehat{\mathrm{ΑΕΔ}}$ και $\widehat{\mathrm{ΔΕΓ}}$.

\begin{center}
\begin{tikzpicture}[scale=0.62]
\draw[thick] (0,0)--(4,0)--(2,4.95)--cycle;
\draw[thick,gblue] (0.8,1.98)--(3.2,1.98);
\node[below left] at (0,0) {$\mathrm{Β}$}; \node[below right] at (4,0) {$\mathrm{Γ}$}; \node[above] at (2,4.95) {$\mathrm{Α}$};
\node[left] at (0.8,1.98) {$\mathrm{Δ}$}; \node[right] at (3.2,1.98) {$\mathrm{Ε}$};
\node at (2,4.2) {\small $44^\circ$};
\end{tikzpicture}
\end{center}

\ex Οι $\varepsilon_1$ και $\varepsilon_2$ είναι παράλληλες και μια τέμνουσα τις τέμνει στα $\mathrm{Α}$ και $\mathrm{Β}$. Οι εντός και επί τα αυτά γωνίες είναι $70^\circ$ στο $\mathrm{Α}$ και $110^\circ$ στο $\mathrm{Β}$. Οι διχοτόμοι τους τέμνονται στο $\mathrm{Κ}$. Βρες τη γωνία $\widehat{\mathrm{ΑΚΒ}}$. Θα έβγαινε το ίδιο με άλλες τιμές γωνιών;

\ex Στο τρίγωνο $\mathrm{ΑΒΓ}$ είναι $\hat{\mathrm{Β}}=30^\circ$ και η γωνία $\hat{\mathrm{Α}}$ είναι τετραπλάσια της $\hat{\mathrm{Γ}}$. Βρες τις γωνίες του τριγώνου και το είδος του ως προς τις γωνίες και τις πλευρές του.

\zoni{Για το σπίτι}{ρίξε μια ματιά}

\ex Κόψε από χαρτί ένα τρίγωνο της αρεσκείας σου, χρωμάτισε τις τρεις γωνίες του με διαφορετικά χρώματα και σκίσε τες. Βάλε τες δίπλα δίπλα με τις κορυφές στο ίδιο σημείο. Τι σχηματίζουν; Δοκίμασε και με ένα αμβλυγώνιο τρίγωνο.

\ex Σχεδίασε τρίγωνο $\mathrm{ΑΒΓ}$ με $\mathrm{ΒΓ}=6$ cm, $\hat{\mathrm{Β}}=50^\circ$ και $\hat{\mathrm{Γ}}=60^\circ$. Υπολόγισε τη γωνία $\hat{\mathrm{Α}}$ και μετά μέτρησέ τη. Συμφωνούν;

\paragrafos{Τετράπλευρα, παραλληλόγραμμα, τραπέζια}{Β.3.10, «Τετράπλευρα, παραλληλόγραμμα, τραπέζια»}

\begin{anagke}
Η θεία Μαρία στρώνει πλακάκια στην αυλή και έχει μπροστά της τρεις στοίβες. Τα πλακάκια της πρώτης έχουν τα απέναντι ζευγάρια πλευρών παράλληλα και εφαρμόζουν τέλεια το ένα δίπλα στο άλλο. Της δεύτερης έχουν μόνο ένα ζευγάρι παράλληλων πλευρών και μοιάζουν με μικρά τραπεζάκια. Της τρίτης δεν έχουν καθόλου παράλληλες πλευρές και δύσκολα «κάθονται» με τα άλλα. «Και τα τρία έχουν τέσσερις πλευρές», λέει η Αγγελική, «αλλά δεν είναι το ίδιο σχήμα». Ο Στέλιος προσπαθεί να τα βάλει όλα στην ίδια κατηγορία, η Ρήνη ρωτά ποια ιδιότητα κάνει τα πλακάκια της πρώτης στοίβας τόσο βολικά, και η Ζωή θέλει να ξέρει γιατί.

Σε αυτή την παράγραφο ονομάζουμε τα τετράπλευρα, μελετάμε το παραλληλόγραμμο και το τραπέζιο και μαθαίνουμε τις ιδιότητές τους.
\end{anagke}

\begin{lexi}{τετράπλευρο, παραλληλόγραμμο, τραπέζιο, διαγώνιος}{ονόματα τετραπλεύρων}
Το \lexiword{τετράπλευρο} έχει τέσσερις πλευρές. Το \lexiword{παραλληλόγραμμο} είναι το σχήμα που ορίζεται από «παράλληλες γραμμές». Το \lexiword{τραπέζιο} έρχεται από το «τράπεζα», το τραπέζι, γιατί μοιάζει με τραπεζάκι. Η \lexiword{διαγώνιος} είναι το τμήμα που περνά «δια γωνίας», από γωνία σε γωνία. Ισοσκελές σημαίνει «με ίσα σκέλη».
\end{lexi}

\begin{metafrasi}
Σχήμα με τέσσερις ευθύγραμμες πλευρές. & τετράπλευρο\\
Τα απέναντι ζεύγη πλευρών είναι παράλληλα. & παραλληλόγραμμο ($\mathrm{ΑΒ}\parallel\mathrm{ΓΔ}$, $\mathrm{ΑΔ}\parallel\mathrm{ΒΓ}$)\\
Μόνο ένα ζεύγος παράλληλων πλευρών. & τραπέζιο\\
Τραπέζιο με ίσες τις μη παράλληλες πλευρές. & ισοσκελές τραπέζιο\\
Η απόσταση ανάμεσα στις παράλληλες πλευρές. & ύψος\\
\end{metafrasi}

\tmima{Η θεωρία}

\begin{orismos}{τετράπλευρο}
\textbf{Τετράπλευρο} λέγεται το ευθύγραμμο σχήμα με τέσσερις πλευρές. Έχει τέσσερις κορυφές και τέσσερις γωνίες. Τα τμήματα που ενώνουν απέναντι κορυφές λέγονται \textbf{διαγώνιοι}.
\end{orismos}

\begin{center}
\begin{tikzpicture}[scale=0.8]
\draw[thick] (0.3,2.2)--(3.6,2.6)--(4,0.2)--(0,0)--cycle;
\draw[gblue,dashed] (0.3,2.2)--(4,0.2); \draw[gblue,dashed] (3.6,2.6)--(0,0);
\node[above left] at (0.3,2.2) {$\mathrm{Α}$}; \node[above right] at (3.6,2.6) {$\mathrm{Β}$};
\node[below right] at (4,0.2) {$\mathrm{Γ}$}; \node[below left] at (0,0) {$\mathrm{Δ}$};
\end{tikzpicture}
\end{center}

\begin{kanonas}{άθροισμα γωνιών τετραπλεύρου}
Μια διαγώνιος χωρίζει το τετράπλευρο σε δύο τρίγωνα, καθένα με άθροισμα γωνιών $180^\circ$. Άρα το άθροισμα των γωνιών κάθε τετραπλεύρου είναι $2\cdot180^\circ=360^\circ$.
\end{kanonas}

\begin{orismos}{παραλληλόγραμμο}
\textbf{Παραλληλόγραμμο} λέγεται το τετράπλευρο που έχει τις απέναντι πλευρές του παράλληλες: $\mathrm{ΑΒ}\parallel\mathrm{ΔΓ}$ και $\mathrm{ΑΔ}\parallel\mathrm{ΒΓ}$.
\end{orismos}

Ειδικές περιπτώσεις παραλληλογράμμου είναι το ορθογώνιο, ο ρόμβος και το τετράγωνο, που θα δούμε στην επόμενη παράγραφο.

\begin{center}
\begin{tikzpicture}[scale=0.9]
\draw[thick] (0,0)--(4,0)--(5.2,2)--(1.2,2)--cycle;
\draw[gblue,dashed] (1.2,2)--(4,0); \draw[gblue,dashed] (5.2,2)--(0,0);
\draw[orange!80!black,thick] (1.2,2)--(1.2,0); \draw (1.2,0.2)--(1.4,0.2)--(1.4,0);
\fill (2.6,1) circle (1.6pt) node[above right] {\small $\mathrm{Ο}$};
\node[above left] at (1.2,2) {$\mathrm{Α}$}; \node[above right] at (5.2,2) {$\mathrm{Β}$};
\node[below right] at (4,0) {$\mathrm{Γ}$}; \node[below left] at (0,0) {$\mathrm{Δ}$};
\node[left,orange!80!black] at (1.2,1) {\small $\upsilon$};
\end{tikzpicture}
\end{center}

\begin{idiotita}{παραλληλογράμμου}
Σε κάθε παραλληλόγραμμο $\mathrm{ΑΒΓΔ}$ ισχύουν:
\begin{itemize}[leftmargin=1.2em,itemsep=0.3ex,topsep=2pt]
\item οι \textbf{απέναντι πλευρές} είναι ίσες: $\mathrm{ΑΒ}=\mathrm{ΓΔ}$ και $\mathrm{ΑΔ}=\mathrm{ΒΓ}$,
\item οι \textbf{απέναντι γωνίες} είναι ίσες: $\hat{\mathrm{Α}}=\hat{\mathrm{Γ}}$ και $\hat{\mathrm{Β}}=\hat{\mathrm{Δ}}$,
\item οι \textbf{διαγώνιοι διχοτομούνται}: $\mathrm{ΟΑ}=\mathrm{ΟΓ}$ και $\mathrm{ΟΒ}=\mathrm{ΟΔ}$,
\item οι \textbf{διαδοχικές γωνίες} είναι παραπληρωματικές: $\hat{\mathrm{Α}}+\hat{\mathrm{Δ}}=180^\circ$.
\end{itemize}
\end{idiotita}

Η τελευταία ιδιότητα προκύπτει από όσα ξέρουμε για τις παράλληλες: οι $\hat{\mathrm{Α}}$ και $\hat{\mathrm{Δ}}$ είναι εντός και επί τα αυτά γωνίες των παραλλήλων $\mathrm{ΑΒ}$, $\mathrm{ΔΓ}$ με τέμνουσα την $\mathrm{ΑΔ}$.

Το σημείο τομής $\mathrm{Ο}$ των διαγωνίων λέγεται \textbf{κέντρο} του παραλληλογράμμου. Κάθε πλευρά μπορεί να ονομαστεί \textbf{βάση}, και η απόσταση της βάσης από την απέναντι πλευρά λέγεται \textbf{ύψος} ($\upsilon$ στο σχήμα). Η \textbf{περίμετρος} είναι το άθροισμα των πλευρών, δηλαδή $2\cdot(\mathrm{ΑΒ}+\mathrm{ΒΓ})$.

\begin{orismos}{τραπέζιο}
\textbf{Τραπέζιο} λέγεται το τετράπλευρο που έχει \emph{μόνο δύο} πλευρές παράλληλες. Οι παράλληλες λέγονται \textbf{βάσεις} (η μεγάλη και η μικρή) και η απόστασή τους \textbf{ύψος}. Αν οι άλλες δύο πλευρές είναι ίσες, το τραπέζιο λέγεται \textbf{ισοσκελές}.
\end{orismos}

\begin{center}
\begin{tikzpicture}[scale=0.75]
\draw[thick] (0,0)--(5,0)--(3.6,2)--(1,2)--cycle;
\node[below left] at (0,0) {$\mathrm{Δ}$}; \node[below right] at (5,0) {$\mathrm{Γ}$};
\node[above right] at (3.6,2) {$\mathrm{Β}$}; \node[above left] at (1,2) {$\mathrm{Α}$};
\node at (2.5,-0.7) {\small τραπέζιο ($\mathrm{ΑΒ}\parallel\mathrm{ΔΓ}$)};
\begin{scope}[xshift=7.5cm]
\draw[thick] (0,0)--(5,0)--(4,2)--(1,2)--cycle;
\draw[dashed,gblue] (2.5,2.5)--(2.5,-0.5);
\fill (2.5,2) circle (1.4pt) node[above right] {\small $\mathrm{Μ}_1$}; \fill (2.5,0) circle (1.4pt) node[below right] {\small $\mathrm{Μ}_2$};
\node[below left] at (0,0) {$\mathrm{Δ}$}; \node[below right] at (5,0) {$\mathrm{Γ}$};
\node[above right] at (4,2) {$\mathrm{Β}$}; \node[above left] at (1,2) {$\mathrm{Α}$};
\node at (2.5,-1.05) {\small ισοσκελές τραπέζιο};
\end{scope}
\end{tikzpicture}
\end{center}

\begin{idiotita}{ισοσκελούς τραπεζίου}
Σε ισοσκελές τραπέζιο $\mathrm{ΑΒΓΔ}$ ($\mathrm{ΑΒ}\parallel\mathrm{ΔΓ}$, $\mathrm{ΑΔ}=\mathrm{ΒΓ}$):
\begin{itemize}[leftmargin=1.2em,itemsep=0.3ex,topsep=2pt]
\item οι γωνίες που βρίσκονται στην ίδια βάση είναι ίσες: $\hat{\mathrm{Δ}}=\hat{\mathrm{Γ}}$ και $\hat{\mathrm{Α}}=\hat{\mathrm{Β}}$,
\item η ευθεία των μέσων $\mathrm{Μ}_1$, $\mathrm{Μ}_2$ των βάσεων είναι \textbf{άξονας συμμετρίας} και μεσοκάθετος των δύο βάσεων.
\end{itemize}
\end{idiotita}

Σε κάθε τραπέζιο οι γωνίες που βρίσκονται στην ίδια μη παράλληλη πλευρά είναι παραπληρωματικές (εντός και επί τα αυτά γωνίες των βάσεων): $\hat{\mathrm{Α}}+\hat{\mathrm{Δ}}=180^\circ$ και $\hat{\mathrm{Β}}+\hat{\mathrm{Γ}}=180^\circ$.

\begin{prosoxi}
\begin{itemize}[leftmargin=1.2em,itemsep=0.3ex,topsep=2pt]
\item Το τραπέζιο έχει \emph{ακριβώς} ένα ζεύγος παράλληλων πλευρών. Αν είχε και το δεύτερο, θα ήταν παραλληλόγραμμο.
\item Το ισοσκελές τραπέζιο έχει ίσες τις μη παράλληλες πλευρές, αλλά δεν είναι παραλληλόγραμμο.
\item Οι διαγώνιοι ενός παραλληλογράμμου διχοτομούνται, αλλά δεν είναι γενικά ίσες ούτε κάθετες.
\item Το ύψος είναι \emph{απόσταση}, κάθετη στις βάσεις, και συνήθως δεν είναι πλευρά του σχήματος.
\end{itemize}
\end{prosoxi}

\tmima{Λυμένα παραδείγματα}

\paradeigma{Πλευρές και περίμετρος}
\tixero Παραλληλόγραμμο $\mathrm{ΑΒΓΔ}$ με $\mathrm{ΑΒ}=7$ cm και $\mathrm{ΒΓ}=4$ cm.\par
\tithelo Οι άλλες δύο πλευρές και η περίμετρος.\par
\possindeo Οι απέναντι πλευρές είναι ίσες: $\mathrm{ΓΔ}=\mathrm{ΑΒ}=7$ cm και $\mathrm{ΑΔ}=\mathrm{ΒΓ}=4$ cm. Η περίμετρος είναι $7+4+7+4=22$ cm.\par
\apantisi $\mathrm{ΓΔ}=7$ cm, $\mathrm{ΑΔ}=4$ cm, περίμετρος $22$ cm.\par
\elegxos $2\cdot(7+4)=22$ cm.

\paradeigma{Γωνίες παραλληλογράμμου}
\tixero Παραλληλόγραμμο $\mathrm{ΑΒΓΔ}$ με $\hat{\mathrm{Α}}=65^\circ$.\par
\tithelo Οι άλλες γωνίες.\par
\possindeo Η $\hat{\mathrm{Γ}}$ είναι απέναντι από την $\hat{\mathrm{Α}}$, άρα $\hat{\mathrm{Γ}}=65^\circ$. Οι $\hat{\mathrm{Β}}$ και $\hat{\mathrm{Δ}}$ είναι διαδοχικές της $\hat{\mathrm{Α}}$, άρα $\hat{\mathrm{Β}}=\hat{\mathrm{Δ}}=180^\circ-65^\circ=115^\circ$.\par
\apantisi $65^\circ,115^\circ,65^\circ,115^\circ$.\par
\elegxos $65+115+65+115=360$.

\paradeigma{Οι διαγώνιοι}
\tixero Στο παραλληλόγραμμο $\mathrm{ΑΒΓΔ}$ οι διαγώνιοι τέμνονται στο $\mathrm{Ο}$ και $\mathrm{ΟΑ}=3{,}5$ cm, $\mathrm{ΟΒ}=2{,}4$ cm.\par
\tithelo Τα μήκη των διαγωνίων.\par
\possindeo Οι διαγώνιοι διχοτομούνται, άρα $\mathrm{ΑΓ}=2\cdot3{,}5=7$ cm και $\mathrm{ΒΔ}=2\cdot2{,}4=4{,}8$ cm.\par
\apantisi $\mathrm{ΑΓ}=7$ cm και $\mathrm{ΒΔ}=4{,}8$ cm.\par
\elegxos Το $\mathrm{Ο}$ είναι μέσο και των δύο διαγωνίων.

\paradeigma{Ισοσκελές τραπέζιο}
\tixero Ισοσκελές τραπέζιο $\mathrm{ΑΒΓΔ}$ με $\mathrm{ΑΒ}\parallel\mathrm{ΔΓ}$ και $\hat{\mathrm{Δ}}=70^\circ$.\par
\tithelo Οι άλλες γωνίες.\par
\possindeo Οι γωνίες της βάσης $\mathrm{ΔΓ}$ είναι ίσες: $\hat{\mathrm{Γ}}=70^\circ$. Οι $\hat{\mathrm{Α}}$ και $\hat{\mathrm{Δ}}$ είναι παραπληρωματικές, άρα $\hat{\mathrm{Α}}=180^\circ-70^\circ=110^\circ$, και $\hat{\mathrm{Β}}=\hat{\mathrm{Α}}=110^\circ$.\par
\apantisi $\hat{\mathrm{Α}}=\hat{\mathrm{Β}}=110^\circ$ και $\hat{\mathrm{Γ}}=\hat{\mathrm{Δ}}=70^\circ$.\par
\elegxos $110+110+70+70=360$.

\paradeigma{Τι είδους τετράπλευρο;}
\tixero Στο τετράπλευρο $\mathrm{ΑΒΓΔ}$ είναι $\hat{\mathrm{Α}}=105^\circ$, $\hat{\mathrm{Β}}=100^\circ$ και $\hat{\mathrm{Δ}}=75^\circ$.\par
\tithelo Η $\hat{\mathrm{Γ}}$ και το είδος του τετραπλεύρου.\par
\possindeo Το άθροισμα είναι $360^\circ$, άρα $\hat{\mathrm{Γ}}=360^\circ-(105^\circ+100^\circ+75^\circ)=80^\circ$. Με τέμνουσα την $\mathrm{ΑΔ}$ οι $\hat{\mathrm{Α}}$, $\hat{\mathrm{Δ}}$ έχουν άθροισμα $180^\circ$, άρα $\mathrm{ΑΒ}\parallel\mathrm{ΔΓ}$. Με τέμνουσα την $\mathrm{ΑΒ}$ οι $\hat{\mathrm{Α}}$, $\hat{\mathrm{Β}}$ έχουν άθροισμα $205^\circ\neq180^\circ$, άρα οι $\mathrm{ΑΔ}$ και $\mathrm{ΒΓ}$ δεν είναι παράλληλες.\par
\apantisi $\hat{\mathrm{Γ}}=80^\circ$ και το τετράπλευρο είναι τραπέζιο (όχι παραλληλόγραμμο).\par
\elegxos Στο παραλληλόγραμμο θα ήταν $\hat{\mathrm{Α}}+\hat{\mathrm{Β}}=180^\circ$.

\tmima{Ασκήσεις}

\zoni{Το κατάλαβα}{ονόματα και ιδιότητες}

\ex Σωστό ή λάθος; α) Κάθε παραλληλόγραμμο είναι τετράπλευρο. β) Κάθε τετράπλευρο είναι παραλληλόγραμμο. γ) Το τραπέζιο έχει ακριβώς ένα ζεύγος παράλληλων πλευρών. δ) Οι διαγώνιοι του παραλληλογράμμου διχοτομούνται. ε) Οι διαγώνιοι κάθε παραλληλογράμμου είναι ίσες. στ) Οι απέναντι γωνίες παραλληλογράμμου είναι ίσες.

\ex Σχεδίασε ένα παραλληλόγραμμο $\mathrm{ΑΒΓΔ}$, τις διαγωνίους του και το κέντρο $\mathrm{Ο}$. Μέτρησε τα $\mathrm{ΟΑ}$ και $\mathrm{ΟΓ}$, και φέρε ένα ύψος του.

\ex Στο παραλληλόγραμμο $\mathrm{ΑΒΓΔ}$ είναι $\mathrm{ΑΒ}=9$ cm και $\mathrm{ΒΓ}=5$ cm. Βρες τις $\mathrm{ΓΔ}$, $\mathrm{ΑΔ}$ και την περίμετρο.

\ex Στο παραλληλόγραμμο $\mathrm{ΑΒΓΔ}$ είναι $\hat{\mathrm{Α}}=72^\circ$. Βρες τις άλλες γωνίες.

\ex Οι διαγώνιοι παραλληλογράμμου $\mathrm{ΑΒΓΔ}$ τέμνονται στο $\mathrm{Ο}$ και είναι $\mathrm{ΑΓ}=10$ cm, $\mathrm{ΒΔ}=8$ cm. Βρες τα $\mathrm{ΟΑ}$, $\mathrm{ΟΓ}$, $\mathrm{ΟΒ}$, $\mathrm{ΟΔ}$.

\ex Στο ισοσκελές τραπέζιο $\mathrm{ΑΒΓΔ}$ ($\mathrm{ΑΒ}\parallel\mathrm{ΔΓ}$) είναι $\hat{\mathrm{Δ}}=68^\circ$. Βρες τις άλλες γωνίες.

\ex Στο τραπέζιο $\mathrm{ΑΒΓΔ}$ ($\mathrm{ΑΒ}\parallel\mathrm{ΔΓ}$) είναι $\hat{\mathrm{Α}}=95^\circ$ και $\hat{\mathrm{Β}}=110^\circ$. Βρες τις $\hat{\mathrm{Δ}}$ και $\hat{\mathrm{Γ}}$.

\zoni{Το χρησιμοποιώ}{πλακάκια και ζαρντινιέρες}

\ex \textbf{Το πλακάκι.} Ένα πλακάκι της θείας Μαρίας είναι παραλληλόγραμμο με πλευρές $30$ cm και $20$ cm. Θέλει να κολλήσει ταινία γύρω από την άκρη του. Πόση ταινία χρειάζεται;

\ex \textbf{Η ζαρντινιέρα.} Η πλαϊνή όψη μιας ζαρντινιέρας είναι ισοσκελές τραπέζιο με μικρή βάση $40$ cm, μεγάλη βάση $60$ cm και μη παράλληλες πλευρές $25$ cm. Οι γωνίες στη μεγάλη βάση είναι $70^\circ$. Βρες τις άλλες γωνίες και την περίμετρο.

\ex \textbf{Ο μηχανισμός του Ηλία.} Ο Ηλίας έχει τέσσερις ράβδους που ενώνονται με καρφιά σε παραλληλόγραμμο. Στην αρχή η $\hat{\mathrm{Α}}$ είναι $90^\circ$. Τον πιέζει ώστε να γίνει $\hat{\mathrm{Α}}=60^\circ$. Βρες τις άλλες γωνίες. Αλλάζουν τα μήκη των πλευρών;

\ex \textbf{Η περίμετρος.} Ένα παραλληλόγραμμο οικόπεδο έχει περίμετρο $46$ m και η μια πλευρά του είναι κατά $3$ m μεγαλύτερη από την άλλη. Βρες τις πλευρές του.

\ex \textbf{Το ύψος.} Σχεδίασε παραλληλόγραμμο $\mathrm{ΑΒΓΔ}$ με $\mathrm{ΑΒ}=5$ cm, $\mathrm{ΑΔ}=3$ cm και $\hat{\mathrm{Α}}=60^\circ$. Φέρε το ύψος από το $\mathrm{Δ}$ προς την $\mathrm{ΑΒ}$ και μέτρησέ το.

\zoni{Το σκέφτομαι}{γιατί και πώς}

\ex Εξήγησε γιατί σε κάθε παραλληλόγραμμο $\mathrm{ΑΒΓΔ}$ οι γωνίες $\hat{\mathrm{Α}}$ και $\hat{\mathrm{Δ}}$ έχουν άθροισμα $180^\circ$.

\ex Το τετράπλευρο $\mathrm{ΑΒΓΔ}$ έχει $\hat{\mathrm{Α}}=110^\circ$ και $\hat{\mathrm{Δ}}=65^\circ$. Μπορεί να είναι παραλληλόγραμμο; Δικαιολόγησε.

\ex Στο τετράπλευρο $\mathrm{ΕΖΗΔ}$ είναι $\hat{\mathrm{Δ}}=\hat{\mathrm{Η}}=90^\circ$, $\mathrm{ΔΕ}=4$ cm και $\mathrm{ΗΖ}=7$ cm. Δείξε ότι $\mathrm{ΔΕ}\parallel\mathrm{ΗΖ}$ και ότι το τετράπλευρο είναι τραπέζιο και όχι παραλληλόγραμμο.

\ex Στο ισοσκελές τραπέζιο $\mathrm{ΑΒΓΔ}$ ($\mathrm{ΑΒ}\parallel\mathrm{ΔΓ}$) είναι $\hat{\mathrm{Δ}}=\hat{\mathrm{Γ}}$. Εξήγησε γιατί $\hat{\mathrm{Α}}=\hat{\mathrm{Β}}$.

\zoni{Για το σπίτι}{ρίξε μια ματιά}

\ex Βρες στο σπίτι ή στον δρόμο τρία αντικείμενα που έχουν σχήμα παραλληλογράμμου ή τραπεζίου (παράθυρο, πλακάκι, ράφι, κερκίδα). Σχεδίασέ τα και σημείωσε ποιες πλευρές είναι παράλληλες.

\ex Σχεδίασε παραλληλόγραμμο $\mathrm{ΑΒΓΔ}$ με $\mathrm{ΑΒ}=6$ cm, $\mathrm{ΑΔ}=4$ cm και $\hat{\mathrm{Α}}=70^\circ$. Βρες με υπολογισμό και με μέτρηση τις άλλες γωνίες του.

\paragrafos{Ορθογώνια, ρόμβοι, τετράγωνα}{Β.3.11, «Είδη παραλληλογράμμων και ιδιότητες: ορθογώνια, ρόμβοι, τετράγωνα»}

\begin{anagke}
Η Αγγελική διπλώνει χαρτί σε τρία σχήματα: ένα ορθογώνιο, έναν ρόμβο και ένα τετράγωνο. Για το καθένα ψάχνει με πόσους τρόπους μπορεί να το διπλώσει ώστε τα δύο μισά να κάθονται ακριβώς το ένα πάνω στο άλλο. «Το ορθογώνιο δύο, ο ρόμβος δύο, το τετράγωνο τέσσερις!», λέει. Ο Στέλιος αναφωνεί: «Το τετράγωνο τα έχει όλα, άρα ό,τι ισχύει για τα άλλα ισχύει και για αυτό!». Η Ρήνη όμως ρωτά: «Δηλαδή το τετράγωνο είναι και ορθογώνιο και ρόμβος; Δεν βγάζει νόημα». Η Φανή γράφει στο τετράδιό της και ρωτά τι ακριβώς κάνει κάθε σχήμα «ξεχωριστό».

Σε αυτή την παράγραφο ορίζουμε το ορθογώνιο, τον ρόμβο και το τετράγωνο, βρίσκουμε τις ιδιότητές τους και καταλαβαίνουμε πώς σχετίζονται.
\end{anagke}

\begin{lexi}{ορθογώνιο, ρόμβος, τετράγωνο}{ονόματα παραλληλογράμμων}
Το \lexiword{ορθογώνιο} έχει όλες τις γωνίες \emph{ορθές} (από το «ορθός» και «γωνία»). Ο \lexiword{ρόμβος} ήταν στην αρχαιότητα η σβούρα, που «γυρίζει». Το \lexiword{τετράγωνο} έχει «τέσσερις γωνίες» ίσες, και τέσσερις ίσες πλευρές.
\end{lexi}

\begin{metafrasi}
Παραλληλόγραμμο με όλες τις γωνίες ορθές. & ορθογώνιο\\
Παραλληλόγραμμο με όλες τις πλευρές ίσες. & ρόμβος\\
Ορθογώνιο και ρόμβος μαζί. & τετράγωνο\\
Η ευθεία που διπλώνει το σχήμα σε δύο ίσα μέρη. & άξονας συμμετρίας\\
\end{metafrasi}

\tmima{Η θεωρία}

Τα τρία αυτά σχήματα είναι παραλληλόγραμμα, άρα έχουν όλες τις ιδιότητες του παραλληλογράμμου (ίσες απέναντι πλευρές και γωνίες, διχοτομούμενες διαγώνιες). Έχουν όμως και κάτι παραπάνω.

\begin{orismos}{ορθογώνιο}
\textbf{Ορθογώνιο} (παραλληλόγραμμο) λέγεται το παραλληλόγραμμο που έχει όλες τις γωνίες του ορθές.
\end{orismos}

\begin{idiotita}{ορθογωνίου}
Εκτός από τις ιδιότητες του παραλληλογράμμου, στο ορθογώνιο:
\begin{itemize}[leftmargin=1.2em,itemsep=0.3ex,topsep=2pt]
\item οι \textbf{διαγώνιοι είναι ίσες}: $\mathrm{ΑΓ}=\mathrm{ΒΔ}$,
\item οι \textbf{μεσοκάθετοι των πλευρών} του είναι \textbf{άξονες συμμετρίας} (έχει δύο άξονες).
\end{itemize}
\end{idiotita}

\begin{orismos}{ρόμβος}
\textbf{Ρόμβος} λέγεται το παραλληλόγραμμο που έχει όλες τις πλευρές του ίσες.
\end{orismos}

\begin{idiotita}{ρόμβου}
Εκτός από τις ιδιότητες του παραλληλογράμμου, στον ρόμβο:
\begin{itemize}[leftmargin=1.2em,itemsep=0.3ex,topsep=2pt]
\item οι \textbf{διαγώνιοι τέμνονται κάθετα}: $\mathrm{ΑΓ}\perp\mathrm{ΒΔ}$,
\item οι \textbf{διαγώνιοι διχοτομούν τις γωνίες} του,
\item οι \textbf{ευθείες των διαγωνίων} είναι \textbf{άξονες συμμετρίας} (έχει δύο άξονες).
\end{itemize}
\end{idiotita}

\begin{orismos}{τετράγωνο}
\textbf{Τετράγωνο} λέγεται το παραλληλόγραμμο που έχει όλες τις γωνίες του ορθές και όλες τις πλευρές του ίσες.
\end{orismos}

\begin{idiotita}{τετραγώνου}
Το τετράγωνο είναι ταυτόχρονα \textbf{ορθογώνιο και ρόμβος}, άρα έχει όλες τις ιδιότητές τους. Οι διαγώνιοι του είναι ίσες, κάθετες και διχοτομούν τις γωνίες του (κάθε γωνία της διαγωνίου με την πλευρά είναι $45^\circ$). Έχει \textbf{τέσσερις άξονες συμμετρίας}: τις δύο διαγωνίους και τις δύο μεσοκαθέτους των πλευρών.
\end{idiotita}

\begin{center}
\begin{tikzpicture}[scale=0.62]
\draw[thick] (0,0)--(4,0)--(4,2.6)--(0,2.6)--cycle;
\draw[gblue,dashed] (-0.4,1.3)--(4.4,1.3); \draw[gblue,dashed] (2,-0.4)--(2,3);
\node at (2,-1) {\small ορθογώνιο};
\begin{scope}[xshift=7cm]
\draw[thick] (2,-0.3)--(4,1.3)--(2,2.9)--(0,1.3)--cycle;
\draw[gblue,dashed] (-0.4,1.3)--(4.4,1.3); \draw[gblue,dashed] (2,-0.7)--(2,3.3);
\node at (2,-1) {\small ρόμβος};
\end{scope}
\begin{scope}[xshift=14cm]
\draw[thick] (0.2,0)--(3.8,0)--(3.8,3.6)--(0.2,3.6)--cycle;
\draw[gblue,dashed] (0,0)--(4,3.6); \draw[gblue,dashed] (0,3.6)--(4,0);
\draw[gblue,dashed] (-0.2,1.8)--(4.2,1.8); \draw[gblue,dashed] (2,-0.2)--(2,3.8);
\node at (2,-1) {\small τετράγωνο};
\end{scope}
\end{tikzpicture}
\end{center}

Το διάγραμμα δείχνει πώς συνδέονται όλα τα τετράπλευρα που μάθαμε.

\begin{center}
\begin{tikzpicture}[scale=0.9,node distance=0.8cm,
  box/.style={draw,rounded corners=2pt,fill=gblueL,inner sep=4pt,font=\small,align=center}]
\node[box] (t) at (3,3.6) {τετράπλευρα};
\node[box] (p) at (0.8,2.1) {παραλληλόγραμμα};
\node[box] (tr) at (5.2,2.1) {τραπέζια};
\node[box] (o) at (-0.6,0.6) {ορθογώνια};
\node[box] (r) at (3,0.6) {ρόμβοι};
\node[box] (q) at (1.2,-0.8) {τετράγωνα};
\draw[->,thick] (t)--(p); \draw[->,thick] (t)--(tr);
\draw[->,thick] (p)--(o); \draw[->,thick] (p)--(r);
\draw[->,thick] (o)--(q); \draw[->,thick] (r)--(q);
\end{tikzpicture}
\end{center}

Οι διαγώνιοι χωρίζουν τα σχήματα σε τρίγωνα: στο ορθογώνιο σε τέσσερα ισοσκελή τρίγωνα (γιατί $\mathrm{ΟΑ}=\mathrm{ΟΒ}=\mathrm{ΟΓ}=\mathrm{ΟΔ}$), στον ρόμβο σε τέσσερα ορθογώνια τρίγωνα (γιατί οι διαγώνιοι είναι κάθετες) και στο τετράγωνο σε τέσσερα ορθογώνια και ισοσκελή τρίγωνα. Η περίμετρος του ορθογωνίου είναι $2\cdot(\mathrm{ΑΒ}+\mathrm{ΒΓ})$ και του ρόμβου και του τετραγώνου $4\cdot(\text{πλευρά})$.

Σχετικά με το ύψος: όλες οι πλευρές του ορθογωνίου και του τετραγώνου είναι και ύψη, γιατί είναι κάθετες μεταξύ τους. Στον ρόμβο τα ύψη που φέρνουμε από μια κορυφή είναι ίσα, ενώ σε ένα τυχαίο παραλληλόγραμμο δεν είναι αναγκαστικά ίσα.

\begin{prosoxi}
\begin{itemize}[leftmargin=1.2em,itemsep=0.3ex,topsep=2pt]
\item Κάθε τετράγωνο είναι ρόμβος και ορθογώνιο, αλλά \emph{όχι} κάθε ρόμβος ή ορθογώνιο είναι τετράγωνο.
\item Οι διαγώνιοι του ορθογωνίου είναι ίσες αλλά \emph{όχι κάθετες}, ενώ του ρόμβου είναι κάθετες αλλά \emph{όχι ίσες}.
\item Το ορθογώνιο έχει δύο άξονες συμμετρίας και όχι τέσσερις· οι διαγώνιοί του \emph{δεν} είναι άξονες.
\item Ο ρόμβος δεν έχει ορθές γωνίες, εκτός αν είναι τετράγωνο.
\end{itemize}
\end{prosoxi}

\tmima{Λυμένα παραδείγματα}

\paradeigma{Το ορθογώνιο}
\tixero Ορθογώνιο $\mathrm{ΑΒΓΔ}$ με $\mathrm{ΑΒ}=8$ cm, $\mathrm{ΒΓ}=6$ cm και $\mathrm{ΑΓ}=10$ cm.\par
\tithelo Η περίμετρος, η $\mathrm{ΒΔ}$ και τα τμήματα $\mathrm{ΟΑ}$, $\mathrm{ΟΒ}$ ($\mathrm{Ο}$ το κέντρο).\par
\possindeo Η περίμετρος είναι $2\cdot(8+6)=28$ cm. Οι διαγώνιοι του ορθογωνίου είναι ίσες, άρα $\mathrm{ΒΔ}=\mathrm{ΑΓ}=10$ cm. Διχοτομούνται, άρα $\mathrm{ΟΑ}=\mathrm{ΟΒ}=10:2=5$ cm.\par
\apantisi Περίμετρος $28$ cm, $\mathrm{ΒΔ}=10$ cm, $\mathrm{ΟΑ}=\mathrm{ΟΒ}=5$ cm.\par
\elegxos Το τρίγωνο $\mathrm{ΟΑΒ}$ είναι ισοσκελές ($\mathrm{ΟΑ}=\mathrm{ΟΒ}$).

\paradeigma{Ο ρόμβος και οι γωνίες του}
\tixero Ρόμβος $\mathrm{ΑΒΓΔ}$ με $\hat{\mathrm{Α}}=40^\circ$.\par
\tithelo Οι γωνίες του και η γωνία $\widehat{\mathrm{ΑΒΟ}}$, όπου $\mathrm{Ο}$ το σημείο τομής των διαγωνίων.\par
\possindeo Η $\hat{\mathrm{Γ}}=40^\circ$ (απέναντι) και $\hat{\mathrm{Β}}=\hat{\mathrm{Δ}}=180^\circ-40^\circ=140^\circ$. Η διαγώνιος $\mathrm{ΑΓ}$ διχοτομεί την $\hat{\mathrm{Α}}$, άρα $\widehat{\mathrm{ΟΑΒ}}=20^\circ$. Οι διαγώνιοι είναι κάθετες, άρα $\widehat{\mathrm{ΑΟΒ}}=90^\circ$. Στο τρίγωνο $\mathrm{ΑΟΒ}$: $\widehat{\mathrm{ΑΒΟ}}=180^\circ-20^\circ-90^\circ=70^\circ$.\par
\apantisi $40^\circ,140^\circ,40^\circ,140^\circ$ και $\widehat{\mathrm{ΑΒΟ}}=70^\circ$.\par
\elegxos Η $\widehat{\mathrm{ΑΒΟ}}$ είναι η μισή της $\hat{\mathrm{Β}}$: $140^\circ:2=70^\circ$.

\paradeigma{Ρόμβος με διαγωνίους}
\tixero Ρόμβος με πλευρά $6{,}5$ cm και διαγωνίους $12$ cm και $5$ cm.\par
\tithelo Η περίμετρος του ρόμβου και η περίμετρος του τριγώνου $\mathrm{ΑΟΒ}$.\par
\possindeo Η περίμετρος του ρόμβου είναι $4\cdot6{,}5=26$ cm. Οι διαγώνιοι διχοτομούνται: $\mathrm{ΟΑ}=12:2=6$ cm και $\mathrm{ΟΒ}=5:2=2{,}5$ cm. Το τρίγωνο $\mathrm{ΑΟΒ}$ έχει πλευρές $6$, $2{,}5$ και $6{,}5$.\par
\apantisi $26$ cm και $6+2{,}5+6{,}5=15$ cm.\par
\elegxos Η περίμετρος του ρόμβου είναι $4$ φορές την πλευρά.

\paradeigma{Το τετράγωνο με άγνωστο}
\tixero Στο τετράγωνο $\mathrm{ΑΒΓΔ}$ είναι $\mathrm{ΑΒ}=3x+1$ και $\mathrm{ΒΓ}=2x+6$.\par
\tithelo Το $x$, η πλευρά και η περίμετρος.\par
\possindeo Όλες οι πλευρές του τετραγώνου είναι ίσες, άρα $3x+1=2x+6$, δηλαδή $x=5$. Η πλευρά είναι $3\cdot5+1=16$ και η περίμετρος $4\cdot16=64$.\par
\apantisi $x=5$, πλευρά $16$, περίμετρος $64$.\par
\elegxos $2\cdot5+6=16$.

\paradeigma{Τρίγωνα από τις διαγωνίους}
\tixero Ένα ορθογώνιο, ένας ρόμβος και ένα τετράγωνο με τις διαγωνίους τους.\par
\tithelo Σε τι τρίγωνα χωρίζονται;\par
\possindeo Στο ορθογώνιο οι διαγώνιοι είναι ίσες και διχοτομούνται, άρα τα τέσσερα τρίγωνα έχουν δύο ίσες πλευρές: ισοσκελή. Στον ρόμβο οι διαγώνιοι είναι κάθετες: ορθογώνια τρίγωνα. Στο τετράγωνο και τα δύο: ορθογώνια και ισοσκελή.\par
\apantisi Ορθογώνιο: ισοσκελή. Ρόμβος: ορθογώνια. Τετράγωνο: ορθογώνια και ισοσκελή.\par
\elegxos Το τετράγωνο έχει τις ιδιότητες και των δύο, οπότε βρίσκουμε και τα δύο είδη.

\tmima{Ασκήσεις}

\zoni{Το κατάλαβα}{ποιο σχήμα;}

\ex Σωστό ή λάθος; α) Κάθε τετράγωνο είναι ρόμβος. β) Κάθε ρόμβος είναι τετράγωνο. γ) Οι διαγώνιοι του ρόμβου είναι κάθετες. δ) Οι διαγώνιοι του ορθογωνίου είναι κάθετες. ε) Το ορθογώνιο έχει τέσσερις άξονες συμμετρίας. στ) Το τετράγωνο έχει τέσσερις άξονες συμμετρίας.

\ex Για κάθε ιδιότητα πες σε ποια από τα τρία σχήματα (ορθογώνιο, ρόμβος, τετράγωνο) ισχύει: α) όλες οι γωνίες ορθές β) όλες οι πλευρές ίσες γ) οι διαγώνιοι ίσες δ) οι διαγώνιοι κάθετες ε) οι διαγώνιοι διχοτομούν τις γωνίες στ) τέσσερις άξονες συμμετρίας.

\ex Σχεδίασε: α) ένα ορθογώνιο $5$ cm επί $3$ cm, β) έναν ρόμβο με πλευρά $3$ cm και μία γωνία $60^\circ$, γ) ένα τετράγωνο με πλευρά $4$ cm. Φέρε τις διαγωνίους τους και πες σε τι τρίγωνα χωρίζονται.

\ex Ένας ρόμβος έχει $\hat{\mathrm{Α}}=70^\circ$. Βρες τις άλλες γωνίες του.

\ex Βρες: α) την περίμετρο ορθογωνίου με πλευρές $12$ cm και $7$ cm, β) την πλευρά τετραγώνου με περίμετρο $36$ cm, γ) την πλευρά ρόμβου με περίμετρο $30$ cm.

\ex Στο ορθογώνιο $\mathrm{ΑΒΓΔ}$ η διαγώνιος $\mathrm{ΑΓ}$ είναι $14$ cm. Βρες τη $\mathrm{ΒΔ}$ και τα $\mathrm{ΟΑ}$, $\mathrm{ΟΒ}$, $\mathrm{ΟΓ}$, $\mathrm{ΟΔ}$.

\ex Στον ρόμβο $\mathrm{ΑΒΓΔ}$ είναι $\hat{\mathrm{Α}}=50^\circ$. Βρες τις γωνίες $\widehat{\mathrm{ΟΑΒ}}$ και $\widehat{\mathrm{ΟΒΑ}}$.

\zoni{Το χρησιμοποιώ}{αυλές, παράθυρα, χαρταετοί}

\ex \textbf{Η αυλή.} Η αυλή της θείας Μαρίας είναι τετράγωνη με πλευρά $4{,}5$ m. Πόσα μέτρα περιφράκτη χρειάζεται και πόσους άξονες συμμετρίας έχει το σχήμα της;

\ex \textbf{Το παράθυρο.} Ο παππούς Γιάννης ελέγχει αν ένα παραλληλόγραμμο παράθυρο είναι ορθογώνιο μετρώντας τις διαγωνίους του: $130$ cm και $128$ cm. Είναι ορθογώνιο; Γιατί;

\ex \textbf{Ο χαρταετός.} Ο Νίκος φτιάχνει χαρταετό σε σχήμα ρόμβου με πλευρά $40$ cm. Οι ξύλινες ράβδοι πάνω στις διαγωνίους είναι $64$ cm και $48$ cm. Βρες την περίμετρο του χαρταετού και την περίμετρο ενός από τα τέσσερα τρίγωνα που σχηματίζονται.

\ex \textbf{Τα πλακάκια.} Η Φανή τοποθετεί $5$ τετράγωνα πλακάκια πλευράς $25$ cm σε μία σειρά. Τι σχήμα σχηματίζουν και ποια είναι η περίμετρός του;

\ex \textbf{Ο ισχυρισμός του Στέλιου.} Ο Στέλιος λέει: «Οι διαγώνιοι του ρόμβου είναι κάθετες, άρα και οι διαγώνιοι του ορθογωνίου». Έχει δίκιο; Φέρε ένα παράδειγμα με ορθογώνιο $5$ cm επί $3$ cm.

\zoni{Το σκέφτομαι}{γιατί και πώς}

\ex Εξήγησε με αναδίπλωση γιατί οι διαγώνιοι του ρόμβου τέμνονται κάθετα.

\ex Στο τετράγωνο $\mathrm{ΑΒΓΔ}$ εξήγησε γιατί η διαγώνιος $\mathrm{ΒΔ}$ σχηματίζει γωνία $45^\circ$ με την πλευρά $\mathrm{ΑΒ}$.

\ex Στο ορθογώνιο $\mathrm{ΑΒΓΔ}$ με κέντρο $\mathrm{Ο}$ είναι $\widehat{\mathrm{ΑΒΔ}}=30^\circ$. Βρες τις $\widehat{\mathrm{ΔΒΓ}}$, $\widehat{\mathrm{ΑΟΒ}}$, $\widehat{\mathrm{ΒΟΓ}}$ και πες τι είδους τρίγωνο είναι το $\mathrm{ΟΒΓ}$.

\ex Ένας ρόμβος έχει περίμετρο $52$ cm και μία πλευρά του είναι $3x-2$ cm. Βρες το $x$.

\zoni{Για το σπίτι}{ρίξε μια ματιά}

\ex Κόψε από χαρτί ένα τετράγωνο και ένα ορθογώνιο. Δίπλωσέ τα για να βρεις τους άξονες συμμετρίας τους. Πόσους βρήκες σε κάθε σχήμα;

\ex Βρες στο σπίτι ένα αντικείμενο σχήματος ορθογωνίου, ένα σχήματος τετραγώνου και ένα σχήματος ρόμβου (π.χ. στο σκάκι, σε μοτίβο, σε μπλούζα). Μέτρησε τις διαγωνίους τους και πες ποιες είναι ίσες.

\paragrafos{Κύκλος και στοιχεία του}{Β.3.12, «Κύκλος και στοιχεία του»}

\begin{anagke}
Ο μπαμπάς της Ρήνης τοποθετεί έναν δρομολογητή Wi-Fi στο κέντρο της αυλής του σχολείου, και το σήμα του φτάνει ακριβώς $15$ μέτρα προς κάθε κατεύθυνση. Ο Ηλίας ρωτά: «Πού μπορούμε να καθίσουμε και να έχουμε σήμα;». Ο Νίκος τρέχει $15$ μέτρα από τον δρομολογητή και σταματά. Ο Στέλιος πηγαίνει $20$ μέτρα μακριά, και η Φανή μένει $10$ μέτρα μακριά. «Όλοι όσοι είναι στα $15$ μέτρα σχηματίζουν μια γραμμή», λέει η Αγγελική και με ένα σχοινί την ζωγραφίζει στο χώμα: στρογγυλή. Η Ζωή ρωτά πώς λέγονται τα διάφορα κομμάτια του σχήματος και πόσο «μεγάλο» είναι ένα τμήμα του.

Σε αυτή την παράγραφο μαθαίνουμε τον κύκλο και τα στοιχεία του, μετράμε τόξα με μοίρες και φτιάχνουμε σχήματα με τον διαβήτη.
\end{anagke}

\begin{lexi}{κύκλος, ακτίνα, χορδή, τόξο}{τα στοιχεία του κύκλου}
Ο \lexiword{κύκλος} είναι αυτό που «κυλάει». Η \lexiword{ακτίνα} είναι η «αχτίδα», μια γραμμή που ξεκινά από το κέντρο όπως οι αχτίδες του ήλιου. Το \lexiword{τόξο} είναι το όπλο που τεντώνεται, και η \lexiword{χορδή} είναι το νήμα του τόξου, που ενώνει τα δύο άκρα του. Η \lexiword{διάμετρος} είναι αυτό που μετρά «δια μέσου» του κέντρου.
\end{lexi}

\begin{metafrasi}
Όλα τα σημεία σε ίδια απόσταση από το κέντρο. & κύκλος $(\mathrm{Ο},\rho)$\\
Η απόσταση από το κέντρο. & ακτίνα $\rho$\\
Τμήμα που ενώνει δύο σημεία του κύκλου. & χορδή\\
Χορδή που περνά από το κέντρο. & διάμετρος $\delta=2\rho$\\
Μέρος του κύκλου ανάμεσα σε δύο σημεία. & τόξο $\overset{\frown}{\mathrm{ΑΒ}}$\\
\end{metafrasi}

\tmima{Η θεωρία}

\begin{orismos}{κύκλος}
\textbf{Κύκλος} λέγεται το σύνολο όλων των σημείων του επιπέδου που απέχουν την ίδια απόσταση από ένα σταθερό σημείο $\mathrm{Ο}$. Το $\mathrm{Ο}$ λέγεται \textbf{κέντρο} και η κοινή απόσταση \textbf{ακτίνα} $\rho$. Ο κύκλος με κέντρο $\mathrm{Ο}$ και ακτίνα $\rho$ συμβολίζεται $(\mathrm{Ο},\rho)$ και σχεδιάζεται με τον \textbf{διαβήτη}.
\end{orismos}

Δύο κύκλοι είναι \textbf{ίσοι} όταν έχουν ίσες ακτίνες. Δύο κύκλοι με το ίδιο κέντρο λέγονται \textbf{ομόκεντροι}.

\begin{center}
\begin{tikzpicture}[scale=0.95]
\draw[thick] (0,0) circle (1.8);
\fill (0,0) circle (1.8pt) node[below left] {\small $\mathrm{Ο}$};
\draw[gblue,thick] (0,0)--(1.27,1.27) node[midway,above left] {\small $\rho$};
\fill (1.27,1.27) circle (1.6pt) node[above right] {$\mathrm{Α}$};
\draw (-1.8,0)--(1.8,0); \fill (-1.8,0) circle (1.6pt) node[left] {$\mathrm{Β}$}; \fill (1.8,0) circle (1.6pt) node[right] {$\mathrm{Γ}$};
\node[above] at (-0.7,0.02) {\small διάμετρος};
\draw[orange!80!black,thick] (-1.5,-0.995)--(0.9,-1.56);
\fill (-1.5,-0.995) circle (1.6pt) node[left] {$\mathrm{Δ}$}; \fill (0.9,-1.56) circle (1.6pt) node[below right] {$\mathrm{Ε}$};
\node[orange!80!black] at (-0.2,-1.0) {\small χορδή};
\draw[red!70!black,very thick] (1.8,0) arc (0:45:1.8);
\node[red!70!black] at (2.35,0.75) {\small τόξο};
\end{tikzpicture}
\end{center}

\begin{itemize}[leftmargin=1.2em,itemsep=0.3ex,topsep=2pt]
\item \textbf{Χορδή} είναι το ευθύγραμμο τμήμα που ενώνει δύο σημεία του κύκλου.
\item \textbf{Διάμετρος} είναι η χορδή που περνά από το κέντρο. Είναι η μεγαλύτερη χορδή, ίση με το διπλάσιο της ακτίνας ($\delta=2\rho$), και το κέντρο είναι το μέσο της.
\item Δύο σημεία $\mathrm{Α}$, $\mathrm{Β}$ του κύκλου τον χωρίζουν σε δύο μέρη, καθένα από τα οποία λέγεται \textbf{τόξο} με άκρα τα $\mathrm{Α}$ και $\mathrm{Β}$, $\overset{\frown}{\mathrm{ΑΒ}}$. Ένα τόξο που είναι το μισό του κύκλου λέγεται \textbf{ημικύκλιο}.
\item Ο \textbf{κυκλικός δίσκος} είναι ο κύκλος μαζί με το εσωτερικό του. Ένα σημείο $\mathrm{Μ}$ του δίσκου έχει $\mathrm{ΟΜ}\le\rho$: πάνω στον κύκλο όταν $\mathrm{ΟΜ}=\rho$, μέσα του όταν $\mathrm{ΟΜ}<\rho$, έξω όταν $\mathrm{ΟΜ}>\rho$.
\end{itemize}

\begin{orismos}{επίκεντρη γωνία και μέτρο τόξου}
Μια γωνία λέγεται \textbf{επίκεντρη} όταν η κορυφή της είναι το κέντρο του κύκλου. Οι πλευρές της τέμνουν τον κύκλο στα $\mathrm{Α}$, $\mathrm{Β}$ και το τόξο $\overset{\frown}{\mathrm{ΑΒ}}$ που περιέχεται στο εσωτερικό της λέγεται \textbf{αντίστοιχο τόξο}. Το \textbf{μέτρο ενός τόξου} (σε μοίρες) είναι ίσο με το μέτρο της αντίστοιχης επίκεντρης γωνίας.
\end{orismos}

\begin{center}
\begin{tikzpicture}[scale=0.95]
\draw[thick] (0,0) circle (1.7);
\fill (0,0) circle (1.8pt) node[below] {\small $\mathrm{Ο}$};
\draw[thick] (0,0)--(1.7,0) (0,0)--(0.85,1.47);
\fill (1.7,0) circle (1.6pt) node[right] {$\mathrm{Α}$}; \fill (0.85,1.47) circle (1.6pt) node[above right] {$\mathrm{Β}$};
\draw[red!70!black,very thick] (1.7,0) arc (0:60:1.7);
\draw[gblue] (0.45,0) arc (0:60:0.45); \node[gblue] at (0.72,0.3) {\small $\omega$};
\end{tikzpicture}
\end{center}

\begin{idiotita}{τόξα και επίκεντρες γωνίες}
\begin{itemize}[leftmargin=1.2em,itemsep=0.3ex,topsep=2pt]
\item Ο κύκλος είναι ένα τόξο $360^\circ$ και το ημικύκλιο ένα τόξο $180^\circ$.
\item Δύο κάθετες διάμετροι χωρίζουν τον κύκλο σε τέσσερα ίσα τόξα των $90^\circ$.
\item Στον ίδιο κύκλο (ή σε ίσους κύκλους) δύο ίσες επίκεντρες γωνίες έχουν ίσα αντίστοιχα τόξα και αντίστροφα: δύο ίσα τόξα έχουν ίσες επίκεντρες γωνίες.
\end{itemize}
\end{idiotita}

Ο διαβήτης είναι και εργαλείο κατασκευής. Αν θέλουμε τρίγωνο με πλευρές $a,b,c$, σχεδιάζουμε ένα τμήμα $\mathrm{ΑΒ}=c$ και δύο κύκλους, $(\mathrm{Α},b)$ και $(\mathrm{Β},a)$. Το σημείο τομής τους είναι η τρίτη κορυφή.

\begin{prosoxi}
\begin{itemize}[leftmargin=1.2em,itemsep=0.3ex,topsep=2pt]
\item Ο κύκλος είναι μόνο η γραμμή· το εσωτερικό μαζί με τη γραμμή είναι ο κυκλικός δίσκος.
\item Κάθε διάμετρος είναι χορδή, αλλά όχι κάθε χορδή είναι διάμετρος. Καμία χορδή δεν είναι μεγαλύτερη από τη διάμετρο.
\item Το \emph{μέτρο} ενός τόξου (σε μοίρες) δεν είναι το \emph{μήκος} του. Δύο τόξα με το ίδιο μέτρο σε κύκλους με διαφορετική ακτίνα έχουν διαφορετικό μήκος.
\item Δεν υπάρχει πάντα τρίγωνο με τρία δοσμένα μήκη: οι κύκλοι μπορεί να μη τέμνονται.
\end{itemize}
\end{prosoxi}

\tmima{Λυμένα παραδείγματα}

\paradeigma{Ακτίνα και διάμετρος}
\tixero Ένας κύκλος έχει ακτίνα $3{,}5$ cm. Ένας άλλος έχει διάμετρο $5$ cm.\par
\tithelo Η διάμετρος του πρώτου και η ακτίνα του δεύτερου.\par
\possindeo $\delta=2\rho=2\cdot3{,}5=7$ cm. $\rho=\delta:2=5:2=2{,}5$ cm.\par
\apantisi $7$ cm και $2{,}5$ cm.\par
\elegxos Η διάμετρος είναι διπλάσια της ακτίνας.

\paradeigma{Μέσα, πάνω ή έξω;}
\tixero Κύκλος $(\mathrm{Ο},4)$ cm και σημεία με $\mathrm{ΟΑ}=3$ cm, $\mathrm{ΟΒ}=4$ cm, $\mathrm{ΟΓ}=5$ cm.\par
\tithelo Η θέση κάθε σημείου.\par
\possindeo Συγκρίνουμε με την ακτίνα $4$. Το $\mathrm{Α}$ έχει $3<4$, το $\mathrm{Β}$ έχει $4=4$, το $\mathrm{Γ}$ έχει $5>4$.\par
\apantisi Το $\mathrm{Α}$ είναι μέσα, το $\mathrm{Β}$ πάνω στον κύκλο, το $\mathrm{Γ}$ έξω.\par
\elegxos Μόνο το $\mathrm{Β}$ ανήκει στον κύκλο· τα $\mathrm{Α}$, $\mathrm{Β}$ ανήκουν στον κυκλικό δίσκο.

\paradeigma{Επίκεντρη γωνία και τόξο}
\tixero Σε κύκλο με κέντρο $\mathrm{Ο}$ η επίκεντρη γωνία $\widehat{\mathrm{ΑΟΒ}}$ είναι $70^\circ$.\par
\tithelo Το μέτρο του τόξου $\overset{\frown}{\mathrm{ΑΒ}}$ και του άλλου τόξου με άκρα $\mathrm{Α}$, $\mathrm{Β}$.\par
\possindeo Το τόξο που αντιστοιχεί στη γωνία έχει το ίδιο μέτρο, $70^\circ$. Το υπόλοιπο τόξο συμπληρώνει τον κύκλο των $360^\circ$.\par
\apantisi $70^\circ$ και $360^\circ-70^\circ=290^\circ$.\par
\elegxos $70^\circ+290^\circ=360^\circ$.

\paradeigma{Κατασκευή τριγώνου}
\tixero Θέλουμε τρίγωνο $\mathrm{ΑΒΓ}$ με $\mathrm{ΑΒ}=6$ cm, $\mathrm{ΑΓ}=5$ cm, $\mathrm{ΒΓ}=4$ cm. Υπάρχει τρίγωνο με πλευρές $2$, $3$ και $6$ cm;\par
\tithelo Η κατασκευή και η απάντηση για το δεύτερο τρίγωνο.\par
\possindeo Σχεδιάζουμε $\mathrm{ΑΒ}=6$ cm. Γράφουμε κύκλο $(\mathrm{Α},5)$ και κύκλο $(\mathrm{Β},4)$. Ένα από τα σημεία τομής τους είναι το $\mathrm{Γ}$. Με πλευρές $2$, $3$, $6$ θα χρειαζόμασταν τους κύκλους $(\mathrm{Α},3)$ και $(\mathrm{Β},2)$ με $\mathrm{ΑΒ}=6$. Οι δύο ακτίνες μαζί φτάνουν μόνο $3+2=5$ cm, λιγότερο από τα $6$ cm, άρα οι κύκλοι δεν τέμνονται.\par
\apantisi Το πρώτο τρίγωνο κατασκευάζεται, το δεύτερο δεν υπάρχει.\par
\elegxos Για να τέμνονται οι δύο κύκλοι πρέπει η τρίτη πλευρά να είναι μικρότερη από το άθροισμα των άλλων δύο.

\begin{center}
\begin{tikzpicture}[scale=0.55]
\draw[gray] (0,0) circle (5); \draw[gray] (6,0) circle (4);
\draw[thick] (0,0)--(6,0)--(3.75,3.31)--cycle;
\fill (0,0) circle (2pt) node[below left] {$\mathrm{Α}$};
\fill (6,0) circle (2pt) node[below right] {$\mathrm{Β}$};
\fill (3.75,3.31) circle (2pt) node[above] {$\mathrm{Γ}$};
\node at (3,-0.5) {\small $6$ cm}; \node at (1.5,2.0) {\small $5$ cm}; \node at (5.5,2.0) {\small $4$ cm};
\end{tikzpicture}
\end{center}

\paradeigma{Κύκλος που περνά από τρία σημεία}
\tixero Τρία σημεία $\mathrm{Α}$, $\mathrm{Β}$, $\mathrm{Γ}$ που δεν είναι στην ίδια ευθεία.\par
\tithelo Ο κύκλος που περνά από τα τρία.\par
\possindeo Το κέντρο $\mathrm{Ο}$ απέχει το ίδιο από το $\mathrm{Α}$ και το $\mathrm{Β}$, άρα βρίσκεται στη μεσοκάθετο του $\mathrm{ΑΒ}$. Ομοίως βρίσκεται στη μεσοκάθετο του $\mathrm{ΑΓ}$. Το κέντρο είναι το σημείο τομής των δύο μεσοκαθέτων. Ακτίνα είναι το $\mathrm{ΟΑ}$.\par
\apantisi Κέντρο η τομή των μεσοκαθέτων, ακτίνα $\mathrm{ΟΑ}$.\par
\elegxos Ο κύκλος περνά και από το $\mathrm{Γ}$, γιατί $\mathrm{ΟΓ}=\mathrm{ΟΑ}$.

\begin{center}
\begin{tikzpicture}[scale=0.8]
\draw[gblue,thick] (2,1) circle (2.236);

\draw[dashed] (2,-0.8)--(2,3.2); \draw[dashed] (-0.4,1.8)--(3.2,0.6);
\draw (0,0)--(4,0)--(1,3)--cycle;
\fill (0,0) circle (1.6pt) node[below left] {$\mathrm{Α}$};
\fill (4,0) circle (1.6pt) node[below right] {$\mathrm{Β}$};
\fill (1,3) circle (1.6pt) node[above] {$\mathrm{Γ}$};
\fill (2,1) circle (1.8pt) node[above right] {$\mathrm{Ο}$};
\end{tikzpicture}
\end{center}

\tmima{Ασκήσεις}

\zoni{Το κατάλαβα}{στοιχεία του κύκλου}

\ex Συμπλήρωσε: α) Ο κύκλος είναι ένα τόξο \ldots\ μοιρών. β) Το ημικύκλιο είναι ένα τόξο \ldots\ μοιρών. γ) Δύο κάθετες διάμετροι χωρίζουν τον κύκλο σε $4$ ίσα τόξα των \ldots\ μοιρών. δ) Αν $\rho=7$ cm, τότε $\delta=\ldots$ cm. ε) Αν $\delta=18$ cm, τότε $\rho=\ldots$ cm.

\ex Σωστό ή λάθος; α) Δύο κύκλοι είναι ίσοι όταν έχουν ίσες ακτίνες. β) Κάθε χορδή είναι διάμετρος. γ) Η διάμετρος είναι η μεγαλύτερη χορδή. δ) Όλα τα σημεία του κυκλικού δίσκου απέχουν από το κέντρο το πολύ όσο η ακτίνα. ε) Το μέτρο ενός τόξου είναι ίσο με το μέτρο της αντίστοιχης επίκεντρης γωνίας. στ) Μια χορδή μπορεί να είναι μεγαλύτερη από τη διάμετρο.

\ex Σχεδίασε ένα σημείο $\mathrm{Ο}$. Με τρία χρώματα βάψε: τα σημεία που απέχουν λιγότερο από $2$ cm από το $\mathrm{Ο}$, αυτά που απέχουν ακριβώς $2$ cm και αυτά που απέχουν περισσότερο από $2$ cm.

\ex Σχεδίασε κύκλο $(\mathrm{Ο},3\text{ cm})$ και πάρε πάνω του τρία σημεία $\mathrm{Α}$, $\mathrm{Β}$, $\mathrm{Γ}$. Σχεδίασε τη διάμετρο $\mathrm{ΑΑ'}$, την επίκεντρη γωνία $\widehat{\mathrm{ΑΟΒ}}$ και τη χορδή $\mathrm{ΒΓ}$.

\ex Δύο ομόκεντροι κύκλοι έχουν ακτίνα $2$ cm και διάμετρο $6$ cm αντίστοιχα. Βρες την ακτίνα του δεύτερου και πόσο απέχουν οι δύο κύκλοι.

\ex Ένας κύκλος έχει ακτίνα $6{,}5$ cm. Πόση είναι η διάμετρός του; Ποια είναι η μεγαλύτερη δυνατή χορδή του; Μπορεί να έχει χορδή $14$ cm;

\ex Σε κύκλο με κέντρο $\mathrm{Ο}$ είναι $\widehat{\mathrm{ΑΟΒ}}=125^\circ$. Βρες το μέτρο των δύο τόξων με άκρα $\mathrm{Α}$ και $\mathrm{Β}$.

\zoni{Το χρησιμοποιώ}{σήμα, ρολόγια, πίτσες}

\ex \textbf{Το Wi-Fi.} Ο δρομολογητής της αυλής έχει εμβέλεια $15$ m. Ο Ηλίας κάθεται $12$ m μακριά, η Ρήνη ακριβώς $15$ m και ο Στέλιος $18$ m. Ποιοι έχουν σήμα; Πόση είναι η μέγιστη απόσταση δύο ανθρώπων που έχουν σήμα όταν κάθονται αντίθετα, εκατέρωθεν του δρομολογητή;

\ex \textbf{Το ρολόι.} Η Αγγελική ζωγραφίζει ένα ρολόι. Πόσες μοίρες έχει η επίκεντρη γωνία ανάμεσα σε δύο διαδοχικούς αριθμούς; Πόσες ανάμεσα στο $12$ και στο $3$; Πόσο είναι το τόξο από το $12$ μέχρι το $5$; Πόσες μοίρες διαγράφει ο λεπτοδείκτης σε $10$ λεπτά;

\ex \textbf{Η πίτσα.} Ο Στέλιος κόβει μια στρογγυλή πίτσα σε $8$ ίσα κομμάτια. Πόση είναι η γωνία της κορυφής κάθε κομματιού στο κέντρο; Πόσο είναι το τόξο $3$ κομματιών;

\ex \textbf{Το παρτέρι.} Το κυκλικό παρτέρι του παππού Γιάννη έχει ακτίνα $2{,}4$ m. Πόσο μήκος έχει η μεγαλύτερη ευθεία διαδρομή μέσα στο παρτέρι; Μπορεί να τραβήξει σχοινί $3$ m ανάμεσα σε δύο σημεία του κύκλου; Και $5$ m;

\ex \textbf{Ο στόχος.} Ο Βασίλης παίζει βελάκια σε στόχο με ομόκεντρους κύκλους ακτίνας $5$, $10$ και $15$ cm. Το βελάκι του καρφώνεται $12$ cm από το κέντρο. Ανάμεσα σε ποιους δύο κύκλους βρίσκεται;

\zoni{Το σκέφτομαι}{κατασκευές}

\ex α) Κατασκεύασε τρίγωνο με πλευρές $5$, $4$ και $3$ cm. β) Υπάρχει τρίγωνο με πλευρές $7$, $2$ και $4$ cm; Δικαιολόγησε.

\ex Κατασκεύασε ισόπλευρο τρίγωνο πλευράς $5$ cm μόνο με χάρακα και διαβήτη. Εξήγησε γιατί οι τρεις πλευρές του είναι ίσες.

\ex Η Ζωή βρήκε ένα σπασμένο κυκλικό πιάτο και θέλει να βρει το κέντρο του. Πώς μπορεί, με τρία σημεία στην άκρη του πιάτου και διαβήτη, να το βρει;

\ex Η Πολυξένη σχεδιάζει κύκλο ακτίνας $3$ cm και, με το διαβήτη ανοιχτό $3$ cm, πηγαίνει από σημείο σε σημείο πάνω στον κύκλο. Μετά από πόσα βήματα επιστρέφει στην αρχή; Εξήγησε με τα τρίγωνα που σχηματίζονται.

\zoni{Για το σπίτι}{ρίξε μια ματιά}

\ex Βρες τρία κυκλικά αντικείμενα στο σπίτι. Μέτρησε τη διάμετρό τους και βρες την ακτίνα τους. Πώς θα βρεις το κέντρο ενός χάρτινου κύκλου διπλώνοντάς τον;

\ex Σχεδίασε κύκλο με διάμετρο $7$ cm, χώρισέ τον σε $6$ ίσα τόξα και φτιάξε ένα διακοσμητικό «λουλούδι».

\paragrafos{Ευθεία και κύκλος}{Β.3.13, «Σχετικές θέσεις ευθείας και κύκλου»}

\begin{anagke}
Στο κέντρο του κήπου η θεία Μαρία έχει στρογγυλό ποτιστικό που ρίχνει νερό σε απόσταση $4$ μέτρων προς όλες τις κατευθύνσεις. Κοντά στο ποτιστικό περνά ένα ευθύγραμμο μονοπάτι και η Ζωή ρωτά: «Θα βραχούμε περπατώντας;». Ο Ηλίας μετρά την απόσταση του μονοπατιού από το ποτιστικό: $2{,}5$ μέτρα. Ο Βασίλης λέει ότι θα βραχεί σε δύο σημεία, ο Νίκος ότι σε ένα και η Φανή ότι θα προσέξει μήπως δεν βραχεί καθόλου. «Εξαρτάται», λέει η κυρία Αλεξάνδρα, «από το πόσο μακριά περνά το μονοπάτι». Η Ρήνη βρίσκει ότι υπάρχει μια ειδική απόσταση όπου το μονοπάτι μόλις και μετά βίας ακουμπά το νερό.

Σε αυτή την παράγραφο βρίσκουμε πόσα κοινά σημεία μπορεί να έχουν μια ευθεία και ένας κύκλος και πώς σχετίζεται αυτό με την απόσταση της ευθείας από το κέντρο.
\end{anagke}

\begin{lexi}{τέμνουσα, εφαπτομένη, εξωτερική}{θέσεις ευθείας και κύκλου}
Η \lexiword{τέμνουσα} «κόβει» τον κύκλο. Η \lexiword{εφαπτομένη} είναι αυτή που «εφάπτεται», δηλαδή απλώς αγγίζει (το ρήμα «άπτομαι» σημαίνει «αγγίζω»). Η \lexiword{εξωτερική} ευθεία μένει «έξω» από τον κύκλο.
\end{lexi}

\begin{metafrasi}
Η ευθεία και ο κύκλος έχουν δύο κοινά σημεία. & τέμνουσα, $\mathrm{ΟΜ}<\rho$\\
Η ευθεία και ο κύκλος έχουν ένα κοινό σημείο. & εφαπτομένη, $\mathrm{ΟΜ}=\rho$\\
Η ευθεία και ο κύκλος δεν έχουν κοινό σημείο. & εξωτερική, $\mathrm{ΟΜ}>\rho$\\
\end{metafrasi}

\tmima{Η θεωρία}

Ένας κύκλος $(\mathrm{Ο},\rho)$ και μια ευθεία $\varepsilon$ μπορεί να έχουν \textbf{δύο}, \textbf{ένα} ή \textbf{κανένα} κοινό σημείο. Ονομάζουμε $\mathrm{ΟΜ}$ την \textbf{απόσταση} του κέντρου $\mathrm{Ο}$ από την ευθεία $\varepsilon$, δηλαδή το κάθετο τμήμα από το $\mathrm{Ο}$ προς την $\varepsilon$.

\begin{center}
\begin{tikzpicture}[scale=0.85]
\draw[thick] (0,0) circle (1.3);
\fill (0,0) circle (1.6pt) node[above left] {\small $\mathrm{Ο}$};
\draw[thick,gblue] (-2,-0.7)--(2,-0.7); \draw[dashed] (0,0)--(0,-0.7);
\fill (0,-0.7) circle (1.4pt) node[below right] {\small $\mathrm{Μ}$};
\node at (0,-2.9) {\small τέμνουσα, $\mathrm{ΟΜ}<\rho$};
\begin{scope}[xshift=4.8cm]
\draw[thick] (0,0) circle (1.3);
\fill (0,0) circle (1.6pt) node[above left] {\small $\mathrm{Ο}$};
\draw[thick,gblue] (-2,-1.3)--(2,-1.3); \draw[dashed] (0,0)--(0,-1.3);
\fill (0,-1.3) circle (1.4pt) node[below right] {\small $\mathrm{Μ}$};
\node at (0,-2.9) {\small εφαπτομένη, $\mathrm{ΟΜ}=\rho$};
\end{scope}
\begin{scope}[xshift=9.6cm]
\draw[thick] (0,0) circle (1.3);
\fill (0,0) circle (1.6pt) node[above left] {\small $\mathrm{Ο}$};
\draw[thick,gblue] (-2,-1.9)--(2,-1.9); \draw[dashed] (0,0)--(0,-1.9);
\fill (0,-1.9) circle (1.4pt) node[below right] {\small $\mathrm{Μ}$};
\node at (0,-2.9) {\small εξωτερική, $\mathrm{ΟΜ}>\rho$};
\end{scope}
\end{tikzpicture}
\end{center}

\begin{kanonas}{σχετικές θέσεις ευθείας και κύκλου}
Αν $d=\mathrm{ΟΜ}$ είναι η απόσταση του κέντρου από την ευθεία $\varepsilon$, τότε:
\begin{itemize}[leftmargin=1.2em,itemsep=0.3ex,topsep=2pt]
\item αν $d<\rho$, η $\varepsilon$ είναι \textbf{τέμνουσα} του κύκλου και έχουν \textbf{δύο} κοινά σημεία,
\item αν $d=\rho$, η $\varepsilon$ είναι \textbf{εφαπτομένη} του κύκλου και έχουν \textbf{ένα} κοινό σημείο, το \textbf{σημείο επαφής},
\item αν $d>\rho$, η $\varepsilon$ είναι \textbf{εξωτερική} και \textbf{δεν έχουν} κοινό σημείο.
\end{itemize}
\end{kanonas}

Όταν η $\varepsilon$ εφάπτεται στο σημείο $\mathrm{Μ}$, το $\mathrm{ΟΜ}$ είναι ακτίνα και η απόστασή της από το $\mathrm{Ο}$, άρα:

\begin{idiotita}{εφαπτομένη}
Η εφαπτομένη ενός κύκλου σε ένα σημείο του είναι \textbf{κάθετη στην ακτίνα} που καταλήγει στο σημείο επαφής. Από κάθε σημείο του κύκλου περνά \textbf{μία μόνο} εφαπτομένη.
\end{idiotita}

Αν από ένα σημείο $\mathrm{Μ}$ έξω από τον κύκλο φέρουμε τις δύο εφαπτόμενες που ακουμπούν στα $\mathrm{Α}$ και $\mathrm{Β}$, τα τμήματα $\mathrm{ΜΑ}$ και $\mathrm{ΜΒ}$ λέγονται \textbf{εφαπτόμενα τμήματα}.

\begin{prosoxi}
\begin{itemize}[leftmargin=1.2em,itemsep=0.3ex,topsep=2pt]
\item Η \emph{απόσταση} από το κέντρο είναι το κάθετο τμήμα $\mathrm{ΟΜ}$, όχι οποιοδήποτε τμήμα προς την ευθεία.
\item Όσο μικρότερη είναι η απόσταση, τόσο περισσότερο «μπαίνει» η ευθεία στον κύκλο: μικρή απόσταση σημαίνει τέμνουσα, όχι εξωτερική.
\item Η ευθεία που περνά από το κέντρο είναι τέμνουσα (απόσταση $0$) και όχι εφαπτομένη.
\item Η εφαπτομένη είναι κάθετη στην \emph{ακτίνα} του σημείου επαφής, όχι σε οποιαδήποτε ακτίνα.
\end{itemize}
\end{prosoxi}

\tmima{Λυμένα παραδείγματα}

\paradeigma{Η θέση της ευθείας}
\tixero Κύκλος $(\mathrm{Ο},5)$ cm και τρεις ευθείες με αποστάσεις από το $\mathrm{Ο}$: $3$ cm, $5$ cm και $7$ cm.\par
\tithelo Η θέση κάθε ευθείας ως προς τον κύκλο.\par
\possindeo Συγκρίνουμε την απόσταση με την ακτίνα $\rho=5$.\par
\apantisi Η ευθεία με $3<5$ είναι τέμνουσα (δύο κοινά σημεία). Η ευθεία με $5=5$ είναι εφαπτομένη (ένα κοινό σημείο). Η ευθεία με $7>5$ είναι εξωτερική (κανένα).\par
\elegxos Όσο η ευθεία απομακρύνεται, τα κοινά σημεία μειώνονται: $2$, $1$, $0$.

\paradeigma{Εφαπτομένη σε σημείο του κύκλου}
\tixero Κύκλος $(\mathrm{Ο},\rho)$ και σημείο $\mathrm{Α}$ πάνω του.\par
\tithelo Σχεδίαση της εφαπτομένης στο $\mathrm{Α}$.\par
\possindeo Φέρνουμε την ακτίνα $\mathrm{ΟΑ}$. Στο $\mathrm{Α}$ σχεδιάζουμε την ευθεία $\varepsilon$ που είναι κάθετη στην $\mathrm{ΟΑ}$ (με γνώμονα).\par
\apantisi Η $\varepsilon$ είναι η εφαπτομένη στο $\mathrm{Α}$.\par
\elegxos Η απόσταση του $\mathrm{Ο}$ από την $\varepsilon$ είναι $\mathrm{ΟΑ}=\rho$, άρα η $\varepsilon$ εφάπτεται.

\paradeigma{Κύκλος που εφάπτεται σε ευθεία}
\tixero Ευθεία $\varepsilon$ και σημείο $\mathrm{Α}$ πάνω της.\par
\tithelo Κύκλος που εφάπτεται στην $\varepsilon$ στο $\mathrm{Α}$.\par
\possindeo Φέρνουμε την κάθετη στην $\varepsilon$ στο $\mathrm{Α}$. Παίρνουμε πάνω της ένα οποιοδήποτε σημείο $\mathrm{Κ}$ και γράφουμε τον κύκλο $(\mathrm{Κ},\mathrm{ΚΑ})$.\par
\apantisi Ο κύκλος $(\mathrm{Κ},\mathrm{ΚΑ})$.\par
\elegxos Η απόσταση του $\mathrm{Κ}$ από την $\varepsilon$ είναι $\mathrm{ΚΑ}$, ίση με την ακτίνα. Υπάρχουν άπειροι τέτοιοι κύκλοι, ανάλογα με τη θέση του $\mathrm{Κ}$.

\paradeigma{Εφαπτόμενες στα άκρα χορδής}
\tixero Κύκλος $(\mathrm{Ο},\rho)$ και χορδή $\mathrm{ΑΒ}$.\par
\tithelo Οι εφαπτόμενες στα $\mathrm{Α}$ και $\mathrm{Β}$.\par
\possindeo Σχεδιάζουμε τις ακτίνες $\mathrm{ΟΑ}$ και $\mathrm{ΟΒ}$. Στο $\mathrm{Α}$ φέρνουμε κάθετη στην $\mathrm{ΟΑ}$ και στο $\mathrm{Β}$ κάθετη στην $\mathrm{ΟΒ}$. Αυτές είναι οι εφαπτόμενες και τέμνονται σε ένα σημείο $\mathrm{Μ}$.\par
\apantisi Οι εφαπτόμενες τέμνονται στο $\mathrm{Μ}$ και τα $\mathrm{ΜΑ}$, $\mathrm{ΜΒ}$ είναι τα εφαπτόμενα τμήματα.\par
\elegxos Το $\mathrm{Μ}$ είναι έξω από τον κύκλο.

\begin{center}
\begin{tikzpicture}[scale=0.9]
\draw[thick] (0,0) circle (1.6);
\fill (0,0) circle (1.6pt) node[above] {\small $\mathrm{Ο}$};
\fill (-1.028,-1.226) circle (1.6pt) node[left] {$\mathrm{Α}$};
\fill (1.028,-1.226) circle (1.6pt) node[right] {$\mathrm{Β}$};
\fill (0,-2.09) circle (1.6pt) node[below] {$\mathrm{Μ}$};
\draw (-1.028,-1.226)--(1.028,-1.226);
\draw[gblue,thick] (-1.75,-0.02)--(0,-2.09)--(1.75,-0.02);
\draw[dashed] (0,0)--(-1.028,-1.226) (0,0)--(1.028,-1.226);
\end{tikzpicture}
\end{center}

\tmima{Ασκήσεις}

\zoni{Το κατάλαβα}{τέμνουσα, εφαπτομένη, εξωτερική}

\ex Για κύκλο με ακτίνα $\rho=6$ cm και ευθείες με απόσταση από το κέντρο $d=4$ cm, $d=6$ cm και $d=9$ cm, πες αν η ευθεία είναι τέμνουσα, εφαπτομένη ή εξωτερική και πόσα κοινά σημεία έχει με τον κύκλο.

\ex Βρες πόσα κοινά σημεία έχουν ο κύκλος και η ευθεία όταν: α) $\rho=8$, $d=3$ β) $\rho=4$, $d=4$ γ) $\rho=3$, $d=6$ δ) $\rho=2{,}5$, $d=2{,}5$ ε) $\rho=10$, $d=12$.

\ex Σωστό ή λάθος; α) Η εφαπτομένη έχει ένα μόνο κοινό σημείο με τον κύκλο. β) Η εφαπτομένη είναι κάθετη στην ακτίνα του σημείου επαφής. γ) Αν μια ευθεία είναι εξωτερική, η απόστασή της από το κέντρο είναι μικρότερη από την ακτίνα. δ) Από ένα σημείο του κύκλου περνά μία μόνο εφαπτομένη. ε) Μια ευθεία μπορεί να έχει τρία κοινά σημεία με έναν κύκλο. στ) Κάθε ευθεία που περνά από το κέντρο είναι τέμνουσα.

\ex Σχεδίασε κύκλο $(\mathrm{Ο},4\text{ cm})$ και τρεις ευθείες που απέχουν από το $\mathrm{Ο}$ $3$ cm, $4$ cm και $5$ cm. Πες τη θέση της καθεμιάς.

\ex Σχεδίασε έναν κύκλο, πάρε ένα σημείο $\mathrm{Α}$ πάνω του και σχεδίασε την εφαπτομένη στο $\mathrm{Α}$.

\ex Κύκλος $(\mathrm{Ο},3\text{ cm})$ και ευθεία $\varepsilon$ με απόσταση $3$ cm από το $\mathrm{Ο}$. Πόσα κοινά σημεία έχουν; Πώς πρέπει να αλλάξει η απόσταση για να έχουν δύο κοινά σημεία;

\zoni{Το χρησιμοποιώ}{κήπος, τροχοί, μπάλες}

\ex \textbf{Το ποτιστικό.} Το ποτιστικό της θείας Μαρίας βρέχει σε ακτίνα $4$ m. Τρία ευθύγραμμα μονοπάτια απέχουν από αυτό $2{,}5$ m, $4$ m και $6$ m. Για κάθε μονοπάτι πες πόσα σημεία του βρέχονται.

\ex \textbf{Ο τροχός.} Ο Ηλίας κινεί το ποδήλατό του σε ευθύγραμμο δρόμο. Ο τροχός έχει ακτίνα $35$ cm. Πόσα σημεία του τροχού ακουμπούν στον δρόμο και πόσο ψηλά από το έδαφος βρίσκεται ο άξονας (το κέντρο) του τροχού;

\ex \textbf{Η μπάλα.} Ο Νίκος ακουμπά μια μπάλα του μπάσκετ ακτίνας $12$ cm στη γωνία δαπέδου και τοίχου. Πόσο απέχει το κέντρο της από τον τοίχο και πόσο από το δάπεδο; Τι ισχύει για τον τοίχο και το δάπεδο ως προς τον κύκλο που σχηματίζει η μπάλα στο πλάι;

\ex \textbf{Η λίμνη.} Στο πάρκο υπάρχει κυκλική λίμνη ακτίνας $50$ m. Ένας δρόμος περνά σε απόσταση $38$ m από το κέντρο της. Είναι εφικτός ο δρόμος; Πόση πρέπει να είναι η απόσταση ώστε ο δρόμος να περνά ακριβώς στην άκρη της λίμνης, χωρίς να τη διασχίζει;

\ex \textbf{Το σήμα.} Ο δρομολογητής Wi-Fi της Ρήνης έχει εμβέλεια $25$ m. Ο δρόμος απέναντι απέχει $30$ m από αυτόν. Φτάνει το σήμα στον δρόμο; Αν ο δρόμος ήταν σε απόσταση $25$ m, σε πόσα σημεία θα έφτανε; Και αν ήταν σε απόσταση $20$ m;

\zoni{Το σκέφτομαι}{κατασκευές και συμπεράσματα}

\ex Σε κύκλο $(\mathrm{Ο},\rho)$ φέρε δύο κάθετες διαμέτρους $\mathrm{ΑΒ}$ και $\mathrm{ΓΔ}$ και σχεδίασε τις εφαπτόμενες στα $\mathrm{Α}$, $\mathrm{Β}$, $\mathrm{Γ}$, $\mathrm{Δ}$. Τι σχήμα σχηματίζουν οι τέσσερις εφαπτόμενες και πόση είναι η πλευρά του;

\ex Σχεδίασε κύκλο $(\mathrm{Ο},4{,}5\text{ cm})$ και μια διάμετρο $\mathrm{ΑΒ}$. Σχεδίασε τις εφαπτόμενες του κύκλου που είναι παράλληλες προς την $\mathrm{ΑΒ}$. Πόσες είναι;

\ex Οι ευθείες $\varepsilon_1$ και $\varepsilon_2$ είναι εφαπτόμενες ενός κύκλου στα σημεία $\mathrm{Α}$ και $\mathrm{Β}$. Περίγραψε πώς μπορούμε να βρούμε το κέντρο του κύκλου.

\ex Πόσες εφαπτόμενες του κύκλου περνούν από ένα σημείο: α) πάνω στον κύκλο; β) έξω από τον κύκλο; γ) μέσα στον κύκλο; Δικαιολόγησε.

\zoni{Για το σπίτι}{ρίξε μια ματιά}

\ex Πάρε ένα κυκλικό καπάκι και έναν χάρακα. Τοποθέτησε τον χάρακα ώστε να εφάπτεται στο καπάκι. Μέτρησε την απόσταση του κέντρου από την άκρη του χάρακα. Τι βρίσκεις;

\ex Σχεδίασε ένα ποδήλατο: δύο ίσους κύκλους που εφάπτονται στην ίδια ευθεία (τον δρόμο). Πού βρίσκονται τα κέντρα τους ως προς τον δρόμο;

\epanalipsi{Β3}

\tmima{Λεξικό του κεφαλαίου}

\renewcommand{\arraystretch}{1.5}
\begin{tabular*}{\linewidth}{@{\extracolsep{\fill}} >{\raggedright\arraybackslash}p{0.28\linewidth} >{\raggedright\arraybackslash}p{0.62\linewidth}}
\toprule
\textbf{Λέξη ή σύμβολο} & \textbf{Με μία πρόταση}\\
\midrule
ισόπλευρο, ισοσκελές, σκαληνό & Τρεις, δύο ή καμία ίση πλευρά.\\
οξυγώνιο, ορθογώνιο, αμβλυγώνιο & Όλες οι γωνίες οξείες, μία ορθή, μία αμβλεία.\\
διάμεσος, ύψος, διχοτόμος & Προς το μέσο της απέναντι πλευράς, κάθετο στην απέναντι πλευρά, χωρίζει τη γωνία στη μέση.\\
$\hat{\mathrm{Α}}+\hat{\mathrm{Β}}+\hat{\Gamma}=180^\circ$ & Το άθροισμα των γωνιών τριγώνου.\\
εξωτερική γωνία & Ίση με το άθροισμα των δύο απέναντι εσωτερικών.\\
παραλληλόγραμμο & Απέναντι πλευρές παράλληλες, ίσες· διαγώνιες διχοτομούνται.\\
ορθογώνιο, ρόμβος, τετράγωνο & Ορθές γωνίες· ίσες πλευρές· και τα δύο μαζί.\\
τραπέζιο, ισοσκελές τραπέζιο & Μόνο δύο παράλληλες πλευρές· και ίσες οι άλλες δύο.\\
\bottomrule
\end{tabular*}

\tmima{Ασκήσεις επανάληψης}

\ex Τρίγωνο έχει πλευρές $9$, $9$ και $12$ cm. Ποιο είναι το είδος του και η περίμετρός του;

\ex Βρες την τρίτη γωνία τριγώνου με γωνίες $58^\circ$ και $64^\circ$. Τι είδους είναι;

\ex Ισοσκελές τρίγωνο έχει γωνία κορυφής $36^\circ$. Βρες τις γωνίες βάσης.

\ex Σε τρίγωνο η εξωτερική γωνία στο $\mathrm{Β}$ είναι $130^\circ$ και η $\hat{\mathrm{Α}}=45^\circ$. Βρες τις $\hat{\mathrm{Β}}$ και $\hat{\Gamma}$.

\ex Οι οξείες ενός ορθογωνίου τριγώνου είναι $x$ και $x+20^\circ$. Βρες τες.

\ex Πόσα ύψη έχει ένα τρίγωνο; Πού βρίσκονται σε αμβλυγώνιο;

\ex Βρες την περίμετρο ρόμβου με πλευρά $7{,}5$ cm και παραλληλογράμμου με πλευρές $9$ και $6$ cm.

\ex Ορθογώνιο $9\times12$ cm. Βρες τη διαγώνιο και την απόσταση του κέντρου συμμετρίας από μια κορυφή.

\ex Ρόμβος με διαγώνιες $14$ και $48$ cm. Βρες την πλευρά και την περίμετρό του.

\ex Παραλληλόγραμμο έχει γωνία $75^\circ$. Βρες τις άλλες.

\ex Ισοσκελές τραπέζιο έχει γωνίες βάσης $58^\circ$. Βρες τις άλλες.

\ex Πόσους άξονες συμμετρίας έχουν: ισόπλευρο τρίγωνο, ορθογώνιο, ρόμβος, τετράγωνο, ισοσκελές τραπέζιο;

\ex Η περίμετρος ισοσκελούς τραπεζίου είναι $46$ cm, οι βάσεις $16$ και $10$ cm. Βρες τις ίσες πλευρές.

\ex Ποιο παραλληλόγραμμο έχει ίσες και κάθετες διαγώνιες;

\ex \textbf{Ερωτήσεις αυτοαξιολόγησης (τρίγωνα).} Σωστό ή λάθος; α) Το ισόπλευρο έχει όλες τις γωνίες $60^\circ$. \quad β) Τρίγωνο έχει δύο ορθές γωνίες. \quad γ) Οι γωνίες βάσης ισοσκελούς είναι ίσες. \quad δ) Η εξωτερική γωνία ισούται με το άθροισμα των δύο απέναντι.

\ex \textbf{Ερωτήσεις αυτοαξιολόγησης (τετράπλευρα).} Σωστό ή λάθος; α) Κάθε ρόμβος είναι παραλληλόγραμμο. \quad β) Οι διαγώνιες ρόμβου είναι ίσες. \quad γ) Οι διαγώνιες ορθογωνίου είναι ίσες. \quad δ) Το τραπέζιο έχει τέσσερις παράλληλες πλευρές.

\tmima{Ο χάρτης της επανάληψης}

Μετά τις λύσεις, δες σε ποιες ασκήσεις δυσκολεύτηκες. Ο πίνακας σου λέει ποια παράγραφο να ξαναδιαβάσεις.

\renewcommand{\arraystretch}{1.5}
\begin{tabular*}{\linewidth}{@{\extracolsep{\fill}} >{\raggedright\arraybackslash}p{0.38\linewidth} >{\raggedright\arraybackslash}p{0.52\linewidth}}
\toprule
\textbf{Αν δυσκολεύτηκες στις} & \textbf{ξαναδιάβασε τις παραγράφους}\\
\midrule
ΕΒ3.1, ΕΒ3.6 & 3.1 (στοιχεία και είδη)\\
ΕΒ3.2, ΕΒ3.3, ΕΒ3.4, ΕΒ3.5 & 3.2 (άθροισμα γωνιών)\\
ΕΒ3.7 & 3.3 (είδη τετραπλεύρων)\\
ΕΒ3.8, ΕΒ3.9, ΕΒ3.10, ΕΒ3.11, ΕΒ3.13 & 3.4 (ιδιότητες)\\
ΕΒ3.12, ΕΒ3.14 & 3.4 και Β2 (συμμετρία)\\
ΕΒ3.15, ΕΒ3.16 & 3.1 έως 3.4\\
\bottomrule
\end{tabular*}

